\chapter{绪论}

\section{研究背景及意义}
近年来，随着智能控制、人工智能与机器人技术的迅速发展，以无人机和无人车为代表的自主智能无人系统应运而生，并在“低空经济”等新兴产业形态的推动下，逐步从实验室走向大规模实际应用，极大程度地提高了人类的生产效率和生活质量。无人机（Unmanned Aerial Vehicle, UAV）和无人车（Unmanned Ground Vehicle, UGV）作为重要的自主机器人，已经在多个领域展现出巨大的应用潜力。无人机以其高机动性和灵活性\cite{uavsurvey}，特别适用于执行空中任务，如农业维保\cite{agricultureusage}、灾害救援\cite{disasterusage}等；而无人地面车辆则凭借其高负载能力和稳定的地面作业性能，广泛应用于复杂地形中的运输、环境信息采集及基础设施巡检\cite{ugvinspection}等任务，其典型应用场景如\autoref{fig:uav_ugv_usage}所示。

\begin{figure}[htbp]
    \centering
    \begin{subfigure}[b]{0.4\textwidth}
        \centering
        \includegraphics[width=\textwidth]{chapter_1/Agricultural_robots.jpg}
        \caption{农业}
        % \label{fig:sub_a}
    \end{subfigure}
    \hspace{0.01\textwidth}
    \begin{subfigure}[b]{0.4\textwidth}
        \centering
        \includegraphics[width=\textwidth]{chapter_1/savesearch.png}
        \caption{搜救}
        % \label{fig:sub_b}
    \end{subfigure}
    \vspace{0.5cm}
    \begin{subfigure}[b]{0.4\textwidth}
        \centering
        \includegraphics[width=\textwidth]{chapter_1/search_ugv.jpg}
        \caption{巡检}
        % \label{fig:sub_c}
    \end{subfigure}
    \hspace{0.01\textwidth}
    \begin{subfigure}[b]{0.4\textwidth}
        \centering
        \includegraphics[width=\textwidth]{chapter_1/ugv_trans.jpg}
        \caption{运输}
        % \label{fig:sub_d}
    \end{subfigure}

    \caption{无人系统应用场景}
    \label{fig:uav_ugv_usage}
\end{figure}

进而，越来越多的研究和应用开始将无人机与无人车组合构成空地协同机器人系统，以充分发挥两者在空间机动性与地面作业能力上的互补优势，从而在执行任务时实现更高的效率与灵活性\cite{collaborationsurvey}。空地系统的主要优势体现在两个方面：一方面，无人机具备显著的空中视角优势，能够在不受地形强约束的条件下快速覆盖大范围区域，灵活绕避障碍物，完成环境感知与信息采集等任务；另一方面，无人车则依托其较强的负载能力和稳定的地面运动性能，可在复杂地形中长时间持续作业，执行物资运输、近地作业与精细操作等功能。因此，将无人机与无人车有机结合，不仅可以弥补单一平台在感知范围、负载能力或续航等方面的性能局限，还能够在高负荷、高风险任务中提升整体系统的安全性与可靠性。然而，尽管空地协同系统展现出显著优势，其在实际复杂环境中的大规模落地仍面临诸多挑战。特别是在期望实现真正意义上的多机器人自主协同作业时，状态估计、环境感知、轨迹规划与跟踪控制等模块共同构成了移动机器人系统的核心基础能力。如何在不依赖先验环境信息或外部基础设施的条件下，通过上述功能模块实现多机器人系统内部的准确相对定位、稳定感知反馈以及高效协同决策，仍是当前空地协同系统研究中的关键难题之一。

\begin{figure}[htbp]
    \centering
    \begin{subfigure}[b]{0.3\textwidth}
        \centering
        \includegraphics[width=\textwidth]{chapter_1/lidar.jpg}
        \caption{激光雷达}
        % \label{fig:sub_a}
    \end{subfigure}
    \hspace{0.01\textwidth}
    \begin{subfigure}[b]{0.3\textwidth}
        \centering
        \includegraphics[width=\textwidth]{chapter_1/t265.jpg}
        \caption{视觉传感器}
        % \label{fig:sub_b}
    \end{subfigure}
    \hspace{0.01\textwidth}
    \begin{subfigure}[b]{0.3\textwidth}
        \centering
        \includegraphics[width=\textwidth]{chapter_1/uwb芯片.jpg}
        \caption{UWB定位芯片}
        % \label{fig:sub_c}
    \end{subfigure}
    \caption{传感器硬件示意图}
    \label{fig:loc_and_per}
\end{figure}

在机器人定位与感知领域，状态估计的精度直接决定了其导航与任务执行能力，而环境感知的分辨率与全面性则约束了规划阶段可用地图信息的细节程度。目前，定位传感器的典型应用场景大致可分为室外与室内两类，感知硬件示意如\autoref{fig:loc_and_per}所示。对于开阔的室外环境，全球定位系统（Global Positioning System, GPS）\cite{gps} 凭借其广覆盖与易用性，长期以来一直是移动机器人定位的首选方案。然而，在室内空间、城市峡谷等遮挡严重的 GPS 拒止环境中，GPS 信号易出现严重衰减甚至完全丢失，导致定位精度显著下降。针对室内高精度定位需求，运动捕捉系统\cite{motioncapture} 虽然可以提供毫米级的高精度位姿估计，但其价格昂贵、部署复杂且强依赖固定场地，难以适应大范围、动态变化的未知环境。因此，移动机器人在实际应用中更常搭载激光雷达（Light Detection And Ranging, LiDAR）与视觉传感器，并与内置的惯性测量单元（Inertial Measurement Unit, IMU）进行多传感器融合，形成视觉惯性里程计（Visual Inertial Odometry, VIO）和激光雷达惯性里程计（LiDAR Inertial Odometry, LIO）等技术路线。在此基础上发展出的即时定位与建图（Simultaneous Localization and Mapping, SLAM）技术，使得机器人能够仅依赖机载传感器在未知环境中实现自主建图与自定位。然而，现有 SLAM 研究大多聚焦于单机器人场景，面向（异构）机器人集群的协同 SLAM 仍面临诸多挑战，例如多机器人间的通信与数据共享机制、初始位姿标定策略、动态环境下的鲁棒性保证以及集中式与分布式计算架构的取舍等。

近年来，超宽带（Ultra-Wideband, UWB）技术\cite{uwbsurvey} 凭借其高定位精度、低功耗以及较强的抗多径干扰能力，逐渐成为 GPS 拒止环境下无人系统定位的重要研究方向。通过在环境中部署 UWB 基站并在移动平台上安装 UWB 标签，可实现基于飞行时间测距的高精度位置估计。与运动捕捉系统相比，UWB 系统在成本和部署复杂度方面具有显著优势，可视为一种“平价”的室内高精度定位替代方案。但传统 UWB 系统同样依赖预先标定好的固定基站布局，其有效工作区域通常局限于事先部署完成的封闭场景，难以直接推广到大尺度、结构复杂或未知的任务环境。将 UWB 与 VIO/LIO 等机载 SLAM 技术进行多传感器融合，是提升定位精度与鲁棒性的一条有效途径，尤其在协同定位任务中，UWB 测距信息可以为多机器人间提供可靠的相对位置约束。然而，由于 UWB 测量对视距链路（Line-of-Sight, LoS）较为敏感，在存在严重遮挡或环境结构复杂的情况下，其测量性能容易退化。当前多数工作或假设空旷无遮挡环境，或依赖固定 UWB 基站作为静态锚点，这在很大程度上限制了其在基础设施未知或遮挡严重场景中的可扩展性。

在机器人运动规划领域，如何在复杂环境中生成既安全又高效的行动轨迹，是实现多机器人系统自主导航的关键。为此，大量研究聚焦于异构机器人系统，尤其是空地协作系统，利用空中与地面平台在感知范围、机动性能和负载能力上的互补优势，共同完成环境探索\cite{agexplore1,agexplore2}与地图重建\cite{agmap1,agmap2}等任务，如\autoref{fig:explore_task}所示，从而弥补单一平台在整体性能与任务适应性上的固有限制。在此类协作任务中，多机器人通常依赖系统内部的无线通信实现信息交互与行动协调，但受限于机载通信设备功率以及环境中障碍物的遮挡，链路质量与通信可靠性往往难以保证。特别是在部署于未知环境时，机器人之间的相互可观测性不仅关系到相对定位与实时态势感知的准确性，也是维持协同控制与全局任务推进的基础。现有研究中，多数工作基于理想化的范围通信模型\cite{communication1,communication2,communication3,communication4}，通常假设在一定距离内通信链路始终可用，而忽略了障碍物对视距链路（Line-of-Sight, LoS）的阻断及其对无线信号质量的显著影响，这在实际环境中容易导致系统内部的通信中断与观测丢失\cite{liu2023relative}。因此，一个核心挑战是在存在障碍物的复杂场景下，既要完成外部给定的导航与任务目标，又要在任务执行过程中持续维护多机器人之间的 LoS 连通性\cite{multirobotvis}。进一步地，当机器人部署于完全未知环境时，环境结构与障碍物分布具有不可预知性，这不仅加剧了对 LoS 连通的破坏风险，同时也使得传统连通性控制方法中常用的“一致、简洁障碍物集合”建模假设难以成立；在分布式感知与异构传感配置下，为整个多机器人系统构建统一的环境障碍物表示尤为困难\cite{multirobotvis}。

\begin{figure}[htbp]
    \centering
    \includegraphics[width=.7\linewidth]{chapter_1/cerberus.png}
    \caption{多机器人环境探索与地图重建任务}
    \label{fig:explore_task}
\end{figure}

综上所述，异构多机器人在未知环境中的相对定位与协同规划仍存在关键短板，尤其缺乏一种能够将定位精度诉求与任务目标驱动的协同规划过程进行统一建模并加以联合优化的系统化框架。在复杂遮挡环境中，如何在保证外部任务推进的同时，维持多机器人间的有效可见性与可靠相对定位，进而提升整体协同运行效率与鲁棒性，仍是有待深入研究的重要问题。基于此，本文构建了一套空地协同平台，并提出面向可见性编队的协同定位与规划框架，形成“定位—感知—规划”的耦合机制，在复杂遮挡环境下显著提升系统的定位稳定性与协同作业能力，具有重要的理论与工程应用意义。

围绕上述目标，本文首先搭建基于动态 UWB 基站的空地协同定位平台，结合多边定位与多源传感信息融合实现空地系统间的相对定位，并在自行搭建的硬件平台中完成验证，为后续规划模块提供稳定可靠的状态与环境输入。在此基础上，本文面向 LoS 约束与可见性需求，设计了兼顾避障安全性与协同效率的空地协同规划方法，将多机器人间的可见性维护与任务轨迹生成有机融合，形成“全局引导—局部优化”的两阶段求解思路。实验结果表明，所构建的空地机器人平台能够在复杂遮挡环境下实现较高精度的相对定位，并在保证轨迹安全与可执行性的前提下有效提升编队协同效率。

\section{国内外研究现状}
尽管针对多机器人系统的定位与感知框架已有大量研究工作，但现有方法在工程应用层面普遍存在初始化流程复杂、系统部署门槛较高等问题；同时，将定位感知问题与协同规划任务进行统一建模并加以联合优化的研究仍相对有限，定位与规划模块往往以松耦合的方式分别设计。为此，本节将从多机器人定位与感知、多机器人协同规划等几个方面，对国内外相关研究现状进行系统梳理与分析。

\subsection{无人系统定位研究现状}
在异构机器人系统中，依据不同平台的机械与载荷特性选择传感器完成环境感知与定位已成共识。空地协同场景下，轻载无人机通常配备深度相机与 IMU，并在部分场景中辅以 UWB 芯片，实现机载感知与状态估计；而载荷与供电能力更强的无人车或大型无人机则常配备激光雷达，以补充环境几何细节并提升定位鲁棒性。

以视觉为核心的定位方法经历了从基于滤波到基于优化的演进。Mourikis 等人\cite{mourikis2007msckf} 提出的 MSCKF 方法在滑动窗口 EKF 框架下实现低延迟、低资源占用的 VIO；Leutenegger 等人\cite{leutenegger2015okvis} 提出的 OKVIS 则采用关键帧驱动的非线性优化，将 IMU 预积分与图像重投影残差紧耦合，显著提升精度与稳健性；Qin 等人提出的 VINS-Mono\cite{qin2018vins} 与 VINS-Fusion\cite{qin2019vinsfusion} 在优化式框架内引入鲁棒特征管理、在线外参与时间偏差自标定、多传感器融合及回环闭环，使 VIO 在长时运行与跨场景部署方面更具通用性；Campos 等人\cite{campos2021orbslam3} 提出的 ORB-SLAM3 面向单目/双目/鱼眼与视觉–惯性统一建模，并通过多地图管理与重定位机制提升跨运动模式与长时间运行的稳定性。尽管上述方法在静态、纹理充足场景中已能提供较高精度，但在低纹理、强光照变化、快速机动等条件下，视觉定位仍易退化并产生尺度漂移，且系统对初始化、时间同步与外参标定较为敏感。

为抑制纯视觉定位方法的退化，部分研究引入了全局几何约束，UWB 因低成本、轻量化和较强抗多径能力被广泛用于为视觉–惯性估计提供距离锚定。Queralta 等人\cite{queralta2020uwb} 在 GPS 拒止场景下发布 UWB 数据集，验证单/多点测距可为无人机导航提供厘米–分米级距离约束；Lin 等人\cite{lin2023gnssuwb} 在室内环境中实证表明将 UWB 引入 VIO 可显著提升定位精度与稳定性。围绕锚点布局与量测质量，Zhao 等人\cite{zhao2022finding} 给出基于 TDoA 的锚点几何优化准则，Delama 等人提出的 UVIO 框架\cite{delama2023uvio} 则在锚点标定不完美的情况下，通过偏置补偿与紧耦合建模稳定引入 UWB 距离约束。Nguyen 等人\cite{nguyen2021range} 提出的“以测距为中心”的相机–IMU–UWB 紧耦合方案以及 Li 等人\cite{li2024uwbvio_fg} 的因子图建模方法，均利用 UWB 提供绝对尺度与弱全局约束以抑制单目尺度漂移；Jiang 等人\cite{jiang2024visionuwb} 和 Li 等人\cite{li2024uvioakf} 进一步结合 LoS/NLoS 判别与自适应权重调整，以提升在遮挡和多径条件下的鲁棒性。总体来看，UWB 为 VIO/VSLAM 提供了有力的距离锚定，但其效果仍严重依赖锚点–标签间的 LoS 质量与几何可观测性。

激光雷达虽然体积与重量相对较大，但其高精度点云在感知与定位中具有重要价值。Xu 等人\cite{xu2021fastlio} 提出的 FAST-LIO 框架采用紧耦合 ESKF 将 LiDAR 特征与 IMU 融合，在快速机动与噪声场景下实现高频、稳健的里程计估计；FAST-LIO2\cite{xu2022fastlio2} 进一步避免显式特征提取，直接构造点到局部几何面的残差，并引入增量式 iKD-Tree 实现高效建图。在图优化与回环方向，Zhang 等人\cite{zhang2014loam} 提出的 LOAM 里程计与建图实现了解耦，使用角点与面点约束实现高精度实时定位与增量建图，Shan 等人\cite{shan2020liosam,kim2018scancontext} 提出的 LIO-SAM 将 IMU 预积分、LiDAR 残差与回环检测统一于因子图平滑，并结合 Scan Context 方法进行全局重定位，有效提升长时定位鲁棒性与回访位姿一致性。但值得指出的是，当 LiDAR 视场内几何约束不足（如近似共面、共线或远距离稀疏回波）时，定位系统仍存在退化现象，需借助回环与外源约束缓解漂移结果。

在应对几何退化与长时一致性方面，LiDAR-SLAM 与 UWB 的融合被证明是有效途径。Park 等人\cite{liro2021} 提出的 LIRO 将 UWB 测距与 LiDAR–IMU 紧耦合于里程计前端，在滑窗内最小化点云几何残差与测距残差，以增强走廊、空旷等弱约束场景的可观测性；Nguyen 等人\cite{viral2021} 提出的 VIRAL-SLAM 则在因子图框架内统一相机、IMU、LiDAR 与 UWB 的两阶段紧耦合建模，并结合回环与全局图优化，显著降低纯视觉/纯激光在光照极端或结构稀疏下的漂移。更早的协作式方法\cite{uwb_lidar_coop_slam2012} 已验证“UWB 测距锚定 + 几何匹配”在多机相对定位与地图对齐上的可行性。

在多机与异构协同场景中，VIO+UWB 的协同定位与致密重建展示了“地面–空中”信息互补的潜力：地面平台（或固定锚点）提供中尺度几何约束，无人机提供高频自运动与局部稠密感知，二者联合完成相对定位与场景建模\cite{queralta2022vio}。面向“少锚/无锚”设置，Zhang 等人\cite{zhang2022anchorfree} 给出锚点自由的 UWB/INS 相对定位方案，Papadimitriou 等人\cite{papadimitriou2022rangeaid} 则在协作位姿估计中引入“距辅助”的多机器人相对位姿图模型。针对空中集群，Xu 等人提出去中心化的全向视觉–惯性–UWB 状态估计系统\cite{xu2022omni}，并在后续工作中发展为面向空中集群的分布式 LIO 方案 Swarm-LIO/Swarm-LIO2\cite{zhu2022swarmlio,zhu2024swarmlio2}，如 \autoref{fig:swarm-lio2} 所示，各无人机在本体上执行高精度自定位，并通过目标检测、跨机轨迹匹配与因子图优化实现协作约束融合，展示了分布式 VIO/LIO–UWB 融合在大规模系统中的可扩展性。

\begin{figure}[htbp]
    \centering
    \includegraphics[width=.5\linewidth]{chapter_1/多机slam.jpg}
    \caption{面向空中集群的分布式 LIO}
    \label{fig:swarm-lio2}
\end{figure}

在更复杂的地下与灾害环境中，Chang 等人\cite{chang2022lamp2} 提出的 LAMP 2.0 为异构多机器人给出“可插拔前端与鲁棒后端”的多机 SLAM 方案，强调在通信受限与感知退化条件下的稳健性；Cramariuc 等人\cite{cramariuc2022maplab2} 的 maplab 2.0 则以多模态、多会话和多机器人映射为目标，重构模块化接口，支持视觉–惯性与非视觉地标、语义目标的统一地图表达，为异构系统间地图互操作提供了研究型平台。

综合来看，当前异构多机器人系统在定位与感知上仍存在共性瓶颈：视觉与 VIO 在低纹理与强机动条件下易退化，UWB 受 LoS 与锚点布局制约且存在 NLoS 偏差，LiDAR 在弱几何与远距稀疏回波场景中会退化；在多机场景中还叠加通信带宽与时延限制、跨平台时间同步与外参自标定难题。更为关键的是，现有感知与定位大多作为被动估计模块独立设计，鲜有工作将“可见性与连通性需求”显式纳入规划层进行联合建模与优化，从而难以主动维持对任务有利的观测条件与几何约束。

\subsection{无人系统协同规划研究现状}

在多机器人协同规划中，同构团队的研究多采用集中式或分布式框架，典型目标包括避障与安全约束、编队保持、覆盖控制及任务分配等。对于异构团队（如空地系统），由于不同平台在动力学约束、作业高度与感知能力上的差异，规划问题进一步复杂化：UAV 具备高机动性与高位姿可达性，UGV 则受地形与非完整约束限制。要充分发挥空地系统在信息获取与任务执行上的互补优势，协同规划必须在信息增益、队形稳定以及通信/感知连通维护之间进行权衡。

在“感知驱动”的协同规划框架中，核心问题可以概括为“由谁、在何处、何时获取关键环境信息”，并据此反推各载体的机动与分工策略。Li 等人\cite{li2024colag} 将 UAV 视作“移动感知服务器”，在空侧运行 SLAM 并共享局部地图与相对位姿估计；当 UGV 的碰撞风险或定位不确定性上升时，通过带时间窗的调度模型触发空侧支援，以压缩等待时间与额外航程。Chen 等人\cite{chen2024costeffective} 在“少传感器、近距离协同”的条件下提出 SEAL（空侧顺序探索与定位援助）与 CAT（地侧自适应避碰轨迹）两级策略，使仅依赖里程计与相对位姿的 UGV 仍能完成稳定覆盖与编队。为系统化上述思路，典型做法将“感知服务”形式化为组合优化问题：空侧基于任务价值选择援助序列，地侧在安全与能量约束下规划会合与避障；调度层常借鉴带时间窗车辆路径问题（VRPTW）\cite{solomon1987vrptw}，在时限与资源约束下优化支援次序与停靠点分配。

\begin{figure}[htbp]
    \centering
    \includegraphics[width=.9\linewidth]{chapter_1/空地感知驱动.jpg}
    \caption{Chen 等人\cite{chen2024costeffective}提出的空地系统规划方法示意图}
    \label{fig:costefficient}
\end{figure}

在“感知点位选择”方面，信息增益与“下一最佳视点”（Next-Best-View, NBV）类策略为空侧生成高价值视点序列提供了通用范式。Popovic 等人\cite{popovic2020ipp} 提出面向无人机主动感知的 IPP 框架，在占据概率地图上以熵减与互信息作为效用函数，结合传感器可见空间与飞行时间代价，采用滚动时域重规划生成高价值航迹；Charrow 等人\cite{charrow2015it} 则从信息论出发，将互信息梯度融入轨迹优化，在统一优化问题中权衡“信息获取效率—机动代价—动力学可达性”；Bircher 等人\cite{bircher2016nbv} 通过在线构建有限深度的视点树结构，评估各候选视点的熵减收益并采用滚动式执行策略，仅执行首段轨迹后基于更新地图重复规划，从而兼顾在线性与探索效率。在“角色分配”层面，多机器人信息获取中的 Leader–Follower 架构\cite{dicarro2021mripp} 将强感知平台派往信息密度高的区域执行 IPP/NBV 规划，弱感知平台在其“感知带宽与覆盖阴影”下完成局部任务与安全会合，使得在传感资源不均的条件下整体信息增益最大化。概括而言，协感知驱动的协同规划通常可抽象为“信息采集—机动同步”的两层耦合结构：上层以地图熵、互信息或目标可探测概率度量“感知价值”，生成空侧候选视点与受援对象排序；下层在滚动时域或两阶段规划框架中，将上述决策落实为可执行的多机器人轨迹与会合时程。通过引入带时间窗调度与在线重规划，可在感知预算稀缺或分布不均时显著降低队形停滞和频繁全局重规划的需求，在不增加硬件投入的前提下提升系统整体效率与鲁棒性。

当感知受遮挡且需跨平台协同时，规划中必须显式考虑视线互观（Line of Sight, LoS）问题。Amigoni 等人\cite{amigoni2017survey} 在综述受通信限制环境下的多机器人探索时指出，仅以距离近似连通性难以应对障碍物遮挡。连续视线连通要求系统在整个运动过程中保持网络连通性。Giordano 等人\cite{giordano2013passivity} 以图拉普拉斯的 Fiedler 特征值为指标，构造基于势函数的分散控制律，并以“LoS 线段到最近障碍的距离”来刻画遮挡，从控制层保证连通性；Nestmeyer 等人\cite{nestmeyer2017decentralized} 在此基础上实现了“多目标去向—保持连通”的并行分布式探索。为便于嵌入优化，Yang 等人\cite{yang2023mvee} 采用最小体积包络椭球（MVEE）对 LoS 线段进行解析近似，获得可微且封闭式的约束形式；Chen 等人\cite{chen2023multiuavlos} 则在多 UAV 投放中以分离超平面定义安全区，将互见性转化为位置可行域，从而在障碍密集场景中保障 LoS。与此相对，离散视线连通仅要求在特定时间步保持连通。Hollinger 等人\cite{hollinger2012periodic} 提出的周期连通框架首先生成满足时序约束的连通拓扑，再在其上进行机动规划；Stump 等人\cite{stump2011visibility} 通过对环境进行多边形分解，将邻近机器人约束在公共多面体内以确保互见；Dutta 等人\cite{dutta2019continuousconn} 在信息采集任务中显式加入连续连通约束，使路径在信息增益与连通维护之间进行权衡；Xia 等人\cite{xia2023relink} 则提出两阶段方法：先在离散层搜索连通的目标拓扑，再在连续层对各机器人位姿进行局部细化，以兼顾里程与信息目标。

\begin{figure}[htbp]
    \centering
    \includegraphics[width=.75\linewidth]{chapter_1/realm.jpg}
    \caption{Bai 等人\cite{bai2025realm} 提出的 Realm 无人车协同规划方法示意图}
    \label{fig:realm}
\end{figure}

在未知环境中，依赖先验地图的可见性建模往往失效，需要直接从在线感知中度量 LoS。Bai 等人\cite{bai2025realm} 提出的 Realm 方法在任务代价中引入可微的 LoS 距离度量，基于占据表示实时计算协同梯度，使编队在执行跟踪任务的同时主动绕避遮挡并维持互观；控制律以分布式形式求解，仅需一跳通信即可扩展到多机。与仅基于通信半径或静态拓扑指标的做法相比，该类方法直接反映真实几何结构与临时遮挡，实现了由“感知驱动”向“可见性约束导航”的闭环过渡，为后续将 LoS 约束与任务规划深度耦合提供了基础。

\section{论文主要研究内容}
\subsection{研究内容与创新点}
本文围绕未知环境下空地系统的相对定位与协同规划问题展开研究。现有工作在总体思路上往往将“定位与感知”和“协同规划”进行松耦合处理：一方面，定位与感知侧普遍假设机器人之间始终具备良好的视线（LoS）条件，或依赖预先部署的外部基础设施，从而在方法设计时弱化了对障碍物遮挡、临时通信中断等真实情况的建模与应对；另一方面，规划侧通常采用基于距离的简化连通模型刻画系统组织结构，忽略了遮挡对互观质量与通信链路稳定性的直接影响，使得在障碍密集和动态不确定环境中难以同时维持足够的定位精度与协同效率。

从定位与感知的角度看，视觉或激光里程计在低纹理、弱几何结构或强机动场景中易发生退化，使用 UWB 测距又受非视距传播和锚点几何布局的制约而引入系统性偏差；不少协同定位方案默认能够持续获得跨平台的互相观测视野，但缺少对可见性与量测可用性的在线评估与补偿机制，一旦互观被遮挡或信道退化，误差便难以及时抑制，并在系统层面逐步累积。从协同规划的角度看，同构集群在覆盖、编队与探索等任务上已形成较为成熟的框架，但多停留在顶层目标驱动与简化连通约束层面；对于异构空地系统，由于动力学特性与感知能力存在差异显著，已有工作多采用固定的角色分工与预定义协作流程，缺乏根据实时定位情况与感知状况进行动态重分工、会合调度与轨迹修正的能力。

综上，目前尚缺乏一套贯通“定位—感知—规划”的中间协同机制，将定位质量、可见性水平与感知供给等低层状态以可解释的形式变换为规划层的约束与目标，指导在线协同决策；同时也缺乏利用规划结果反向服务于定位与通信，通过保持 LoS、提升几何编队结构一致性与稳固通信链路来抑制误差扩散的闭环设计。针对上述问题，本文面向未知环境中的空地系统，围绕“动态 UWB 支撑的相对定位”和“显式可见性约束的协同规划”两个核心环节开展研究，主要内容与创新点体现在以下三个方面：

\begin{itemize}
    \item \textbf{构建基于移动 UWB 基站的空地相对定位框架。}
    针对传统 UWB 依赖固定基站且需精确标定的局限，本文提出一种基于移动 UWB 基站的空地相对定位框架：将 UWB 基站部署于运行高精度 LiDAR-SLAM 的 UGV 之上，以 UGV 的 SLAM 全局坐标系作为系统参考，UAV 仅携带轻量化 UWB 标签，通过多基站进行一对一测距。为提高测距频率并降低标签功耗，使用高效多基站 TWR 协议，实现百赫兹级测距数据收集；在多基站几何解给出的位置初值基础上，采用误差状态扩展卡尔曼滤波（Error-State Extended Kalman Filter, ESKF）将 IMU 与位置初值结果进行数据融合，输出与 UGV 全局地图一致的高频、无偏 UAV 相对位姿。与此同时，利用多个 UGV SLAM 生成的栅格占据图实现空地间的高效地图共享与协同感知。实验表明，该框架在保证全局一致性的前提下兼顾高频率与低功耗，为后续协同规划提供稳定可靠的状态与环境输入。
    
    \item \textbf{提出基于点云聚类的可见性驱动协同规划方法。}
    在稳定相对定位的基础上，本文提出基于点云聚类的可见性驱动协同规划方法：首先利用基于密度的带噪声空间聚类算法（Density-Based Spatial Clustering of Applications with Noise, DBSCAN）\cite{dbscan} 对环境点云进行聚类，并用最小包络椭圆柱对每一簇进行凸近似，将复杂障碍物统一表示为适于优化求解的凸几何体；在此基础上，结合最小包络椭圆柱与离散控制障碍函数（Discrete Control Barrier Function, DCBF）\cite{dcbf} 建立几何避障约束，保证轨迹的碰撞安全；并围绕椭圆柱构造视线连通质量（LoS）优化目标，将空地系统间的互观比例显式纳入代价函数。规划过程中采用 A* 获得局部目标点，将地面编队结构与运行能耗一并纳入非线性模型预测控制（Nonlinear Model Predictive Control, NMPC）的滚动时域优化框架中。经仿真与实物实验，结果表明该方法能够实时生成同时满足避障、视线连通与能耗约束的协同轨迹，且由于障碍物采用凸近似，优化问题具有良好的可解性与数值稳定性。
    
    \item \textbf{提出基于序列可视空间的两阶段空地协同规划架构。}
    为增强在非结构化与高密度障碍场景下的泛化性与稳定性，本文进一步提出基于序列可视空间的两阶段协同规划方法。首先采用混合 A* 搜索满足无人机动力学约束的引导路径，并在引导路径的若干采样点处构建单点可视空间，串联形成可视走廊以刻画通路周围的可见域；随后将上述的单一 NMPC 问题分解为全局层与局部层：全局层负责生成带时间戳的同步路径点，局部层各机器人基于接收的路径点采用 MINCO 参数化方法生成满足动力学与避障约束的最小能量局部轨迹。在局部层引入障碍物“碰撞力”估计，在不破坏路径点时间同步性的前提下自适应调整轨迹。运动控制层中，针对 UGV 和 UAV 的动力学特性分别设计前馈-反馈轨迹跟踪器。该两阶段架构显著降低了单次优化规模与计算负担，提高了在线可解性，并在仿真与实物验证中展现出对复杂环境的更强适应性与协同稳定性，同时保持良好的定位精度与能耗优化效果。
\end{itemize}

% \begin{itemize}
%     \item 针对传统 UWB 依赖固定基站且需精确标定的局限，本文提出将 UWB 基站动态部署于运动稳定的地面机器人上，以地面机器人运行高精度 LiDAR-SLAM 的全局坐标系作为系统参考。系统中每架 UGV 搭载激光雷达完成局部 SLAM 并在初始化阶段对多车位姿进行一致化；UAV 端仅携带轻量化 UWB 标签用于与多基站的一对一测距。为提高测距频率与降低标签能耗，本文设计并实现一种高效的多基站 TWR（Two-Way Ranging）方法，减少标签发送次数并实现毫秒级测距数据汇聚，使测距频率可达百赫兹量级。在获得多基站—标签的测距数据后，本研究构建 Tri-ESKF 相对位姿估计方法：先基于多基站几何解（多边定位）获得初值，再以误差状态扩展卡尔曼滤波（ESKF）融合 IMU 与测距残差实现高频、无偏的位姿估计，最终输出与 UGV LiDAR-SLAM 全局地图一致的 UAV 相对位姿。此外，基于 UGV-SLAM 产出的栅格占据图实现空地间的高效地图共享与协同感知，保证定位与地图的一致性与可追溯性。经过实物实验，该框架兼顾高频率、低功耗与全局一致性。
%     \item 在获得稳定的相对定位输入后，本文提出一种基于点云聚类的可见性驱动协同规划方法。方法首先采用 DBSCAN 对场景点云进行聚类，将每一簇障碍物用最小包络椭圆柱进行凸几何拟合并提取其参数化表示，从而把复杂的障碍物集合用若干凸几何体近似表达，便于优化求解。基于聚类得到的几何表示，构建两类约束：一是结合最小包络椭圆柱与离散控制障碍函数（Discrete Control Barrier Function，DCBF）建立的几何避障约束，保证轨迹的碰撞安全；二是以椭圆柱为基础构造的视线连通（LoS）优化目标，用以提高空地系统间的互观比例。路径生成层通过 A* 在当前栅格地图上搜索局部目标点并形成目标点约束；在编队层，根据 UWB 定位误差成因设计几何队形约束以增强相对定位可观测性。控制量的积分被作为能量消耗目标，促使生成轨迹既平滑又节能。总体将上述目标与约束纳入非线性模型预测控制（NMPC）框架，并采用滚动时域优化策略实现实时求解。该方法在结构化仿真场景中能够实时生成满足避障、视线连通与能耗约束的协同轨迹，且因障碍物表达为凸集合，优化问题具良好可解性与数值稳定性。
%     \item 为提高方法在非结构化、高密度障碍或障碍“粘连”场景下的泛化能力，本本将点云聚类的障碍表达替换为面向 single-view 可视性的可视走廊（Visual Corridor）表示，以更细粒度刻画从某位置出发的可见域并保证数值可解性。首先采用混合 A*（Kino-dynamic A*）进行动力学可行的引导路径搜索，得到满足动力学约束的粗路径；其次基于该引导路径在每个引导点处生成对应的可视空间，并将相邻可视空间串联形成序列可视走廊，刻画沿通路两侧的可见区域。基于此表示，将原先的单一 NMPC 问题分解为两阶段求解以降低计算负担并提升可行性：第一级（全局）以空地可见性、队形几何与任务要求为目标求解未来若干时刻的同步路径点（带时间戳）；第二级（局部）由每个机器人接收对应路径点并在本地高频运行局部规划器，局部规划器以所接收路径点作为轨迹骨架，采用 MINCO 轨迹参数化方法生成满足动力学和避障约束的最小能量轨迹。为了处理避障，局部规划层中引入障碍物碰撞力估计，在不破坏同步点时序与同步性前提下将轨迹从障碍区域推离。控制方面，对 UGV 采用前馈—反馈控制实现 XY 轴独立速度跟踪；对 UAV 利用其微分平坦性设计前馈—反馈跟踪器以保证对同步点的精确到达与鲁棒性。该两阶段架构在仿真与实物试验中显示出对复杂环境的更好适应性、更高的可解性与更稳定的协作性能，同时保持了相对定位精度与能耗优化效果。
% \end{itemize}

\subsection{论文总体结构}
对应上述的研究内容，本文完成了三项工作，具体如\autoref{fig:article_structure}所示。本文各章节的研究内容安排概述如下：

\begin{figure}[htbp]
    \centering
    \includegraphics[width=.9\linewidth]{chapter_1/论文总体结构.jpg}
    \caption{总体框架图}
    \label{fig:article_structure}
\end{figure}

第一章，绪论。本章首先介绍多机器人系统的发展背景与应用前景，重点讨论由无人机与无人车组成的空地机器人系统在复杂环境任务中的优势与需求；随后指出当前多机器人系统在相对定位、可见性维护与协同规划方面仍存在的关键短板。接着，从无人系统定位与感知、多机器人协同规划以及空地协同系统三个层面，对国内外相关研究工作进行了系统梳理与分析，总结了现有方法的代表性进展及其局限性。在此基础上，阐述本文的总体研究思路与主要创新点，并对全文的章节结构进行概述。

第二章，基于动态 UWB 基站的协同定位方法研究。本章针对传统 UWB 依赖固定基站且需精确标定的局限，提出将 UWB 基站动态部署于运动稳定的UGV之上，以运行高精度 LiDAR-SLAM 的 UGV 全局坐标系作为系统统一参考。在系统架构上，每台 UGV 搭载激光雷达并运行 SLAM，以相同位姿进行初始化；UAV 端仅携带轻量化的 UWB 标签，用于与多基站进行一对一测距。为提高测距频率并降低标签能耗，本文使用高效的多基站 TWR（Two-Way Ranging）测距方法，减少系统间消息发送次数，实现毫秒级测距数据汇总，使系统测距频率可达百赫兹量级。在获得与多基站的测距数据后，构建一套基于误差状态扩展卡尔曼滤波的相对位姿估计方法：先利用多边定位方法获得 UAV 在 UGV 定义的全局坐标系下的位置初值，再将机载 IMU 输出与位置初值信息在滤波框架内进行融合，输出与 UGV SLAM 地图一致的高频、无偏的 UAV 相对位姿。此外，利用 UGV-SLAM 生成的栅格占据图实现空地间的地图共享与协同感知，保证状态估计与地图表示的统一性。最后通过实验验证该框架在高频率、低功耗与全局一致性之间取得了效果，为后续协同规划提供了可靠的状态与环境信息输入。

第三章，基于点云聚类的协同规划方法研究。本章面向空地系统运行过程中因障碍物遮挡导致的视线连接被遮挡的问题，提出一种基于点云聚类的可见性驱动协同规划方法。首先采用 DBSCAN 对场景点云进行聚类，将每一簇障碍物用最小包络椭圆柱进行凸几何拟合并提取对应参数化表示，从而用有限个凸几何体近似复杂障碍物集合，便于写入优化问题。基于上述几何表示，构建两类关键约束与目标：一方面，将最小包络椭圆柱与离散控制障碍函数（Discrete Control Barrier Function, DCBF）\cite{dcbf} 相结合，构建几何避障约束以保证规划轨迹的碰撞安全；另一方面，以椭圆柱几何关系为基础构造视线连通（LoS）优化目标，将空地系统间的互观程度显式纳入代价函数，从规划层面提升可见性。然后采用混合 A* 搜索方法在当前栅格地图上搜索无碰撞的局部规划目标点，并形成目标点约束；针对地面编队问题，根据 UWB 定位误差成因设计几何队形约束以增强相对定位的计算结果可信度。控制量的积分被作为能量消耗目标，促使生成的轨迹既平滑又节能。最终，将上述目标与约束统一纳入非线性模型预测控制（Nonlinear Model Predictive Control, NMPC）框架中，采用滚动时域优化策略并借助 CasADi\cite{casadi} 与 IPOPT\cite{ipopt} 实现在线求解与控制信号分发。实验结果表明，该方法能够实时生成同时满足避障、视线连通与编队约束的高效空地协同轨迹。

第四章，基于序列可视空间的协同规划方法研究。本章针对空地系统在非结构化环境或障碍物“粘连”场景下，点云聚类表达易出现环境表示不精细的问题，将障碍表示替换为面向单视角可视性的可视走廊（Visual Corridor）表示，以更细粒度刻画从某位置出发的可见域并兼顾数值可解性。具体而言，首先采用混合 A* 进行满足无人机动力学约束的引导路径搜索，得到粗略但动力学可行的全局路径；随后在引导路径的若干采样点处生成对应的可视空间表示，并将连续可视空间串联形成可视走廊，描述沿通路两侧的可见区域分布情况。基于该空间表示，将原先统一求解的 NMPC 问题分解为两阶段以降低计算负担并提升可行性：第一阶段（全局层）以空地可见性、队形几何与跟踪引导轨迹为目标，利用 L-BFGS\cite{LBFGS-Lite} 优化器求解未来若干时刻多机器人的同步路径点；第二阶段（局部层）中，每个机器人接收对应同步路径点信息并在本地运行局部规划器，将同步路径点作为待跟踪的轨迹节点，采用 MINCO\cite{wang2022geometrically} 轨迹参数化生成满足动力学与避障约束的最小能量轨迹。为实现局部轨迹避障，局部规划层中引入基于障碍物分布的“碰撞力”估计，在不破坏同步点时序与同步性的前提下，将待优化轨迹从障碍区域推离。在运动控制实现上，对 UGV 采用前馈-反馈控制实现 XY 轴独立速度跟踪；对 UAV 利用其微分平坦性设计前馈-反馈轨迹跟踪器，以保证对全局同步路径点的精确到达与运行鲁棒性。该两阶段架构在仿真与实物试验中展现出对复杂环境更好的适应性、更高的求解成功率与更稳定的协作性能，同时保持了相对定位精度与能耗优化效果，并通过消融试验验证了各优化目标与约束设计的必要性。

第五章，总结与展望。本章对全文的研究工作和主要创新点进行系统总结，回顾本文在空地相对定位、可见性驱动协同规划与两阶段序列可视空间规划方面的主要贡献，分析当前工作的局限性，并对未来在多机器人协同感知与规划、动态任务分配以及更大规模异构系统应用等方面的研究方向进行展望。


% 强调移动基站定位结果
\chapter{基于动态UWB基站的协同定位方法研究}

\section{引言}
在多机器人系统中，精确、稳定的相对定位是实现协同感知与规划控制的基础环节。对于空地异构系统而言，UAV 和 UGV 在运动特性与传感配置方面存在显著差异，如何在动态、未知环境中实现两类机器人之间的高精度相对定位与统一坐标系对齐，是构建协同体系的关键技术难题。传统定位方案主要依赖固定式 UWB 基站或外部定位系统（如 GNSS 或光学运动捕捉系统），通常需要精确布设基站并进行静态标定，这不仅部署成本高、灵活性不足，而且难以适应任务区域频繁变化或环境结构复杂的场景。此外，当环境中存在障碍物遮挡或视线中断时，固定 UWB 网络的几何可观性会显著下降，测距误差被放大，进一步限制了系统的有效覆盖范围与实时定位性能。

为解决上述问题，本章提出一种基于动态 UWB 基站的空地协同定位方法。通过将 UWB 基站安装在地面移动平台上构建动态测距网络，使定位系统具备随队移动、灵活部署与高频测距的能力；地面 UGV 搭载高精度激光雷达运行 LiDAR-SLAM，用以维护全局一致的坐标系，无人机端仅配备轻量级 UWB 标签与 IMU 模块，在无需外部基础设施的条件下实现空地间相对位姿的实时估计。为提高系统更新频率并降低通信负载，本章使用高效的多基站 TWR 协议，减少测距过程中消息交互次数，实现多基站测距结果的同步汇总与低功耗运行，从而获得百赫兹量级的测距数据。

在此基础上，本章构建基于多边定位与 ESKF 的相对位姿估计方法：首先利用多边定位算法在 UGV 定义的全局坐标系下求解 UAV 的位置初值；随后将机载 IMU 输出与 UWB 测距残差在 ESKF 框架中进行紧耦合融合，得到与 UGV LiDAR-SLAM 地图对齐的高频、低偏差 UAV 相对位姿估计结果。该融合策略在抑制单一传感器退化影响的同时，显著提升了相对定位的稳定性与抗噪性，使空地系统在无外部基准的情况下仍能保持较高的定位精度与一致性。

在环境建图与感知共享方面，本章进一步利用 UGV 端 SLAM 生成并维护的栅格占据图作为共享环境表达载体，构建了空地间的地图同步机制，实现多源感知信息在统一坐标系下的整合与高效传输，为后续协同规划与任务决策提供一致、可复用的环境输入。通过实验验证，本文提出的动态 UWB 协同定位框架在运行中表现出高精度、高频率与较强鲁棒性的综合性能，能够有效支撑后续章节中关于可见性保持与协同规划的研究。

本章结构安排如下：\autoref{sec:uwb_range} 介绍基于 UWB 模块的测距方法，包括基于飞行时间（Time of Flight, ToF）的双向双边测距原理与多基站高效测距策略；\autoref{sec:ag_loc} 详细阐述基于移动基站的空地相对定位方法，包括多边定位算法及基于 ESKF 的位姿估计流程；\autoref{sec:env_share} 说明基于栅格占据图的共享环境感知机制与多机器人地图同步方法；\autoref{sec:localization_result} 给出定位实验与结果分析；\autoref{sec:summary_loc} 对本章内容进行总结。

\section{基于UWB模块的测距方法} \label{sec:uwb_range}

为实现空地协同系统中无人机与地面机器人的高精度相对定位，本文采用 UWB 进行无线测距。UWB 信号具有时间分辨率高、多径分量可分离、抗干扰能力强等优势，能够在短距离场景中提供厘米级的测距精度。由于无人机在任务执行过程中机动频繁、姿态快速变化，且空地平台之间难以保持精确的时钟同步，传统依赖静态基站和单向时间戳的测距方法在同步性和稳定性方面均存在不足。为此，本文在 UWB 测距框架下，结合飞行时间（Time of Flight, TOF）原理采用双边双向测距（Double-Sided Two-Way Ranging, DS-TWR）方法，通过双向多次时序交互消除时钟偏差的影响，并在此基础上扩展为多基站高效测距结构，以满足空地协同定位的实时性与能效要求。

\subsection{基于TOF的双边双向测距方法}
UWB 测距的核心思想是利用无线信号在空间传播的飞行时间估计节点间距离。设标签节点与基站节点间信号传播速度为 $c$（在空中近似为电磁波传播速度），单程传播时间为 $T_{\mathrm{prop}}$，则理想情况下有
\begin{equation}
    d = c \, T_{\mathrm{prop}},
    \label{eq:tof_basic}
\end{equation}
其中 $d$ 为两节点之间的真实距离。然而在实际系统中，标签与基站各自使用本地时钟计时，存在未知的时钟偏移与频率漂移，若仅通过单向或简单的双向通信难以直接得到准确的 $T_{\mathrm{prop}}$。DS-TWR 方法通过在两个节点之间往返交换多条消息，并将本地测得的时间戳打包传输，实现对时钟偏差的消除。

如\autoref{fig:twr_process}所示，设备 A（标签）与设备 B（基站）在一次 DS-TWR 过程中交换三条消息 $M_1$、$M_2$ 与 $M_3$。在该过程中，双方分别记录如下时间戳：
\begin{itemize}
    \item 设备 A：发送消息 $M_1$ 的时间 $T_{M1\_tx}$，接收消息 $M_2$ 的时间 $T_{M2\_rx}$，发送消息 $M_3$ 的时间 $T_{M3\_tx}$；
    \item 设备 B：接收消息 $M_1$ 的时间 $T_{M1\_rx}$，发送消息 $M_2$ 的时间 $T_{M2\_tx}$，接收消息 $M_3$ 的时间 $T_{M3\_rx}$。
\end{itemize}

\begin{figure}[htbp]
    \centering
    \includegraphics[width=0.8\linewidth]{chapter_2/tof示例.png}
    \caption{双边双向测距（DS-TWR）时序示意图}
    \label{fig:twr_process}
\end{figure}

由上述时间戳可以定义两次往返测量中的往返时长（round-trip time）与应答延时（reply delay）：
\begin{equation}
    \begin{aligned}
        T_{\mathrm{round}1} &= T_{M2\_rx} - T_{M1\_tx}, &
        T_{\mathrm{reply}1} &= T_{M2\_tx} - T_{M1\_rx}, \\
        T_{\mathrm{round}2} &= T_{M3\_rx} - T_{M2\_tx}, &
        T_{\mathrm{reply}2} &= T_{M3\_tx} - T_{M2\_rx},
    \end{aligned}
    \label{eq:twr_round_reply_def}
\end{equation}
其中 $T_{\mathrm{round}i}$ 由单端时钟测得，$T_{\mathrm{reply}i}$ 则由另一端时钟测得，因此单独使用任一量都不可避免地包含时钟偏移和漂移的影响。将信号传播过程表示为“单程传播时间 $T_{\mathrm{prop}}$ 与本端应答延时之和”，可得两次往返测量的理想时序关系：
\begin{equation}
    \begin{aligned}
        T_{\mathrm{round}1} &= 2 T_{\mathrm{prop}} + T_{\mathrm{reply}1}, \\
        T_{\mathrm{round}2} &= 2 T_{\mathrm{prop}} + T_{\mathrm{reply}2},
    \end{aligned}
    \label{eq:twr_round_reply_prop}
\end{equation}
即每一次完整的往返时长由两倍的单程传播时间和一次应答延时共同构成。

若直接由\autoref{eq:twr_round_reply_prop} 计算 $T_{\mathrm{prop}}$，仍然需要准确获知 $T_{\mathrm{reply}1}$ 与 $T_{\mathrm{reply}2}$，而这两个量本身受对端时钟偏移影响。DS-TWR 的关键在于利用两次往返测量的组合消除时钟偏差，并得到 $T_{\mathrm{prop}}$ 的闭式估计。首先将\autoref{eq:twr_round_reply_prop} 两式相乘，有
\begin{equation}
    \begin{aligned}
        T_{\mathrm{round}1} T_{\mathrm{round}2}
        &= \bigl(2T_{\mathrm{prop}} + T_{\mathrm{reply}1}\bigr)
           \bigl(2T_{\mathrm{prop}} + T_{\mathrm{reply}2}\bigr) \\
        &= 4T_{\mathrm{prop}}^{2}
           + 2T_{\mathrm{prop}}\bigl(T_{\mathrm{reply}1} + T_{\mathrm{reply}2}\bigr)
           + T_{\mathrm{reply}1} T_{\mathrm{reply}2}.
    \end{aligned}
    \label{eq:twr_round_product}
\end{equation}
将\autoref{eq:twr_round_product} 两边同时减去 $T_{\mathrm{reply}1} T_{\mathrm{reply}2}$ 得到
\begin{equation}
    T_{\mathrm{round}1} T_{\mathrm{round}2} - T_{\mathrm{reply}1} T_{\mathrm{reply}2}
    = 2T_{\mathrm{prop}}\bigl(2T_{\mathrm{prop}} + T_{\mathrm{reply}1} + T_{\mathrm{reply}2}\bigr).
    \label{eq:twr_round_minus_reply}
\end{equation}
另一方面，将\autoref{eq:twr_round_reply_prop} 左右两式相加可得
\begin{equation}
    \begin{aligned}
        T_{\mathrm{round}1} + T_{\mathrm{round}2} + T_{\mathrm{reply}1} + T_{\mathrm{reply}2}
        &= (2T_{\mathrm{prop}} + T_{\mathrm{reply}1})
         + (2T_{\mathrm{prop}} + T_{\mathrm{reply}2})
         + T_{\mathrm{reply}1} + T_{\mathrm{reply}2} \\
        &= 4T_{\mathrm{prop}} + 2\bigl(T_{\mathrm{reply}1} + T_{\mathrm{reply}2}\bigr) \\
        &= 2\bigl(2T_{\mathrm{prop}} + T_{\mathrm{reply}1} + T_{\mathrm{reply}2}\bigr).
    \end{aligned}
    \label{eq:twr_round_sum}
\end{equation}
将\autoref{eq:twr_round_minus_reply} 右端中的公因子
$\bigl(2T_{\mathrm{prop}} + T_{\mathrm{reply}1} + T_{\mathrm{reply}2}\bigr)$
用\autoref{eq:twr_round_sum} 表达，可解得单程传播时间的闭式估计为
\begin{equation}
    \widehat{T}_{\mathrm{prop}}
    = \frac{T_{\mathrm{round}1}\,T_{\mathrm{round}2} - T_{\mathrm{reply}1}\,T_{\mathrm{reply}2}}
    {T_{\mathrm{round}1} + T_{\mathrm{round}2} + T_{\mathrm{reply}1} + T_{\mathrm{reply}2}}.
    \label{eq:Tprop_hat}
\end{equation}
可见，\autoref{eq:Tprop_hat} 仅由两端各自测得的往返时长与应答延时构成，不显式依赖双方时钟的绝对偏移，从而在一定程度上抵消了时钟误差对 TOF 估计的影响。

最后，将\autoref{eq:Tprop_hat} 代入距离模型\autoref{eq:tof_basic}，即可得到单次 DS-TWR 过程中的测距估计值：
\begin{equation}
    \hat{d} = c \, \widehat{T}_{\mathrm{prop}} ,
    \label{eq:distance_hat}
\end{equation}
其中 $\hat{d}$ 为基于 DS-TWR 的 UWB 测距结果。实际系统中，$c$ 根据介质折射率进行有效传播速度标定，$\widehat{T}_{\mathrm{prop}}$ 则由硬件计时分辨率与同步精度共同决定。相比仅依赖单次往返时序的简单 TWR，DS-TWR 在无需全局时钟同步的前提下，能够更稳健地抑制时钟偏移与频率漂移引入的系统性误差，为后续的多基站定位与滤波融合提供了更可靠的测距结果。

\subsection{多基站高效型 TWR 测距方法}
在实际空地协同任务中，为在一轮测距周期内同时获得标签（UAV）到多个基站（UGV）的距离，若对每个基站分别执行一次完整的 DS-TWR，将导致通信开销与标签发射次数随基站数量线性增加，难以满足高动态场景下的高频更新需求。为此，本文在单基站 DS-TWR 的基础上使用了一种“单次轮询、多站顺序应答、单次结束帧”的多基站高效型 TWR 策略。该策略通过消息复用与时隙轮询，在一次测距周期内完成对 $K$ 个基站的测距，使整体消息数由传统的 $3K$ 条降至 $K+2$ 条，从而显著降低空口负载与标签侧发射次数，并将测距频率提升至百赫兹量级。其时序结构如\autoref{fig:multi_anchor_twr} 所示。

\begin{figure}[htbp]
    \centering
    \includegraphics[width=0.8\linewidth]{chapter_2/多基站测距.png}
    \caption{多基站高效型 TWR 测距方法时序示意图}
    \label{fig:multi_anchor_twr}
\end{figure}

多基站高效型 TWR 的消息流程与时隙调度可概括如下。首先，标签轮询：标签在时刻 $t^{\mathrm{A}}_{\mathrm{poll,tx}}$ 发送 \texttt{poll} 广播帧，宣布新一轮测距周期开始，所有基站在接收该帧后进入本轮测距流程。其次，基站顺序应答：第 $i$ 个基站在接收 \texttt{poll} 的时刻 $t^{\mathrm{B}_i}_{\mathrm{poll,rx}}$ 记录本地时戳，并在其预分配的时隙起始时间（相对于 \texttt{poll} 接收时刻延迟 $\Delta_i$）发送 \texttt{resp} 帧，其发送时刻记为 $t^{\mathrm{B}_i}_{\mathrm{resp,tx}}$，标签端接收该帧的时刻为 $t^{\mathrm{A}}_{\mathrm{resp,rx}}(i)$。为避免不同基站应答发生碰撞，$\Delta_i$ 由固定时隙长度 $T_{\mathrm{slot}}$ 与时间裕度 $\tau_{\mathrm{guard}}$ 构成：
\begin{equation}
    \Delta_i = (i-1)\,T_{\mathrm{slot}} + \tau_{\mathrm{guard}}, \quad i = 1,\dots,K.
    \label{eq:slot_delta}
\end{equation}
最后，标签结束帧：标签在时刻 $t^{\mathrm{A}}_{\mathrm{final,tx}}$ 发送 \texttt{final} 帧以结束本轮测距周期，各基站接收该帧的时刻记为 $t^{\mathrm{B}_i}_{\mathrm{final,rx}}$。标签可在 \texttt{final} 中携带自身关键时间戳，以便在标签端或基站端完成双向时间消歧与测距计算。

对任意基站 $i$，为保持与单基站 DS-TWR 一致的量化方式，定义两段“往返时长”（round-trip time）与两段“应答延时”（reply delay），分别对应“标签视角的往返”与“基站视角的往返”：
\begin{equation}
    \begin{aligned}
        T^{(i)}_{\mathrm{round}1} &= t^{\mathrm{A}}_{\mathrm{resp,rx}}(i) - t^{\mathrm{A}}_{\mathrm{poll,tx}}, &
        T^{(i)}_{\mathrm{reply}1} &= t^{\mathrm{B}_i}_{\mathrm{resp,tx}} - t^{\mathrm{B}_i}_{\mathrm{poll,rx}}, \\
        T^{(i)}_{\mathrm{round}2} &= t^{\mathrm{B}_i}_{\mathrm{final,rx}} - t^{\mathrm{B}_i}_{\mathrm{resp,tx}}, &
        T^{(i)}_{\mathrm{reply}2} &= t^{\mathrm{A}}_{\mathrm{final,tx}} - t^{\mathrm{A}}_{\mathrm{resp,rx}}(i).
    \end{aligned}
    \label{eq:round_reply_multi}
\end{equation}
其中，$T^{(i)}_{\mathrm{round}1}$ 与 $T^{(i)}_{\mathrm{round}2}$ 分别表示由标签和基站本地时钟测得的两次往返时长，$T^{(i)}_{\mathrm{reply}1}$ 与 $T^{(i)}_{\mathrm{reply}2}$ 则分别为基站和标签侧的处理/应答延时。由于这些量完全由各自节点本地时钟测得，它们单独含有未知的时钟偏移与频率误差，但通过适当的组合可以消除时钟偏差对飞行时间估计的一级影响。

类似于单基站 DS-TWR 的推导（\autoref{eq:Tprop_hat}），对任意基站 $i$，可得到单程传播时间 $T^{(i)}_{\mathrm{prop}}$ 的闭式估计：
\begin{equation}
    \widehat{T}^{(i)}_{\mathrm{prop}}
    =
    \frac{
        T^{(i)}_{\mathrm{round}1}\,T^{(i)}_{\mathrm{round}2}
        - T^{(i)}_{\mathrm{reply}1}\,T^{(i)}_{\mathrm{reply}2}
    }{
        T^{(i)}_{\mathrm{round}1}
        + T^{(i)}_{\mathrm{round}2}
        + T^{(i)}_{\mathrm{reply}1}
        + T^{(i)}_{\mathrm{reply}2}
    }.
    \label{eq:Tprop_multi}
\end{equation}
\autoref{eq:Tprop_multi} 仅依赖于往返时长和应答延时的代数组合，不显式依赖节点间的绝对时钟偏移，因此在一定程度上抑制了时钟不一致对单程传播时间估计的影响。在获得 $\widehat{T}^{(i)}_{\mathrm{prop}}$ 后，对应的距离估计可写为
\begin{equation}
    \hat{d}_i = c\,\widehat{T}^{(i)}_{\mathrm{prop}},
    \label{eq:dist_multi}
\end{equation}
其中 $c$ 为电磁波在当前介质中的有效传播速度。通过对 $i=1,\dots,K$ 逐一进行同样的计算，标签在一个测距周期内即可同时获得其到所有基站的距离估计 $\hat{d}_i$，为后续多边定位提供距离测量值。

从时间开销角度来看，设本轮包含 $K$ 个基站，应答阶段最大片长近似为 $K\,T_{\mathrm{slot}}$，则单周期时长可估计为
\begin{equation}
    T_{\mathrm{cycle}}
    \approx
    T_{\mathrm{air}}(\texttt{poll})
    + K\,T_{\mathrm{slot}}
    + T_{\mathrm{air}}(\texttt{final})
    + T_{\mathrm{guard}},
    \label{eq:cycle}
\end{equation}
对应的定位更新频率为
\begin{equation}
    f_{\mathrm{upd}} \approx \frac{1}{T_{\mathrm{cycle}}} ,
    \label{eq:freq}
\end{equation}
其中 $T_{\mathrm{air}}(\cdot)$ 表示对应帧在空口传输的时间，$T_{\mathrm{guard}}$ 为整体防碰撞与防误差的安全时间裕度。通过合理设计时隙长度 $T_{\mathrm{slot}}$ 与时间裕度 $T_{\mathrm{guard}}$，在保证基站应答不发生冲突的前提下，可以将 $f_{\mathrm{upd}}$ 提升至百赫兹量级，从而满足空地协同系统在高动态飞行与机动条件下对高频定位更新的需求。

综上，多基站高效型 TWR 在保持对称 DS-TWR 时钟消偏特性的基础上，通过单次轮询与顺序时隙应答机制，实现了测距消息数量与标签侧发射次数的显著压缩，在空口带宽受限和机载功耗受限的条件下，仍能提供高频率、多基站的测距能力，为后续基于动态 UWB 基站的空地协同定位奠定了可靠的测距基础。

\section{基于移动基站的空地相对定位方法} \label{sec:ag_loc}

基于上一节得到的 UWB 测距结果，假设在统一全局坐标系下已知每辆 UGV 在当前测量时刻的位姿，由此可以计算出其上 UWB 基站的三维坐标集合 $\{\mathbf{a}_i\in\mathbb{R}^3\}_{i=1}^M$，其中 $M$ 为可用 UWB 锚点的数量；UAV 上部署的标签则获得与各基站之间的距离量测 $\{\tilde d_i\}_{i=1}^M$。本节在此基础上估计 UAV 在同一全局坐标系中的位置初值 $\mathbf{p}\in\mathbb{R}^3$，作为后续滤波融合的状态输入。

\subsection{基于 UWB 的多边定位方法}

多边定位（Multilateration, MLAT）\cite{sidorenko2016multilateration} 的基本思想，是利用多个已知位置的锚点与未知点之间的距离观测构建几何约束，并由此求解未知点的位置。在本系统中，锚点坐标由 UGV 的 LiDAR-SLAM 给出，未知点为 UAV 上的 UWB 标签。忽略测距噪声的理想情况下，UWB 观测可写为一组球面约束：
\begin{equation}
    \bigl\|\mathbf{p} - \mathbf{a}_i\bigr\|_2 = \hat{d}_i, 
    \qquad i = 1,\dots,M,\quad M \ge 3,
    \label{eq:uwb_range_model}
\end{equation}
其中 $\hat{d}_i$ 为上一节得到的距离估计。在实际应用中，因为存在测距噪声与偏差\autoref{eq:uwb_range_model} 仅在统计意义上成立，后续将通过最小二乘法进行补偿。

为便于理解\autoref{eq:uwb_range_model} 所表达的几何关系，\autoref{fig:mlat_2d} 以二维情形示意了三锚点多边定位的基本原理：三个已知锚点位置 $\{\mathbf{a}_1,\mathbf{a}_2,\mathbf{a}_3\}$ 构成三个以 $\hat{d}_i$ 为半径的圆，它们的交点给出未知点 $\mathbf{p}$ 的候选位置。

\begin{figure}[htbp]
    \centering
    \includegraphics[width=0.7\linewidth]{chapter_2/多边定位.png}
    \caption{二维情况下基于三圆交会的多边定位示意图}
    \label{fig:mlat_2d}
\end{figure}

首先，考虑利用三个锚点构造解析解。当 $M \ge 3$ 时，可从 $\{\mathbf{a}_i,\hat{d}_i\}$ 中任取三个锚点及其观测 $(\mathbf{a}_1,\hat{d}_1)$、$(\mathbf{a}_2,\hat{d}_2)$、$(\mathbf{a}_3,\hat{d}_3)$，在三维空间中构造一组局部正交基。定义
\begin{equation}
    \begin{aligned}
        \mathbf{e}_x &= \frac{\mathbf{a}_2 - \mathbf{a}_1}{\bigl\|\mathbf{a}_2 - \mathbf{a}_1\bigr\|_2}, 
        & d &= \bigl\|\mathbf{a}_2 - \mathbf{a}_1\bigr\|_2, \\
        \alpha &= \mathbf{e}_x^\top(\mathbf{a}_3 - \mathbf{a}_1), 
        & \tilde{\mathbf{e}}_y &= (\mathbf{a}_3 - \mathbf{a}_1) - \alpha\,\mathbf{e}_x, \\
        \mathbf{e}_y &= \frac{\tilde{\mathbf{e}}_y}{\bigl\|\tilde{\mathbf{e}}_y\bigr\|_2}, 
        & \mathbf{e}_z &= \mathbf{e}_x \times \mathbf{e}_y, \\
        \beta &= \mathbf{e}_y^\top(\mathbf{a}_3 - \mathbf{a}_1), &&
    \end{aligned}
    \label{eq:mlat_basis}
\end{equation}
其中 $\mathbf{e}_x,\mathbf{e}_y,\mathbf{e}_z$ 为以 $\mathbf{a}_1$ 为原点建立的右手正交坐标系基向量，$d$ 为锚点 $1$ 与锚点 $2$ 之间的距离，$\alpha,\beta$ 则描述锚点 $3$ 在该局部坐标系中的投影分量。将未知位置 $\mathbf{p}$ 在该局部坐标系中表示为
\begin{equation}
    \mathbf{p} = \mathbf{a}_1 + x \,\mathbf{e}_x + y \,\mathbf{e}_y + z \,\mathbf{e}_z,
    \label{eq:p_local_coords}
\end{equation}
并将其分别代入关于 $(\mathbf{a}_1,\hat{d}_1)$、$(\mathbf{a}_2,\hat{d}_2)$、$(\mathbf{a}_3,\hat{d}_3)$ 的三个球面方程，可得到关于 $(x,y,z)$ 的代数方程组。根据三球交会的经典推导，可得
\begin{equation}
    \begin{aligned}
        x &= \frac{\hat{d}_1^{\,2} - \hat{d}_2^{\,2} + d^2}{2d}, \\
        y &= \frac{\hat{d}_1^{\,2} - \hat{d}_3^{\,2} + \alpha^2 + \beta^2}{2\beta} - \frac{\alpha}{\beta}\,x, \\
        z &= \pm \sqrt{\hat{d}_1^{\,2} - x^2 - y^2}.
    \end{aligned}
    \label{eq:xyz_three_spheres}
\end{equation}
若根号项为负，则说明三球在几何上并不相交，通常意味着锚点几何退化，如近同心、近共线或测距误差过大，此时应更换锚点组合或转而采用线性法。在 $\hat{d}_1^{\,2} - x^2 - y^2 \ge 0$ 时，\autoref{eq:xyz_three_spheres} 给出两个候选解，对应于三球交线上的上下两点。由\autoref{eq:p_local_coords} 可恢复三维空间中的两个候选位置：
\begin{equation}
    \mathbf{p}^{(1,2)} = \mathbf{a}_1 + x\,\mathbf{e}_x + y\,\mathbf{e}_y \pm z\,\mathbf{e}_z.
    \label{eq:two_candidates}
\end{equation}

在存在两个几何候选解的情况下，可引入第四个锚点 $(\mathbf{a}_4,\hat{d}_4)$ 来消除二义性。分别计算
\begin{equation}
    \varepsilon^{(k)} = \bigl|\bigl\|\mathbf{p}^{(k)} - \mathbf{a}_4\bigr\|_2 - \hat{d}_4\bigr|, \quad k=1,2,
    \label{eq:ambiguity_residual}
\end{equation}
并选择残差更小的候选解作为当前的估计位置。

除了解析三球交会方法外，当 $M \ge 4$ 时，还可以通过“平方差消元”将\autoref{eq:uwb_range_model} 线性化，得到线性最小二乘形式。对任意参考锚点索引 $k$，考虑锚点 $i$ 与锚点 $k$ 的两条球面方程，相减并展开可消去 $\|\mathbf{p}\|_2^2$ 项，从而得到
\begin{equation}
    2(\mathbf{a}_i - \mathbf{a}_k)^\top \mathbf{p}
    = \|\mathbf{a}_i\|_2^2 - \|\mathbf{a}_k\|_2^2 + \hat{d}_k^{\,2} - \hat{d}_i^{\,2},
    \qquad i \neq k.
\end{equation}
将上述线性约束对所有 $i\neq k$ 进行堆叠，可写为矩阵形式
\begin{equation}
    \mathbf{A}\,\mathbf{p} = \mathbf{b},
\end{equation}
其中 $\mathbf{A}$ 的每一行为 $2(\mathbf{a}_i - \mathbf{a}_k)^\top$，$\mathbf{b}$ 的对应元素为右端差值。若满足$\mathrm{rank}(\mathbf{A}) = 3$，则标准最小二乘解为
\begin{equation}
    \hat{\mathbf{p}}_{\mathrm{LS}}
    =
    \bigl(\mathbf{A}^\top \mathbf{A}\bigr)^{-1}\mathbf{A}^\top \mathbf{b}.
    \label{eq:ls}
\end{equation}
当锚点近共面导致 $\mathrm{rank}(\mathbf{A}) = 2$ 时，\autoref{eq:ls} 在三维空间中不再具有唯一解。此时可以先在锚点所在平面上求解平面内坐标 $(x,y)$，然后利用距平面上某一锚点的测距值恢复高度分量。例如，以第 $i$ 个锚点为参考，在平面坐标系中已知锚点坐标 $(x_i,y_i)$ 与未知点 $(x,y)$，则高度分量的两个候选值为
\begin{equation}
    z_i^{(\pm)} = \pm \sqrt{\hat{d}_i^{\,2} - (x - x_i)^2 - (y - y_i)^2},
\end{equation}
再通过多锚点之间的一致性检查选取合理的符号，从而恢复三维位置 $\mathbf{p}$。若锚点相对结构出现几乎共线的情况，则需要在重新规划 UGV 编队。

在测距噪声具有非均匀特性的情况下，简单的普通最小二乘法会对噪声较大的观测值赋予过高权重。此时可以引入对角权阵 $\mathbf{W} = \mathrm{diag}(w_i)$，以测距方差为依据设置权重，得到加权最小二乘解
\begin{equation}
    \hat{\mathbf{p}}_{\mathrm{WLS}}
    =
    \bigl(\mathbf{A}^\top \mathbf{W} \mathbf{A}\bigr)^{-1}
    \mathbf{A}^\top \mathbf{W} \mathbf{b}.
\end{equation}

综上所述，本文在实际应用中使用三球交会与加权线性最小二乘两类方法：当前锚点几何结构良好且可用测距数量适中时，优先采用三球交会给出解析候选解，并利用第四个锚点进行二义性解的消除；当出现几何退化、测距数量较多时，则退回到加权最小二乘法计算。整体多边定位流程可概括为\autoref{alg:uwb_mlat} 所示。在实际应用中，UWB 多变定位的结果在z轴上的误差较大难以直接使用，所以常结合激光高度测量计对高度进行修正，以提高z轴方向上的定位精度和稳定性。

\begin{algorithm}[htbp]
  \begin{algorithmic}
    \REQUIRE $\{\mathbf{a}_i\}_{i=1..M}$（由各 UGV 位姿插值得到的锚点坐标），$\{\hat{d}_i\}_{i=1..M}$（UAV 与各锚的 UWB 测距结果）
    \ENSURE $\hat{\mathbf{p}}$（UAV 位置估计）
    \STATE 预处理：剔除无效或明显异常的 $\hat{d}_i$，检测锚点几何是否存在严重退化问题，如同心、近共线或近共面
    \IF{$M \ge 4$ 且几何条件较好}
      \STATE 从中选取若干组三锚点组合，按\autoref{eq:mlat_basis}--\autoref{eq:two_candidates} 计算候选解 $\mathbf{p}^{(k)}$
      \STATE 利用额外锚点依据\autoref{eq:ambiguity_residual} 计算残差，消除二义性与部分退化解
    \ELSIF{$M = 3$ 且三球相交}
      \STATE 按\autoref{eq:mlat_basis}--\autoref{eq:two_candidates} 直接求解，并结合高度先验选取一解
    \ENDIF
    \IF{解析解不存在或可信度不足}
      \STATE 基于全部有效锚点构造线性方程 $\mathbf{A}\mathbf{p}=\mathbf{b}$
      \STATE 根据测距置信度构造权阵 $\mathbf{W}$，求解加权最小二乘解 $\hat{\mathbf{p}}_{\mathrm{WLS}}$
      \IF{锚点接近共面}
        \STATE 在平面内先解 $(x,y)$，再利用测距正根恢复高度分量 $z$
      \ENDIF
      \STATE 将 $\hat{\mathbf{p}}_{\mathrm{WLS}}$ 作为当前位置估计
    \ELSE
      \STATE 将解析解 $\mathbf{p}^{(k^\ast)}$ 作为当前位置估计
    \ENDIF
    \STATE \textbf{return} $\hat{\mathbf{p}}$
  \end{algorithmic}
  \caption{\label{alg:uwb_mlat}基于多边定位的 UAV 位置估计流程}
\end{algorithm}

\subsection{基于 ESKF 的无人机位姿估计}

在上一节多边定位输出的 UAV 世界坐标位置观测基础上，结合机载 IMU 提供的高频惯性测量，本文采用 ESKF \cite{eskf} 进行无人机的实时位姿估计。ESKF 通过使用名义状态与小误差进行建模，将姿态误差限制在线性化局部的最小参数空间，在四元数单位约束与大角度机动条件下仍能保持滤波的一致性；同时，其结构天然适配 IMU 高频预测与UWB 低频、异步绝对位置校正的融合模式，计算负载较低，适合实时应用。滤波最终输出的位姿定义在 UGV 的 SLAM 世界坐标系中，为后续共享地图与协同规划提供统一的状态输入。

首先给出状态与量测模型。对名义状态与误差状态进行分别建模，名义状态定义为
\begin{equation}
    \mathbf{x}
    =
    \bigl[\,
        \mathbf{p},\ 
        \mathbf{v},\ 
        \mathbf{q},\ 
        \mathbf{b}_a,\ 
        \mathbf{b}_g
    \,\bigr],
    \label{eq:eskf_state}
\end{equation}
其中 $\mathbf{p} \in \mathbb{R}^3$、$\mathbf{v} \in \mathbb{R}^3$ 分别为世界坐标系下的位置与速度，$\mathbf{q} \in \mathbb{H}$ 为从机体系到世界系的单位四元数，$\mathbf{b}_a,\mathbf{b}_g \in \mathbb{R}^3$ 分别为加速度计与陀螺仪的偏移。IMU 量测模型写作
\begin{align}
    \tilde{\mathbf{a}} 
    &= \mathbf{R}(\mathbf{q})^\top\bigl(\dot{\mathbf{v}} - \mathbf{g}\bigr) + \mathbf{b}_a + \mathbf{n}_a, 
    \label{eq:imu_acc} \\
    \tilde{\boldsymbol\omega} 
    &= \boldsymbol\omega + \mathbf{b}_g + \mathbf{n}_g,
    \label{eq:imu_gyr}
\end{align}
其中 $\mathbf{R}(\mathbf{q})$ 为由四元数 $\mathbf{q}$ 对应的方向余弦矩阵，$\mathbf{g} = [0,0,-g]^\top$ 为世界系下的重力向量，$\tilde{\mathbf{a}}$、$\tilde{\boldsymbol\omega}$ 分别为加计与陀螺原始数据，$\mathbf{n}_a,\mathbf{n}_g$ 为高斯噪声，$\boldsymbol\omega$ 为机体系下的原始角速度。上一小节多边定位在时刻 $t_k$ 提供 UWB 位置观测，建模为
\begin{equation}
    \mathbf{z}_k^{\mathrm{uwb}}
    =
    \mathbf{p}(t_k) + \mathbf{n}_{\mathrm{uwb}}, 
    \qquad 
    \mathbf{n}_{\mathrm{uwb}} \sim \mathcal{N}(\mathbf{0},\mathbf{R}_{\mathrm{uwb}}),
    \label{eq:uwb_meas}
\end{equation}
其中 $\mathbf{R}_{\mathrm{uwb}}$ 为由经验标定得到的位置观测噪声协方差。

接下来给出名义系统方程及其离散传播形式。将 IMU 读数去偏后得到
\begin{equation}
    \mathbf{a}_m = \tilde{\mathbf{a}} - \mathbf{b}_a,\qquad
    \boldsymbol\omega_m = \tilde{\boldsymbol\omega} - \mathbf{b}_g,
    \label{eq:imu_debias}
\end{equation}
将其作为系统输入，名义状态的连续时间运动学模型为
\begin{align}
    \dot{\mathbf{p}} &= \mathbf{v}, 
    \label{eq:dyn_p} \\
    \dot{\mathbf{v}} &= \mathbf{R}(\mathbf{q})\,\mathbf{a}_m + \mathbf{g}, 
    \label{eq:dyn_v} \\
    \dot{\mathbf{q}} &= \tfrac{1}{2}\,\boldsymbol\Omega(\boldsymbol\omega_m)\,\mathbf{q},
    \qquad
    \boldsymbol\Omega(\boldsymbol\omega)
    =
    \begin{bmatrix}
        0 & -\boldsymbol\omega^\top \\
        \boldsymbol\omega & -[\boldsymbol\omega]_\times
    \end{bmatrix}\!,
    \label{eq:dyn_q} \\
    \dot{\mathbf{b}}_a &= \mathbf{n}_{ba},\qquad 
    \dot{\mathbf{b}}_g = \mathbf{n}_{bg},
    \label{eq:dyn_bias}
\end{align}
其中 $[\cdot]_\times$ 为将三维向量映射为反对称矩阵的算子，$\mathbf{n}_{ba},\mathbf{n}_{bg}$ 为零均值随机游走噪声。在实际实现中，采用零阶保持（ZOH）近似，在采样间隔 $\Delta t$ 内令输入常值，并对\autoref{eq:dyn_p}–\autoref{eq:dyn_bias} 进行数值积分，即可从时刻 $t_{k-1}$ 的后验名义状态 $\hat{\mathbf{x}}_{k-1|k-1}$ 得到时刻 $t_k$ 的先验名义状态 $\hat{\mathbf{x}}_{k|k-1}$。

误差状态用于捕捉名义状态与真实状态之间的小偏差，其定义为
\begin{equation}
    \delta\mathbf{x}
    =
    \bigl[\,
        \delta\mathbf{p},\ 
        \delta\mathbf{v},\ 
        \delta\boldsymbol\theta,\ 
        \delta\mathbf{b}_a,\ 
        \delta\mathbf{b}_g
    \,\bigr],
    \label{eq:error_state}
\end{equation}
其中 $\delta\boldsymbol\theta \in \mathbb{R}^3$ 为小角向量，用左乘扰动的方式表示姿态误差，即真实四元数可近似表示为
\begin{equation}
    \mathbf{q}_{\mathrm{true}}
    \approx
    \exp\bigl(\tfrac{1}{2}\delta\boldsymbol\theta\bigr) \otimes \mathbf{q},
\end{equation}
$\otimes$ 为四元数乘法。线性化后，误差状态的连续时间动力学可写为
\begin{equation}
    \dot{\delta\mathbf{x}}
    =
    \mathbf{F}\,\delta\mathbf{x} + \mathbf{G}\,\mathbf{n},
    \qquad
    \mathbf{n} = \bigl[\mathbf{n}_a,\ \mathbf{n}_g,\ \mathbf{n}_{ba},\ \mathbf{n}_{bg}\bigr],
    \label{eq:error_dyn}
\end{equation}
其中系统矩阵 $\mathbf{F}$ 与噪声输入矩阵 $\mathbf{G}$ 采用块矩阵形式：
\begin{equation}
    \mathbf{F}
    =
    \begin{bmatrix}
        \mathbf{0} & \mathbf{I} & \mathbf{0} & \mathbf{0} & \mathbf{0} \\
        \mathbf{0} & \mathbf{0} & -\mathbf{R}(\mathbf{q})[\mathbf{a}_m]_\times & -\mathbf{R}(\mathbf{q}) & \mathbf{0} \\
        \mathbf{0} & \mathbf{0} & -[\boldsymbol\omega_m]_\times & \mathbf{0} & -\mathbf{I} \\
        \mathbf{0} & \mathbf{0} & \mathbf{0} & \mathbf{0} & \mathbf{0} \\
        \mathbf{0} & \mathbf{0} & \mathbf{0} & \mathbf{0} & \mathbf{0}
    \end{bmatrix},
    \qquad
    \mathbf{G}
    =
    \begin{bmatrix}
        \mathbf{0} & \mathbf{0} & \mathbf{0} & \mathbf{0} \\
        -\mathbf{R}(\mathbf{q}) & \mathbf{0} & \mathbf{0} & \mathbf{0} \\
        \mathbf{0} & -\mathbf{I} & \mathbf{0} & \mathbf{0} \\
        \mathbf{0} & \mathbf{0} & \mathbf{I} & \mathbf{0} \\
        \mathbf{0} & \mathbf{0} & \mathbf{0} & \mathbf{I}
    \end{bmatrix}.
    \label{eq:F_G}
\end{equation}
设连续噪声协方差矩阵为 $\mathbf{Q}_c$，则在采样间隔 $\Delta t$ 内，误差状态的离散转移矩阵与离散噪声协方差可由
\begin{equation}
    \boldsymbol\Phi = \exp(\mathbf{F}\Delta t),\qquad
    \mathbf{Q}_d
    =
    \int_0^{\Delta t}
        \exp(\mathbf{F}\tau)\,\mathbf{G}\mathbf{Q}_c\mathbf{G}^\top
        \exp(\mathbf{F}\tau)^\top
    \,\mathrm{d}\tau
    \label{eq:disc}
\end{equation}
得到。协方差按标准预测公式传播：
\begin{equation}
    \mathbf{P}_{k|k-1}
    =
    \boldsymbol\Phi\,\mathbf{P}_{k-1|k-1}\,\boldsymbol\Phi^\top
    + \mathbf{Q}_d.
    \label{eq:covprop}
\end{equation}

当 UWB 多边定位在时刻 $t_k$ 提供\autoref{eq:uwb_meas} 所示的位置观测时，观测函数对误差状态的线性化形式为
\begin{equation}
    \mathbf{H}_k
    =
    \begin{bmatrix}
        \mathbf{I} & \mathbf{0} & \mathbf{0} & \mathbf{0} & \mathbf{0}
    \end{bmatrix},\qquad
    \mathbf{y}_k
    =
    \mathbf{z}_k^{\mathrm{uwb}} - \hat{\mathbf{p}}_{k|k-1},
    \label{eq:meas_lin}
\end{equation}
即仅对位置分量施加直接观测。根据扩展卡尔曼滤波的更新公式，首先计算创新协方差与卡尔曼增益：
\begin{equation}
    \mathbf{S}_k
    =
    \mathbf{H}_k \mathbf{P}_{k|k-1} \mathbf{H}_k^\top
    + \mathbf{R}_{\mathrm{uwb}},
    \qquad
    \mathbf{K}_k
    =
    \mathbf{P}_{k|k-1}\mathbf{H}_k^\top \mathbf{S}_k^{-1},
    \label{eq:gain}
\end{equation}
然后得到误差状态估计修正量
\begin{equation}
    \delta\hat{\mathbf{x}}_k = \mathbf{K}_k\,\mathbf{y}_k.
\end{equation}
ESKF 通过“误差注入”的方式更新名义状态：位置、速度与零偏采用向量加法，姿态采用小旋量的指数映射左乘原四元数：
\begin{align}
    \hat{\mathbf{p}}_{k|k} &= \hat{\mathbf{p}}_{k|k-1} + \delta\hat{\mathbf{p}}, \\
    \hat{\mathbf{v}}_{k|k} &= \hat{\mathbf{v}}_{k|k-1} + \delta\hat{\mathbf{v}}, \\
    \hat{\mathbf{q}}_{k|k}
    &= 
    \exp\!\bigl(\tfrac{1}{2}\delta\hat{\boldsymbol\theta}\bigr)
    \otimes \hat{\mathbf{q}}_{k|k-1}, \\
    \hat{\mathbf{b}}_{a,k|k} &= \hat{\mathbf{b}}_{a,k|k-1} + \delta\hat{\mathbf{b}}_a, \\
    \hat{\mathbf{b}}_{g,k|k} &= \hat{\mathbf{b}}_{g,k|k-1} + \delta\hat{\mathbf{b}}_g,
    \label{eq:inject}
\end{align}
其中 $\delta\hat{\mathbf{p}}$、$\delta\hat{\mathbf{v}}$、$\delta\hat{\boldsymbol\theta}$、$\delta\hat{\mathbf{b}}_a$、$\delta\hat{\mathbf{b}}_g$ 为 $\delta\hat{\mathbf{x}}_k$ 的各分量，$\exp(\cdot)$ 为 $\mathfrak{so}(3)$ 到 $\mathrm{SO}(3)$ 的指数映射。协方差采用 Joseph 形式更新，以保证数值上的对称性与正定性：
\begin{equation}
    \mathbf{P}_{k|k}
    =
    (\mathbf{I} - \mathbf{K}_k\mathbf{H}_k)\,\mathbf{P}_{k|k-1}\,(\mathbf{I} - \mathbf{K}_k\mathbf{H}_k)^\top
    + \mathbf{K}_k \mathbf{R}_{\mathrm{uwb}} \mathbf{K}_k^\top.
    \label{eq:joseph}
\end{equation}

在初始化与坐标一致性方面，系统首先利用静止阶段的 IMU 数据估计初始姿态：通过加速度计平均值估计俯仰与横滚角，偏航角则根据短时陀螺积分得到；初始速度设置为零，位置初始化为首次多边定位结果。UGV 的 SLAM 世界坐标系被选为统一参考坐标系，UWB 多边定位的结果和 ESKF 的状态估计均在该坐标系下表示，从而保证空地系统间坐标的一致性与可对齐性。

综上所述，本文基于误差状态扩展卡尔曼滤波构建了“IMU 高频预测 + UWB 多边定位低频校正”的空地相对位姿估计模块。该模块既利用了 IMU 在短时间尺度上的高动态响应，又通过 UWB 提供的绝对位置信息抑制长期漂移，实现了在 UGV SLAM 世界坐标系下的实时、高精度 UAV 位姿估计，为后续共享环境建图与可见性驱动协同规划提供稳定可靠的状态输入。

\section{基于栅格占据图的共享环境感知} \label{sec:env_share}

本节在统一世界坐标系下构建并共享三维栅格占据图（3D occupancy grid）。将经过 SLAM 处理并转换到世界系的点云投影到同一体素空间中，采用对数几率（log-odds）形式进行增量贝叶斯更新，从而实现跨平台、可增量同步的三维环境表示。栅格占据图结构简单，易于融合多种传感器输入，并且支持基于体素的可见性查询与碰撞检测，是后续可见性约束与协同规划的基础地图表示。

\subsection{点云到三维占据栅格的增量构建}

输入为经 UGV 端 SLAM 处理后、已转换到统一世界坐标系 $\{W\}$ 的三维点云
\begin{equation}
    \bigl\{\mathbf{p}^{W}\bigr\},
\end{equation}
设体素边长（环境分辨率）为 $r$，当前工作窗口在世界系中的下界坐标为 $\mathbf{l}$，则点 $\mathbf{p}^{W}$ 所属体素的整数索引可写为
\begin{equation}
    \mathbf{k}(\mathbf{p}^{W})
    =
    \Bigl\lfloor\frac{\mathbf{p}^{W}-\mathbf{l}}{r}\Bigr\rfloor
    \in \mathbb{Z}^3,
    \label{eq:voxel_index}
\end{equation}
其中 $\lfloor\cdot\rfloor$ 为按元素向下取整算子，$\mathbf{k} = [k_x,k_y,k_z]^\top$ 表示体素在三维栅格中的离散坐标。

对每个体素 $v_i$，用对数几率 $\ell(i)$ 表示其被占据的概率，记
\begin{equation}
    \ell(i)
    =
    \log\frac{p_i}{1-p_i},
\end{equation}
其中 $p_i$ 为体素 $v_i$ 被占据的概率。对数几率形式便于在多帧观测下进行增量相加更新。对被点云命中的体素，采用固定增量的更新策略：设 $\ell_{\mathrm{hit}}>0$ 为命中增量，$\ell_{\min},\ell_{\max}$ 为对数几率上下限，则对命中体素执行
\begin{equation}
    \ell(i)
    \ \Leftarrow\
    \mathrm{clip}\!\bigl(\,\ell(i) + \ell_{\mathrm{hit}},\ \ell_{\min},\ \ell_{\max}\bigr),
    \label{eq:hit_update_pc_final}
\end{equation}
其中 $\mathrm{clip}(\cdot)$ 表示在区间 $[\ell_{\min},\ell_{\max}]$ 上的饱和裁剪，以避免数值发散。

考虑到传感器安装高度与物体轮廓存在不确定性，直接将命中体素视为障碍可能低估实际占据范围。为增加安全裕度，本文在体素域中对命中体素进行邻域膨胀。记命中点 $\mathbf{p}^{W}$ 对应的体素索引为
\begin{equation}
    \mathbf{k}
    =
    \mathbf{k}(\mathbf{p}^{W})
    \in \mathbb{Z}^3,
\end{equation}
记任一邻域体素索引为 $\mathbf{q}\in\mathbb{Z}^3$，以无穷范数 $\|\cdot\|_{\infty}$ 在体素坐标中度量邻域大小。设 $d_{\mathrm{inf}}$ 为实际空间中的膨胀半径，体素域膨胀半径为
\begin{equation}
    n_{\mathrm{inf}}
    =
    \Bigl\lceil\frac{d_{\mathrm{inf}}}{r}\Bigr\rceil,
\end{equation}
则对所有满足
\begin{equation}
    \bigl\|\mathbf{q}-\mathbf{k}\bigr\|_{\infty} \le n_{\mathrm{inf}}
\end{equation}
的体素，将其标记为“膨胀占据候选”。在占据查询时，综合对数几率阈值与膨胀标记给出二值占据结果。记 $\ell_{\mathrm{occ}}$ 为占据阈值，则体素 $v_i$ 的占据判定为
\begin{equation}
    \mathrm{occ}(i)
    =
    \mathbb{I}\Bigl(
        \ell(i) \ge \ell_{\mathrm{occ}}
        \ \ \text{或}\ \
        v_i\ \text{被膨胀标记}
    \Bigr),
    \label{eq:occ_query_pc_final}
\end{equation}
其中 $\mathbb{I}(\cdot)$ 为指示函数。该判定结果在后续的路径碰撞检查与可见性计算中用作离散障碍表示。

为了在长时间运行中控制地图规模，系统采用滑动窗口的方式维护三维栅格：工作窗口随载体运动在世界系中平移，仅保留位于当前窗口内的体素，将离开窗口的体素内容清空，从而保持写入与查询复杂度近似恒定。此外，为适应环境变化并减弱陈旧观测的影响，周期性地对所有体素施加遗忘机制。设 $\delta>0$ 为遗忘步长，则对每个体素执行
\begin{equation}
    \ell(i)
    \ \Leftarrow\
    \max\bigl(\ell_{\min},\ \ell(i) - \delta\bigr),
    \label{eq:decay_pc_final}
\end{equation}
使得在长期未被新观测命中的区域，对数几率逐渐衰减到 $\ell_{\min}$，从概率意义上反映“环境可能已发生变化”的不确定性。

\subsection{多机器人环境感知的高效共享方法}

在多机器人场景下，如果直接在各载体间传输完整三维占据图，将导致通信负载过高，也不利于在有限带宽下维持地图一致性。为此，本文采用“局部子图裁剪”和“世界坐标系对齐融合”两步组合的共享机制。每辆 UGV 在统一世界坐标系下维护一张以自身当前位置为中心、边长为 $l$ 的局部占据栅格；UAV 端不独立构建全局地图，而是在同一世界系中按需接收并融合各 UGV 的局部子图，从而获取自身感知窗口内的最新环境信息。

\begin{figure}[htbp]
    \centering
    \includegraphics[width=0.8\linewidth]{chapter_2/地图共享.png}
    \caption{环境感知共享方法二维示意图}
    \label{fig:map_share_2d}
\end{figure}

为便于说明，如\autoref{fig:map_share_2d}所示，先考虑二维情形。记 UGV $i$ 与 UAV 的位置分别为
\begin{equation}
    \mathbf{p}^g_i(t),\ \mathbf{p}^a(t) \in \mathbb{R}^2,
\end{equation}
以无穷范数 $\|\cdot\|_{\infty}$ 定义以 UGV $i$ 为中心、边长为 $l$ 的方形感知区域：
\begin{equation}
    Q_i(t)
    =
    \Bigl\{
        \mathbf{p}\in\mathbb{R}^2\ \Big|\ 
        \bigl\|\mathbf{p}-\mathbf{p}^g_i(t)\bigr\|_\infty
        \le \tfrac{l}{2}
    \Bigr\}.
    \label{eq:Qi}
\end{equation}
类似地，以 UAV 为中心、边长为 $l$ 的感知区域为
\begin{equation}
    Q^a(t)
    =
    \Bigl\{
        \mathbf{p}\in\mathbb{R}^2\ \Big|\ 
        \bigl\|\mathbf{p}-\mathbf{p}^a(t)\bigr\|_\infty
        \le \tfrac{l}{2}
    \Bigr\}.
    \label{eq:Qa}
\end{equation}
记 UGV $i$ 在世界系下维护的局部二值占据图为
\begin{equation}
    M_i(\mathbf{p},t) \in \{0,1\},
\end{equation}
其中 $M_i(\mathbf{p},t)=1$ 表示位置 $\mathbf{p}$ 对应的栅格被判定为占据。为降低通信负载，UGV $i$ 仅在 UAV 落入其局部感知区域（$\mathbf{p}^a(t)\in Q_i(t)$）时，向 UAV 发送与 UAV 感知窗口重叠的子图。用集合的指示函数 $\mathbf{1}_\Omega(\mathbf{p})$ 表示 $\mathbf{p}$ 是否属于集合 $\Omega$，则 UGV $i$ 发送的子图可写为
\begin{equation}
    M_i^{s}(\mathbf{p},t)
    =
    M_i(\mathbf{p},t)\,
    \mathbf{1}_{\,Q_i(t)\cap Q^a(t)}(\mathbf{p}),
    \label{eq:submap}
\end{equation}
即仅保留 $Q_i(t)$ 与 $Q^a(t)$ 的交集区域。UAV 在自身对齐于世界坐标系的本地图上，将来自多辆 UGV 的子图进行逐点融合。对二值占据而言，采用“或”规则实现最大化占据：
\begin{equation}
    M_{\mathrm{fused}}(\mathbf{p},t)
    =
    \max_{\,i=1,\dots,N}
    M_i^{s}(\mathbf{p},t),
    \label{eq:binary_fuse}
\end{equation}
其中 $N$ 为参与协同的 UGV 数量。当同一栅格在不同时刻被不同 UGV 更新时，采用“时间戳优先”的策略，在冲突栅格处保留最新一次更新结果，以保证地图的一致性与时效性。

在三维情形下，可将\autoref{eq:Qi} 和 \autoref{eq:Qa} 中的平面方形区域推广为 $\mathbb{R}^3$ 中的边长为 $l$ 的立方体感知窗口，将连续空间 $\mathbf{p}$ 替换为对应的体素索引，实现体素级的子图裁剪与世界系对齐融合。UAV 以自身为中心维护一个三维立方体窗口，按需从各 UGV 拉取与该窗口重叠的局部子图，并在本地将其融合到统一的三维栅格中。借助这一机制，可以在显著降低带宽开销的前提下，维持 UAV 周围局部三维占据图的实时性、完整性和跨平台一致性，为后续可见性评估与协同规划提供可靠的环境输入。

\section{定位实验与结果分析} \label{sec:localization_result}

本节围绕前文提出的基于动态 UWB 基站的空地协同定位方法，开展对比实验，对系统在不同配置下的定位精度与稳定性进行定量评估。实验主要关注两类结果：一类是仅依赖 UWB 多边定位所得的原始位置估计 $\hat{\mathbf{p}}_{\mathrm{uwb}}(t)$，另一类是将多边定位结果与机载 IMU 通过 ESKF 融合后得到的滤波估计 $\hat{\mathbf{p}}(t)$ 与 $\hat{\mathbf{v}}(t)$。在固定基站场景中，通过运动标签实验，分析 UWB 测距误差与几何布局对原始定位精度的影响，并以真值轨迹为基准比较滤波前后的位置误差变化；在移动基站场景中，通过固定标签、UGV 携带基站运动的方式验证动态基站条件下多边定位的几何一致性与误差水平。综合上述结果，从测距残差、位置与速度误差统计以及误差随时间变化等方面，系统评估动态 UWB 基站与 ESKF 融合对提高空地系统相对定位精度和鲁棒性的作用。

\subsection{实验配置与评价指标}

本节在室内搭建 $8\times 8\times 3\,\mathrm{m}^3$ 的无遮挡实验场地，对所提出的动态 UWB 定位框架进行验证。所有实验均在统一的世界坐标系下进行，该坐标系由 UGV 端运行的 FAST-LIO2 \cite{xu2022fastlio2} SLAM 系统定义。UWB 基站统一安装于 UGV 车体上，其在世界坐标系下的位姿由 SLAM 实时估计；UAV 端在质心位置安装单个 UWB 标签，并采用飞控内置的 IMU 模块，IMU 采样频率约为 $200\,\mathrm{Hz}$，UWB 测距更新频率约为 $100\,\mathrm{Hz}$。系统内各数据源通过 ROS 框架完成时间同步，保证在同一时间轴上进行传感器融合与误差评估。

固定基站场景中，将四台 UGV 停放在场地四角，形成边长约为 $6\,\mathrm{m}$ 的正方形布设。四个 UWB 基站分别固定在四台 UGV 上，利用 FAST-LIO2 输出的位姿作为基站坐标，从而在世界系下获得 $\{\mathbf{a}_i\}_{i=0}^3$，基站按逆时针方向依次放置，此时 UGV 不运动，满足“静态锚点”条件。移动基站场景中，保持基站安装方式与放置顺序不变，由 UGV 携带基站在场地内按照预设轨迹运动，SLAM 实时给出 $\mathbf{a}_i(t)$。在这两类场景下，分别设计“标签静止、基站移动”和“标签移动、基站静止”的实验，用于考察固定锚点和动态锚点条件下的定位性能。单次实验期间持续记录 UWB 测距数据、UWB 多边定位结果以及基于 ESKF 的融合位姿估计。

真值位姿由动作捕捉系统提供，得到毫米级精度的位置轨迹 $\mathbf{p}^{\mathrm{gt}}(t)$。速度真值 $\mathbf{v}^{\mathrm{gt}}(t)$ 通过对真值位置进行数值差分获得。为了突出平面内定位性能，评价时重点考察 $x$、$y$ 轴上的误差。

在评价指标方面，针对不同定位方法（例如仅基于 UWB 多边定位的原始解以及基于 UWB–IMU 融合的 ESKF 结果），统一采用以下度量进行比较：一是给出估计轨迹与真值轨迹在平面内的可视化对比；二是统计平面内欧氏位置误差的时序变化、均方根误差（RMSE）与最大误差；三是在存在速度估计的场景下，进一步分析平面内速度误差及其 RMSE。后文中不再区分具体实现形式，统一用 $\hat{\mathbf{p}}(t)$ 表示被评估的估计位置轨迹，用 $\hat{\mathbf{v}}(t)$ 表示对应的速度估计。

\subsection{固定基站条件下的定位性能评估}

在固定基站条件下，为验证所提出 UWB 多边定位与 ESKF 融合方法的静态性能，在室内 $8\times 8\times 3\,\mathrm{m}^3$ 无障碍场地内搭建了四基站 UWB 测试环境。四个基站分别安装在四辆静止的 UGV 上，通过运行 FAST-LIO2 获得并固定其在统一 SLAM 世界坐标系下的位姿，基站几何布设为近似正方形，左下角为 $0$ 号基站，其余依逆时针编号为 $1$–$3$ 号。UAV 在质心处安装 UWB 标签与 IMU 传感器，通过动作捕捉系统获取毫米级精度的位置真值，并由位置差分得到速度真值。UWB 测距更新频率约为 $100\,\mathrm{Hz}$，IMU 频率为 $300\,\mathrm{Hz}$。考虑到 UWB/ESKF 估计与动作捕捉坐标系在平面上存在常偏移，评估时以所选时间段首帧为基准，将真值轨迹在 $xy$ 平面整体平移对齐到 ESKF 估计坐标系下，以消除坐标系偏差的影响。由于 $z$ 向高度由对地激光测距计单独提供且精度远高于平面定位，本节仅针对 $x$、$y$ 两个方向进行误差分析。

\begin{figure}[htbp]
  \centering
  \includegraphics[width=.95\linewidth]{chapter_2/fixed_xy_traj_2x2.pdf}
  \caption{固定基站条件下的 $xy$ 平面轨迹对比}
  \label{fig:fixed_xy_traj}
\end{figure}

在平面轨迹对比方面，\autoref{fig:fixed_xy_traj} 给出了在 $xy$ 平面上 UWB 多边定位结果、ESKF 融合结果以及对齐后的真值轨迹对比，轨迹颜色随时间增加逐渐变深。可以看到，在固定基站且无遮挡的条件下，UWB 多边定位已经能够较好地反映静止标签在安装抖动和环境噪声作用下产生的微小位移，其轨迹整体与真值基本重合，仅在局部存在厘米级抖动。ESKF 融合结果在平面上与 UWB 定位轨迹几乎重合，但曲线更加平滑，说明在当前工况下，IMU 预测主要起到对短时测距缺失的补偿与轨迹平滑作用，对整体平面位置精度的影响较小。

为更直观地展示误差水平，\autoref{tab:fixed_error_stats} 汇总了在所选时间段内各方法的位置与速度误差统计结果。可以看到，UWB 多边定位在 $xy$ 平面上的均方根误差（$\mathrm{RMSE}_{xy}$）约为 $2.97\,\mathrm{cm}$，最大平面位置误差约为 $13.8\,\mathrm{cm}$；ESKF 融合结果的平面均方根误差约为 $2.95\,\mathrm{cm}$，最大误差约为 $11.9\,\mathrm{cm}$，与 UWB 定位处于同一量级，并略有改善。这表明在基站几何布局良好且测距质量稳定的静态场景下，UWB 多边定位本身即可提供厘米级的平面位置精度，ESKF 主要用于增强时间连续性和后续动态场景下的鲁棒性。

\begin{figure}[htbp]
  \centering
  \includegraphics[width=.95\linewidth]{chapter_2/fixed_pos_vel_3sub.pdf}
  \caption{固定基站条件下的位置与速度对比}
  \label{fig:pos_vel_error_compare}
\end{figure}

在速度估计方面，由于 UWB 多边定位不直接输出速度，本节主要比较 ESKF 融合后的速度与真值速度之间的差异。\autoref{fig:pos_vel_error_compare} 的第二与第三张子图分别给出了 $v_x$ 与 $v_y$ 分量随时间的对比曲线，可以看到两者在整个实验时段内变化趋势高度一致，能够准确反映标签在快速运动和系统噪声作用下的速度变化。从定量结果看，$v_x$ 与 $v_y$ 分量的均方根误差分别约为 $0.031\,\mathrm{m/s}$ 和 $0.049\,\mathrm{m/s}$，最大速度误差幅值约为 $0.294\,\mathrm{m/s}$，与 \autoref{tab:fixed_error_stats} 中的统计结果一致，说明在固定基站、无遮挡环境中，ESKF 在保持厘米级位置精度的同时，能够提供平滑且与真值一致的平面速度估计。

\begin{table}[htbp]
    \centering
    \caption{固定基站条件下位置与速度误差统计}
    \label{tab:fixed_error_stats}
    \begin{tabularx}{\linewidth}{c *{5}{X<{\centering}}} 
        \toprule
        方法 & $\mathrm{RMSE}_{xy}$/cm & $e_{\max,xy}$/cm & $\mathrm{RMSE}_{v_x}$/m/s & $\mathrm{RMSE}_{v_y}$/m/s & $e^{v}_{\max}$/m/s \\ 
        \midrule 
        UWB 多边定位 & 2.97 & 13.8 & --    & --    & --    \\

        ESKF 滤波     & 2.95 & 11.9 & 0.031 & 0.049 & 0.294 \\ 
        \bottomrule
    \end{tabularx}
\end{table}

综上，在固定基站条件下，UWB 多边定位结合 ESKF 融合可以实现约 $10\,\mathrm{cm}$ 量级的平面位置精度和 $0.1\,\mathrm{m/s}$ 量级的速度精度，满足后续空地协同规划对静态标定和基准精度的需求，也为后续在移动基站和复杂环境下的性能评估提供了稳定参照。

\subsection{移动基站条件下的定位性能评估}

为评估移动基站场景下 UWB 多边定位与 ESKF 融合方法的鲁棒性，设计了“四车绕行”的实验：四辆搭载 UWB 基站与 LiDAR-SLAM 的 UGV 按照预设脚本沿正方形轨迹运动，UAV 端仅携带 UWB 标签与 IMU，标签在整个实验过程中保持静止。通过改变标签在平面中相对于 UGV 编队的位置，分别构造“标签处于基站构成多边形内部”和“标签处于多边形外部”两种几何工况，比较 ESKF 估计的平面位置在两种工况下的稳定性与误差水平。

\begin{figure}[htbp]
  \centering
  \begin{subfigure}[b]{0.48\linewidth}
    \centering
    \includegraphics[width=\linewidth]{chapter_2/moving_inside_xy_overview_grad_good.pdf}
    \caption{标签位于基站凸包内部}
    \label{fig:moving_xy_inside}
  \end{subfigure}
  \hfill
  \begin{subfigure}[b]{0.48\linewidth}
    \centering
    \includegraphics[width=\linewidth]{chapter_2/moving_inside_xy_overview_grad_bad.pdf}
    \caption{标签位于基站凸包外部}
    \label{fig:moving_xy_outside}
  \end{subfigure}
  \caption{移动基站条件下 UGV 与 UAV 在 $xy$ 平面的轨迹对比}
  \label{fig:moving_xy_inside_outside}
\end{figure}

\autoref{fig:moving_xy_inside_outside} 给出了两种工况下四辆 UGV 与 UAV 在 $xy$ 平面的轨迹分布情况。可以看到，在“包络内”工况中，四辆 UGV 沿正方形轨迹移动，而 UAV 的 ESKF 轨迹始终围绕静止标签的真实位置附近小范围波动，呈现出近似各向同性的微小扰动，且整体偏差有限。在“包络外”工况中，UGV 轨迹保持不变，而标签被布置在正方形外侧右方位置，此时 UAV 的 ESKF 轨迹在该参考点附近出现明显拉长的云团分布，并在靠近基站编队一侧方向上表现出更大的伸展，反映出几何精度稀释（GDOP）增大的影响。

\begin{figure}[htbp]
  \centering
  \begin{subfigure}[b]{0.9\linewidth}
    \centering
    \includegraphics[width=\linewidth]{chapter_2/moving_inside_x_vs_time_good.pdf}
    \caption{标签位于基站凸包内部}
    \label{fig:moving_x_inside}
  \end{subfigure}
  \hfill
  \begin{subfigure}[b]{0.9\linewidth}
    \centering
    \includegraphics[width=\linewidth]{chapter_2/moving_inside_x_vs_time_bad.pdf}
    \caption{标签位于基站凸包外部}
    \label{fig:moving_x_outside}
  \end{subfigure}
  \caption{移动基站条件下各机器人 $x$ 轴位置随时间的对比}
  \label{fig:moving_x_inside_outside}
\end{figure}

\begin{figure}[htbp]
  \centering
  \begin{subfigure}[b]{0.9\linewidth}
    \centering
    \includegraphics[width=\linewidth]{chapter_2/moving_inside_y_vs_time_good.pdf}
    \caption{标签位于基站凸包内部}
    \label{fig:moving_y_inside}
  \end{subfigure}
  \hfill
  \begin{subfigure}[b]{0.9\linewidth}
    \centering
    \includegraphics[width=\linewidth]{chapter_2/moving_inside_y_vs_time_bad.pdf}
    \caption{标签位于基站凸包外部}
    \label{fig:moving_y_outside}
  \end{subfigure}
  \caption{移动基站条件下各机器人 $y$ 轴位置随时间的对比}
  \label{fig:moving_y_inside_outside}
\end{figure}

为进一步分析平面位置随时间的变化规律，\autoref{fig:moving_x_inside_outside} 与 \autoref{fig:moving_y_inside_outside} 分别给出了两种工况下各机器人在 $x$ 轴和 $y$ 轴上的位置随时间演化曲线。对于“包络内”工况，UAV 在 $x$ 轴和 $y$ 轴上的估计值均稳定地围绕静止标签所在位置上下小幅波动，偏差幅度远小于 UGV 的编队运动尺度，表明在基站动态重构的过程中，多边定位与 ESKF 融合能够有效维持静止目标的平面位置一致性。而在“包络外”工况下，UAV 的 $x$ 轴与 $y$ 轴估计随时间呈现更为明显的偏移与缓慢漂移特征，特别是在与基站多边形连线主方向一致的轴向上波动更为剧烈，说明当标签脱离基站构成的凸包后，测距几何条件变差，对单次量测误差与基站运动的不确定性更加敏感。

\begin{figure}[htbp]
  \centering
  \begin{subfigure}[b]{0.85\linewidth}
    \centering
    \includegraphics[width=\linewidth]{chapter_2/moving_inside_uav_eskf_pos_error_good.pdf}
    \caption{标签位于基站凸包内部}
    \label{fig:moving_err_inside}
  \end{subfigure}
  \hfill
  \begin{subfigure}[b]{0.85\linewidth}
    \centering
    \includegraphics[width=\linewidth]{chapter_2/moving_inside_uav_eskf_pos_error_bad.pdf}
    \caption{标签位于基站凸包外部}
    \label{fig:moving_err_outside}
  \end{subfigure}
  \caption{移动基站条件下 UAV ESKF 位置误差随时间的对比}
  \label{fig:moving_err_inside_outside}
\end{figure}

在误差层面，\autoref{fig:moving_err_inside_outside} 比较了两种工况下 ESKF 估计在 $x$、$y$ 两个方向上的位置误差随时间的变化情况。可以观察到，当标签处于基站多边形内部时，$e_x(t)$ 与 $e_y(t)$ 始终围绕零均值小幅振荡，说明在动态锚点条件下仍可维持较高的定位稳定性；而当标签位于多边形外部时，$e_x(t)$ 与 $e_y(t)$ 的波动范围显著增大，且在某些时间段内出现持续的偏置与较大的瞬时峰值，表明此时多边定位对测距噪声与几何退化更为敏感，导致整体相对定位精度明显下降。

综合上述结果，在相同的 UGV 运动轨迹与测距配置下，让标签位于基站凸包内部能够显著改善移动基站场景下的几何可观测性，使得 UWB–ESKF 融合在基站大范围运动时仍能保持较为稳定的平面位置估计；而当标签脱离基站构成的凸包后，定位误差随时间放大并出现明显漂移。

\section{本章小结} \label{sec:summary_loc}

本章围绕空地协同系统中无人机与地面无人车之间的相对定位问题，构建了一套基于动态 UWB 基站的协同定位方法，并从测距建模、多边定位、惯性融合以及环境共享等多个层面进行了系统设计与实验验证。

首先，在\autoref{sec:uwb_range} 中，基于 UWB 信号的飞行时间（TOF）原理，给出了双边双向测距（DS-TWR）的时序关系与单程传播时间的闭式估计公式，并在此基础上设计了“单次轮询、多站顺序应答、单次结束帧”的多基站高效 TWR 测距策略。该策略通过统一轮询与时隙调度机制，在保持对称 DS-TWR 消除时钟偏差能力的前提下，显著减少了标签侧发射次数与总体消息数量，使系统在多基站配置下仍能达到百赫兹量级的测距更新频率，从而满足空地协同定位对高频测距与低能耗的需求。

随后，在\autoref{sec:ag_loc} 中，针对由多辆 UGV 搭载的 UWB 基站与 UAV 标签构成的移动测距网络，给出了基于多边定位的无人机位置初值求解方法，并在此基础上构建了基于 ESKF 的无人机位姿估计框架。在多边定位部分，通过多球交会与线性最小二乘两类方法相结合，显式分析了锚点几何布设对解的可观性与稳定性的影响。在惯性融合部分，引入“名义状态 + 误差状态”建模，将 IMU 高频预测与 UWB 低频绝对位置观测在统一的误差状态 EKF 中紧耦合，实现了在 UGV SLAM 世界坐标系下 UAV 实时位姿的连续输出，为后续协同规划提供了统一且高频的状态输入。

在\autoref{sec:env_share} 中，面向多机器人环境感知与地图共享问题，构建了基于三维栅格占据图的环境表示与同步机制。通过对数几率形式的体素增量更新、膨胀与遗忘机制，相比直接点云叠加能够在保证局部细节的同时抑制历史拖影；通过以统一世界坐标系为参考的“局部子图裁剪 + 世界系对齐融合”策略，在显著降低通信负载的前提下，实现了多辆 UGV 与 UAV 之间的占据地图共享，为后续的可见性约束建模、碰撞检测和栅格级 LoS 判定提供了统一环境载体。

最后，本章在固定基站与移动基站两类典型场景下，对所提出方法的定位性能进行了实验评估。在固定基站条件下，通过四基站静止布设与动作捕捉真值对比，验证了 UWB 多边定位在无遮挡场景中即可达到厘米级平面位置精度，ESKF 在此基础上进一步提供了平滑且与真值一致的速度估计，实现了约 $10\,\mathrm{cm}$ 量级的位置精度和 $0.1\,\mathrm{m/s}$ 量级的速度精度，满足后续协同规划对基准精度的要求。在移动基站条件下，通过四辆 UGV 围绕静止标签运动的实验，分别考察了“标签位于基站凸包内部”和“标签位于凸包外部”两种几何布设。结果表明，当标签处于基站构成的凸包内部时，即便基站大范围运动，UWB–ESKF 融合仍能维持较小的平面位置波动；而当标签脱离凸包后，平面位置误差与漂移显著增大。这一现象验证了动态基站场景中几何可观性的关键作用，也为后续章节在协同规划中引入可见性与定位约束提供了实验依据。

综上所述，本章完成了基于动态 UWB 基站的空地协同定位平台设计与实验验证，在统一世界坐标系下实现了高频、高精度的 UAV 相对位姿估计以及多机器人占据地图共享，为后续内容中基于可见性的协同规划方法提供了可靠的状态估计与环境建图基础。


\chapter{基于点云聚类的协同规划方法研究}

\section{引言}

在未知或半结构环境中实现空地多机器人系统的实时协同规划，核心难点主要体现在两个方面：一是如何在有限算力条件下获得既足够精细又便于约束建模的环境表达；二是如何在同一优化框架内同时处理避障、安全约束、编队保持以及跨平台相对定位等多重目标，并仍然满足在线求解的实时性要求。直接在稠密点云或高分辨率栅格上构建规划问题虽然具有较高几何精度，但会导致决策变量和约束数量急剧膨胀，使在线优化难以收敛或无法满足控制频率；而仅依赖几何启发式的局部方法，则往往难以在多机器人协同场景下兼顾相互约束与全局一致性。

针对上述矛盾，本文提出一种基于点云聚类的协同规划方法。核心思路是基于共享栅格占据图，对障碍点云进行聚类，并采用最小包络椭圆柱对每个聚类簇进行凸集近似，从而在保持障碍物几何轮廓主要特征的前提下，将环境表示压缩为若干低维参数（中心位置、长短轴、朝向及高度），显著降低优化问题的维度与复杂度，便于在滚动时域框架下进行快速求解。

在此基础上，本章首先在统一世界坐标系下使用基于密度的聚类算法 DBSCAN\cite{dbscan} 对栅格映射中的障碍点云进行聚类；随后对每个聚类簇提取二维最小包络椭圆，并沿高度方向扩展为三维椭圆柱，实现复杂障碍物的凸几何参数化表示。然后采用混合 A* 搜索方法在当前栅格地图上搜索无碰撞的局部规划目标点，并形成目标点约束；同时，在椭圆柱凸表示上构建离散控制障碍函数（Discrete Control Barrier Function, DCBF）\cite{dcbf}，形成具有良好可微性的几何避障约束。为体现空地协同特性，进一步在代价函数中引入编队几何与可见性相关指标，将 LoS 质量与相对定位可观测性纳入优化目标；能耗层面则以控制量平方和作为代价，鼓励生成平滑且节能的控制输入。整体问题采用非线性模型预测控制（Nonlinear Model Predictive Control, NMPC）进行建模，并结合滚动时域策略实现在线重规划，以适应环境变化与目标更新。

与直接在栅格占据图上进行规划相比，基于最小包络椭圆柱的障碍凸表示具有封闭形式的距离与法向信息，便于构造数值稳定的约束梯度，从而提升优化问题的可解性与实时性；与仅依赖几何启发的分步式方法相比，本章方法在单一优化框架内统一考虑了航路点约束、避障约束、编队保持与可见性目标，能够系统性地权衡安全性、效率与队形稳定性，并与前文构建的相对定位与共享地图模块自然衔接，形成“感知—定位—规划”的闭环结构。

本章余下内容安排如下：\autoref{sec:dyn_model} 建立 UAV 与麦克纳姆轮 UGV 的动力学模型，并给出统一的离散化形式；\autoref{sec:cluster_path} 介绍基于点云聚类的环境表达与混合 A* 引导路径生成方法；\autoref{sec:dcbf} 在椭圆柱凸表示上构建离散控制障碍函数与几何避障约束；\autoref{sec:traj_opt} 将编队几何、可见性指标与导航目标统一到 NMPC 优化模型中，并给出滚动时域求解策略；\autoref{sec:sim} 给出仿真实验设置与结果分析；最后在 \autoref{sec:summary_opt1} 对本章工作进行总结。

\section{无人系统的动力学建模与离散化} \label{sec:dyn_model}

本节作为建模与优化之间的衔接环节，对空地异构平台的动力学特性进行归纳，并给出统一的离散化表达，为后续优化问题提供一致的状态与控制变量描述。具体而言，首先分别在连续时间域建立 UAV 的微分平坦动力学模型与麦克纳姆轮 UGV 的平面全向运动模型；随后在统一建模层面，将 UAV 视作三维 $\{x,y,z\}$ 空间中的二阶积分系统，将 UGV 视作二维 $\{x,y\}$ 平面内的二阶积分系统，均采用“位置—速度”为状态，“加速度”为控制的形式，并在采样周期 $\Delta t$ 下进行离散化，同时施加不同平台特定的速度与加速度限幅约束，以兼顾轨迹的可执行性与优化的实时性。

\subsection{基于微分平坦的无人机动力学系统}

微分平坦系统的核心在于平坦输出（flat outputs）的存在：存在一个平坦输出向量 $\boldsymbol\sigma(t)$，使得系统状态 $\mathbf{x}(t)$ 与控制输入 $\mathbf{u}(t)$ 可以表示为 $\boldsymbol\sigma(t)$ 及其有限阶导数的代数函数，如下所示：
\begin{equation}
\mathbf{x}(t)
=F\big(\boldsymbol\sigma,\dot{\boldsymbol\sigma},\ldots,\boldsymbol\sigma^{(r-1)}\big),
\qquad
\mathbf{u}(t)
=G\big(\boldsymbol\sigma,\dot{\boldsymbol\sigma},\ldots,\boldsymbol\sigma^{(r)}\big),
\label{eq:flat_outputs}
\end{equation}
其中 $r$ 为系统相对阶。由此可以将原先在状态空间中的动力学与约束转化为平坦空间中的边界与光滑性约束，从而直接在平坦输出上进行轨迹规划与优化。

对于四旋翼 UAV，已有研究表明其具有微分平坦性质\cite{mellinger2011minimum}，常选质心位置与偏航角
\begin{equation}
\boldsymbol\sigma(t)
=\big[x(t),\,y(t),\,z(t),\,\psi(t)\big]^\top
\label{eq:sigma_uav}
\end{equation}
作为平坦输出。由 \autoref{eq:flat_outputs}–\autoref{eq:sigma_uav} 可重构出 UAV 的状态向量（位置 $\mathbf{p}$、速度 $\mathbf{v}$、姿态欧拉角 $\boldsymbol\xi=[\phi,\theta,\psi]^\top$ 等）以及控制输入（总推力 $T$ 与机体系角速度 $\boldsymbol\omega=[p,q,r]^\top$）。这意味着，只要在 $\{x,y,z,\psi\}$ 空间中生成满足一定光滑性与边界条件的轨迹，即可通过代数映射恢复一条可由真实 UAV 动力学跟踪的姿态与控制轨迹，并在平坦空间中直接对速度、加速度等物理量施加约束。

为了从平坦输出构造可执行的姿态与推力，首先根据质心加速度定义辅助向量
\begin{equation}
\mathbf{t}
=\big[\ddot x,\ \ddot y,\ \ddot z+g\big]^{\top},
\qquad 
\mathbf{z}_B
=\frac{\mathbf{t}}{\|\mathbf{t}\|},
\label{eq:zB_def}
\end{equation}
其中 $g$ 为重力加速度，$\mathbf{z}_B$ 为机体系 $z$ 轴在世界坐标系下的方向。给定偏航角 $\psi$，构造世界系下的参考单位向量
\begin{equation}
\mathbf{x}_C
=\big[\cos\psi,\ \sin\psi,\ 0\big]^{\top},
\label{eq:xC_def}
\end{equation}
再由
\begin{equation}
\mathbf{y}_B
=\frac{\mathbf{z}_B\times \mathbf{x}_C}{\|\mathbf{z}_B\times \mathbf{x}_C\|},
\qquad
\mathbf{x}_B
=\mathbf{y}_B\times \mathbf{z}_B,
\label{eq:body_basis}
\end{equation}
得到机体系的三个正交基向量 $\mathbf{x}_B,\mathbf{y}_B,\mathbf{z}_B$，从而构造姿态矩阵
\begin{equation}
{}^{W}\!R_{B}
=\big[\,\mathbf{x}_B\ \ \mathbf{y}_B\ \ \mathbf{z}_B\,\big],
\label{eq:Rwb}
\end{equation}
其中需要保证 $\mathbf{x}_C\times \mathbf{z}_B\neq \mathbf{0}$ 以避免姿态奇异。

总推力由平坦输出的加速度直接给出：
\begin{equation}
u_1
=m\|\mathbf{t}\|
=m\sqrt{\big(\ddot x\big)^{2}+\big(\ddot y\big)^{2}+\big(\ddot z+g\big)^{2}},
\label{eq:thrust}
\end{equation}
其中 $m$ 为 UAV 质量。角速度则由更高阶导数得到。令
\begin{equation}
\mathbf{a}
=\big[\ddot x,\ \ddot y,\ \ddot z\big]^{\top},\qquad
\dot{\mathbf{a}}
=\frac{\mathrm{d}\mathbf{a}}{\mathrm{d}t},
\end{equation}
定义中间量
\begin{equation}
\mathbf{h}_{\omega}
=\frac{m}{u_1}\Big(\dot{\mathbf{a}}-\big(\mathbf{z}_B^{\top}\dot{\mathbf{a}}\big)\mathbf{z}_B\Big),
\label{eq:how_def}
\end{equation}
则机体系角速度分量可写为
\begin{equation}
p=-\mathbf{h}_{\omega}\cdot \mathbf{y}_B,
\qquad 
q=\mathbf{h}_{\omega}\cdot \mathbf{x}_B,
\qquad 
r=\dot{\psi}\,\mathbf{z}_W\cdot \mathbf{z}_B,
\label{eq:ang_vel}
\end{equation}
其中 $\mathbf{z}_W=[0,0,1]^{\top}$ 为世界系竖直方向单位向量。进一步对 \autoref{eq:ang_vel} 求导即可得到角加速度 $\dot{\boldsymbol\omega}$，代入欧拉刚体方程
\begin{equation}
\dot{\boldsymbol{\omega}}
= \mathbf{I}^{-1}\!\left(-\,\boldsymbol{\omega}\times \mathbf{I}\boldsymbol{\omega}+
\begin{bmatrix}u_2\\ u_3\\ u_4\end{bmatrix}\right),
\label{eq:euler}
\end{equation}
其中 $\mathbf{I}$ 为转动惯量矩阵，即可将滚转、俯仰、偏航三个力矩输入 $\{u_2,u_3,u_4\}$ 表达为平坦输出及其导数的代数函数。至此，实现了“平坦输出 $\boldsymbol\sigma(t)$ $\to$ 姿态 ${}^{W}\!R_B$、总推力 $u_1$ 与力矩输入 $\{u_2,u_3,u_4\}$”的完整映射。

综上所述，只要在平坦输出空间为 $\{x(t),y(t),z(t),\psi(t)\}$ 构造足够光滑的参考轨迹，就可以通过 \autoref{eq:zB_def}–\autoref{eq:euler} 唯一恢复与四旋翼动力学一致的姿态与控制输入。后续在规划与优化层面，本文将利用这一微分平坦性，将 UAV 在 $\{x,y,z\}$ 空间中的运动抽象为以位置—速度为状态、加速度为控制的二阶积分模型；在优化中对加速度及其导数施加约束，并通过跟踪平坦输出轨迹保证可执行性，从而实现与真实 UAV 动力学的一致可跟踪性。

\subsection{基于麦克纳姆轮小车的动力学系统}

麦克纳姆轮底盘在理想无打滑条件下（滚子与地面夹角为 $45^\circ$）具备平面全向机动能力。以车体中心 $O$ 为原点建立车体系，如\autoref{fig:mecanum_kin} 所示，车体前进方向记为 $V_x$ 轴，车体左侧为 $V_y$ 轴，绕垂直于平面的转动角速度记为 $\omega_z$。则世界系下 UGV 的位姿与车体系速度可写为
\begin{equation}
\mathbf{x}
=
\begin{bmatrix}
x & y & \theta
\end{bmatrix}^\top,
\qquad
\mathbf{v}_B
=
\begin{bmatrix}
v_x & v_y & \omega_z
\end{bmatrix}^\top,
\label{eq:meca_state_only}
\end{equation}
其中 $x,y$ 为世界系平面坐标，$\theta$ 为车体朝向，$v_x,v_y$ 分别是沿图中 $V_x,V_y$ 方向的平移速度分量。

世界系与车体系速度满足刚体运动关系
\begin{equation}
\begin{bmatrix}
\dot x\\[2pt] \dot y\\[2pt] \dot\theta
\end{bmatrix}
=
{}^{W}\!R_B(\theta)\,
\begin{bmatrix}
v_x\\[2pt] v_y\\[2pt] \omega_z
\end{bmatrix},
\qquad
{}^{W}\!R_B(\theta)
=
\begin{bmatrix}
\cos\theta & -\sin\theta & 0\\
\sin\theta &  \cos\theta & 0\\
0          &  0          & 1
\end{bmatrix},
\label{eq:meca_world_kin_only}
\end{equation}
即只要给定车体系速度 \(\mathbf{v}_B\)，便可以得到世界系下的位姿变化率。

\begin{figure}[htbp]
    \centering
    \includegraphics[width=0.55\linewidth]{chapter_3/麦克纳姆轮小车动力学.jpg}
    \caption{麦克纳姆轮小车底盘及车体系示意图}
    \label{fig:mecanum_kin}
\end{figure}

结合\autoref{fig:mecanum_kin} 的几何参数，记车轮半径为 $r$，车体总长为 $H$、总宽为 $W$，则车体中心 $O$ 到每个车轮在前后方向和左右方向的距离分别为
\[
L_x = \tfrac{H}{2},\qquad L_y = \tfrac{W}{2},
\]
设 A、B、C、D 四个车轮的角速度为
\[
\boldsymbol\omega
=
\begin{bmatrix}
\omega_A & \omega_B & \omega_C & \omega_D
\end{bmatrix}^\top,
\]
其中 A 为左后轮，B 为左前轮，C 为右前轮，D 为右后轮。在常见的麦克纳姆轮滚子布置方式下，前向运动学可以写成
\begin{equation}
\begin{bmatrix}
v_x\\[2pt] v_y\\[2pt] \omega_z
\end{bmatrix}
=
\frac{r}{4}
\begin{bmatrix}
1 & 1 & 1 & 1\\[2pt]
-1 & 1 & 1 & -1\\[4pt]
-\dfrac{1}{L_x+L_y} & \dfrac{1}{L_x+L_y} & -\dfrac{1}{L_x+L_y} & \dfrac{1}{L_x+L_y}
\end{bmatrix}
\boldsymbol\omega.
\label{eq:meca_forward_only}
\end{equation}

由\eqref{eq:meca_world_kin_only}–\eqref{eq:meca_forward_only} 可以看出，在给定车体几何尺寸 $H,W$ 和车轮半径 $r$ 后，通过独立调节四个车轮角速度 \(\omega_A,\omega_B,\omega_C,\omega_D\)，即可在车体系中独立生成任意的平移速度分量 $v_x,v_y$ 以及转动速度 $\omega_z$。换言之，UGV 在平面上具有全向机动能力：一方面可以在保持 $\omega_z=0$ 的同时任意方向平移，另一方面也可以在保持 $v_x=v_y=0$ 的情况下原地旋转。

在后续的规划与优化中，为降低维度并方便与 UAV 统一建模，本文不再显式求解车轮角速度，而是直接在世界系中使用平移速度、加速度作为控制量。在给定合适的速度与加速度限幅后，上述麦克纳姆轮运动学保证了所规划的二阶积分模型轨迹在实际底盘上是可跟踪、可实现的。

\subsection{异构无人系统的统一动力学模型及其离散化}

为在同一优化问题中统一处理空地异构平台，本文在协同规划层面均采用以“位置—速度”为状态、以“加速度”为控制的二阶积分模型。对任一机器人 $r$，记其工作空间维度为 $d_r\in\{2,3\}$（UGV 取 $d_r=2$ 表示平面运动，UAV 取 $d_r=3$ 表示三维运动），在离散时刻 $k$ 定义
\begin{equation}
\mathbf{p}_k^{(r)}\in\mathbb{R}^{d_r},\quad
\mathbf{v}_k^{(r)}\in\mathbb{R}^{d_r},\quad
\mathbf{u}_k^{(r)}\in\mathbb{R}^{d_r},
\label{eq:unified_states}
\end{equation}
其中 $\mathbf{p}_k^{(r)}$ 和 $\mathbf{v}_k^{(r)}$ 分别为世界系位置与速度，$\mathbf{u}_k^{(r)}$ 为加速度控制输入。对连续时间二阶积分系统在采样周期 $\Delta t$ 下采用零阶保持（ZOH）离散化，可得
\begin{equation}
\mathbf{p}_{k+1}^{(r)}=\mathbf{p}_k^{(r)}+\Delta t\,\mathbf{v}_k^{(r)}+\tfrac{1}{2}\Delta t^2\,\mathbf{u}_k^{(r)},\qquad
\mathbf{v}_{k+1}^{(r)}=\mathbf{v}_k^{(r)}+\Delta t\,\mathbf{u}_k^{(r)}.
\label{eq:disc_double_integrator}
\end{equation}

将状态拼接为
\begin{equation}
\mathbf{x}_k^{(r)}=
\begin{bmatrix}
\mathbf{p}_k^{(r)}\\[2pt]
\mathbf{v}_k^{(r)}
\end{bmatrix}\in\mathbb{R}^{2d_r},
\end{equation}
则离散动力学可写为线性系统
\begin{equation}
\mathbf{x}_{k+1}^{(r)}=
\underbrace{\begin{bmatrix}
\mathbf{I} & \Delta t\,\mathbf{I}\\[2pt]
\mathbf{0} & \mathbf{I}
\end{bmatrix}}_{\displaystyle \mathbf{A}_r}
\mathbf{x}_{k}^{(r)}+
\underbrace{\begin{bmatrix}
\tfrac{1}{2}\Delta t^2\,\mathbf{I}\\[2pt]
\Delta t\,\mathbf{I}
\end{bmatrix}}_{\displaystyle \mathbf{B}_r}
\mathbf{u}_{k}^{(r)},
\qquad \mathbf{I}\in\mathbb{R}^{d_r\times d_r},
\label{eq:disc_AB}
\end{equation}
其中 $\mathbf{A}_r,\mathbf{B}_r$ 仅随维度 $d_r$ 与采样周期 $\Delta t$ 变化，可预先计算并在求解过程中复用。

为反映各平台的物理可执行性与安全裕度，在统一模型上对速度与加速度施加盒约束
\begin{equation}
\mathbf{v}_{\min}^{(r)}\ \le\ \mathbf{v}_k^{(r)}\ \le\ \mathbf{v}_{\max}^{(r)},\qquad
\mathbf{u}_{\min}^{(r)}\ \le\ \mathbf{u}_k^{(r)}\ \le\ \mathbf{u}_{\max}^{(r)},
\label{eq:box_constraints}
\end{equation}
其中上下界由 UAV 与 UGV 的最大速度、最大加速度以及安全系数给定，可针对不同平台分别设置。

在异构多机器人场景中，设预测域长度为 $H$（采样周期为 $\Delta t$），机器人集合为 $\mathcal{R}=\{1,\dots,M\}$。对任一 $r\in\mathcal{R}$，在 $k=0,\dots,H$ 上定义
\begin{equation}
\mathbf{x}_k^{(r)}=
\begin{bmatrix}
\mathbf{p}_k^{(r)}\\[2pt]
\mathbf{v}_k^{(r)}
\end{bmatrix}\in\mathbb{R}^{2d_r},\qquad
\mathbf{u}_k^{(r)}\in\mathbb{R}^{d_r},
\end{equation}
则由 \autoref{eq:disc_AB} 得到每个时间步上的动态等式约束
\begin{equation}
\mathbf{x}_{k+1}^{(r)}=\mathbf{A}_r\,\mathbf{x}_k^{(r)}+\mathbf{B}_r\,\mathbf{u}_k^{(r)},
\qquad k=0,\dots,H-1,\qquad \mathbf{x}_0^{(r)}\ \text{已知}.
\label{eq:dyn_mpc_recursion}
\end{equation}

为便于在优化器中统一处理，将未来 $H$ 步的控制与状态按时间堆叠为
\begin{align}
\mathbf{U}^{(r)}&=\big[\mathbf{u}_0^{(r)\top}\ \mathbf{u}_1^{(r)\top}\ \cdots\ \mathbf{u}_{H-1}^{(r)\top}\big]^\top\in\mathbb{R}^{H d_r},\\
\mathbf{X}^{(r)}&=\big[\mathbf{x}_1^{(r)\top}\ \mathbf{x}_2^{(r)\top}\ \cdots\ \mathbf{x}_{H}^{(r)\top}\big]^\top\in\mathbb{R}^{H\cdot 2d_r}.
\end{align}
在多机器人系统下，进一步定义全局堆叠向量
\begin{equation}
\mathcal{U}=\big[\mathbf{U}^{(1)\top}\ \cdots\ \mathbf{U}^{(M)\top}\big]^\top,\qquad
\mathcal{X}=\big[\mathbf{X}^{(1)\top}\ \cdots\ \mathbf{X}^{(M)\top}\big]^\top,
\end{equation}
并以
\begin{equation}
\mathcal{Z}=\big[\mathcal{X}^\top\ \ \mathcal{U}^\top\big]^\top
\end{equation}
作为统一的优化变量集合。由此，\autoref{eq:dyn_mpc_recursion} 以线性等式约束形式保证各机器人在整个预测域上的状态—控制演化一致，\autoref{eq:box_constraints} 则给出每个时刻的速度与加速度限幅。

综上，UAV 端可通过前述微分平坦映射，将在 $\{\mathbf{p},\mathbf{v},\mathbf{u}\}$ 空间中优化得到的轨迹还原为可执行的姿态与推力，从而保证与真实动力学的一致可跟踪性；UGV 端则依托麦克纳姆轮全向底盘，在 $x\!-\!y$ 平面上的二阶积分模型下生成的轨迹均具有可实现性。该统一动力学模型在保留各平台关键运动特性的同时，为后续章节中将导航、避障、编队与可见性等多目标统一纳入同一 NMPC 优化框架提供了连贯的一致表述。

% 简单场景环境下的在线高效优化，低算力平台
\section{基于点云聚类的环境感知与引导路径生成} \label{sec:cluster_path}
为在低算力约束下获得可计算、可约束且可微的环境表达，本节首先对共享全局点云采用 DBSCAN 密度聚类以分离障碍物与噪声；随后对每个聚类簇求解最小包络椭圆，并按高度范围扩展为最小包络椭圆，以少量几何参数紧凑描述障碍物；并基于栅格占据图信息采用混合 A* 启发式搜索\cite{zhou2019ral}在当前栅格地图上搜索无碰撞的局部规划目标点。

\subsection{基于点云聚类的环境感知方法}
设在统一世界坐标系下得到的三维点云集合为
\begin{equation}
P\subset\mathbb{R}^3.
\end{equation}
为在含噪且几何形状不规则的场景中自动分离连通障碍物，本文采用 DBSCAN\cite{dbscan} 方法。给定邻域半径 $\delta>0$ 与最小邻域点数 $\mathrm{Pts}_{\min}\in\mathbb{N}^+$，对任一点 $p\in P$ 定义其 $\delta$ 邻域为
\begin{equation}
N_\delta(p)=\big\{\,q\in P\ \big|\ \mathrm{dist}(p,q)\le\delta\,\big\}.
\label{eq:dbscan_nei}
\end{equation}
当 $|N_\delta(p)|\ge \mathrm{Pts}_{\min}$ 时称 $p$ 为核心点。算法从未处理点出发，若该点为核心点，则将其邻域内的点加入候选集合，并在候选集合中继续查找新的核心点；通过“密度可达”关系（存在序列 $p_1,\dots,p_n\subset P$，满足 $p_1=p,\,p_n=q,\, p_{i+1}\in N_\delta(p_i)$）递归扩展，直至无法继续扩张为止，将所得点集记为一个聚类簇。处于多个核心点邻域交界处的边缘点随机归属到某一簇，未被归入任何簇的点标记为噪声。最终可得到若干簇
\begin{equation}
C_1,\dots,C_{N_c}\subset P,
\end{equation}
每个簇对应空间中由点密度连通的一块障碍结构。

\begin{figure}[htbp]
    \centering
    \includegraphics[width=0.8\linewidth]{chapter_3/ellipes_vis.png}
    \caption{环境聚类为最小包络椭圆的结果}
    \label{fig:mee_cluster}
\end{figure}

为了将聚类结果抽象为轻量且可微的几何表示，本文对每个簇的地平面投影
\begin{equation}
\Pi_{xy}(C_i)=\big\{[x\ y]^\top\ \big|\ (x,y,z)^\top\in C_i\big\}\subset\mathbb{R}^2
\label{eq:cluster_proj}
\end{equation}
求解其最小包络椭圆（minimum enclosing ellipse, MEE）。在平面上采用一般二次曲线进行参数化：
\begin{equation}
r\,x^2 + t\,xy + s\,y^2 + u\,x + v\,y + w = 0,
\label{eq:ellipse_quad}
\end{equation}
其中 $r,s$ 描述 $x$ 与 $y$ 方向的二次项，$t$ 刻画交叉耦合与旋转程度，$u,v$ 对应中心的线性平移项，$w$ 控制整体偏移。将\autoref{eq:ellipse_quad} 写为矩阵形式
\begin{equation}
x^\top \mathbf{A}\,x + 2\,\mathbf{b}^\top x + w = 0,\qquad 
\mathbf{A}=
\begin{bmatrix}
r & t/2\\[2pt]
t/2 & s
\end{bmatrix},\quad 
\mathbf{b}=\tfrac{1}{2}
\begin{bmatrix}
u\\[2pt]
v
\end{bmatrix},\quad 
x=
\begin{bmatrix}
x\\[2pt]
y
\end{bmatrix}.
\label{eq:ellipse_matrix}
\end{equation}
当 $\mathbf{A}\succ 0$ 时，该二次型对应闭合椭圆。椭圆中心由一阶驻点条件给出
\begin{equation}
c=-\mathbf{A}^{-1}\mathbf{b}.
\label{eq:ellipse_center}
\end{equation}
对 $\mathbf{A}$ 做特征分解
\begin{equation}
\mathbf{A}=\mathbf{R}\,\mathrm{diag}(\lambda_1,\lambda_2)\,\mathbf{R}^\top,\qquad 
\lambda_1,\lambda_2>0,\ \mathbf{R}^\top\mathbf{R}=\mathbf{I},
\label{eq:ellipse_eig}
\end{equation}
在主轴坐标下可将二次型化为标准式
\begin{equation}
\big(\mathbf{R}^\top(x-c)\big)^\top 
\mathrm{diag}\!\Big(\tfrac{1}{a^2},\tfrac{1}{b^2}\Big)
\big(\mathbf{R}^\top(x-c)\big)=1,
\label{eq:ellipse_standard}
\end{equation}
由此得到半轴长度 $a\ge b>0$ 以及椭圆朝向角 $\theta$（由 $\mathbf{R}$ 的主轴方向确定）。

在得到平面椭圆后，进一步利用簇点在高度方向的极值将其扩展为“最小包络椭圆柱”。记簇 $C_i$ 的 $z$ 轴最小值与最大值分别为
\begin{equation}
z_i^-=\min_{(x,y,z)\in C_i} z,\qquad 
z_i^+=\max_{(x,y,z)\in C_i} z,
\label{eq:cluster_z_extreme}
\end{equation}
则对应的椭圆柱可表示为
\begin{equation}
\mathcal{C}_i=
\Big\{(x,y,z)\in\mathbb{R}^3\ \Big|\ 
\phi_i\big([x\ y]^\top\big)\le 1,\ 
z\in[z_i^-,\,z_i^+] \Big\},
\label{eq:ellipse_cylinder}
\end{equation}
其中 $\phi_i(\cdot)$ 为\autoref{eq:ellipse_standard} 中的标准化二次型。簇 $C_i$ 的高度为 $h_i=z_i^+-z_i^-$。以集合
\begin{equation}
\{\mathcal{C}_i\}_{i=1}^{N_c}
\end{equation}
表示全部障碍，即得到由参数
\(
\{c_i,a_i,b_i,\theta_i,z_i^-,z_i^+\}
\)
构成的紧凑环境模型。该模型参数维度低、表达连续可微，便于在后续章节中直接构造距离函数、法向量与避障约束，为协同规划提供可靠的几何支撑。\autoref{fig:mee_cluster}展示了DBSCAN 算法在仿真中的点云数据上的聚类效果。算法以点云的空间分布特征为依据，成功识别出了场景中的多个长方体障碍物结构（如图中蓝色椭圆柱所示）。聚类结果表明，该方法能够有效地将空间中离散的点云数据划分为具有明确物理意义的障碍物区域，为后续章节提供了可靠的环境信息。

\subsection{基于混合 A* 的局部引导目标点搜索方法}
本文在二维栅格占据图上采用混合（Kino\mbox{-}Dynamic）A* 生成满足动力学与安全约束的离散可跟踪路径片段，并从中选取局部引导目标点。设当前工作栅格地图为
\[
\mathcal{M}\in\{0,1\}^{N_x\times N_y},
\]
其中 \(0\) 表示自由单元，\(1\) 表示被障碍占据；全局给定的规划起点和终点分别为 \(v_{\mathrm{start}},v_{\mathrm{goal}}\)，对应世界坐标为 \(p(v_{\mathrm{start}})\)、\(p(v_{\mathrm{goal}})\)。当 \(p(v_{\mathrm{goal}})\) 落在当前感知范围之外时，仅依靠局部 NMPC 很难直接找到可行路径，因此需要一个局部参考点来引导优化。

混合 A* 在状态空间中显式引入速度与加速度约束。结合\autoref{sec:dyn_model} 中的统一模型，对 UAV 在平面上的近似动力学取
\begin{equation}
\dot{\mathbf{x}}(t)=A\,\mathbf{x}(t)+B\,\mathbf{u}(t),\qquad
\mathbf{x}(t)=\begin{bmatrix}\mathbf{p}(t)\\ \mathbf{v}(t)\end{bmatrix}\in\mathbb{R}^4,\quad
\mathbf{u}(t)\in\mathbb{R}^2,
\label{eq:kino_ct}
\end{equation}
其中 \(\mathbf{p}=[x,y]^\top\) 为平面位置，\(\mathbf{v}=[v_x,v_y]^\top\) 为平面速度，\(\mathbf{u}=\ddot{\mathbf{p}}\) 为平面加速度，矩阵
\[
A=\begin{bmatrix}\mathbf{0}&\mathbf{I}\\ \mathbf{0}&\mathbf{0}\end{bmatrix},\qquad
B=\begin{bmatrix}\mathbf{0}\\ \mathbf{I}\end{bmatrix}.
\]
搜索图的节点不再仅表示一个栅格索引，而是“到达该栅格单元时的末状态”，记为 \(\mathbf{x}_c\)；边则由满足\autoref{eq:kino_ct} 的“动作基元”（motion primitive）构成。动作基元定义为在固定时长 \(\tau\) 内施加分量恒定的加速度输入
\begin{equation}
\mathbf{u}(t)\equiv \mathbf{u}_d\in U_D\subset[-u_{\max},u_{\max}]^2,\qquad t\in[0,\tau],
\label{eq:kino_primitive_u}
\end{equation}
由此可通过闭式积分得到末状态
\begin{equation}
\mathbf{x}(t+\tau)
= e^{A\tau}\mathbf{x}(t)+\int_0^{\tau} e^{A(\tau-\xi)}B\,\mathbf{u}_d\,\mathrm{d}\xi,
\label{eq:kino_primitive_x}
\end{equation}
并根据得到的末位置确定其落入的目标栅格单元。为在可达性与分支规模之间折中，将每个加速度分量区间 \([-u_{\max},u_{\max}]\) 按均匀步长量化为有限集合，从而每次扩展仅产生有限个候选动作基元。

在扩展过程中，对每个候选动作基元按以下规则筛选：其一，沿动作基元连续轨迹在栅格地图 \(\mathcal{M}\) 上进行碰撞检测，若轨迹穿越占据单元则认为该边不可行；其二，检查末状态是否满足 \(\|\mathbf{v}\|\le v_{\max}\)、\(\|\mathbf{u}_d\|\le u_{\max}\) 等动力学约束；其三，对落入同一栅格单元的多个候选，仅保留评价代价最小者以剪枝，从而有效控制搜索树的规模\cite{zhou2019ral}。

混合 A* 的代价函数同时度量控制能量与时间长度。对持续时长为 \(\tau\)、控制恒为 \(\mathbf{u}_d\) 的动作基元，其边代价取为
\begin{equation}
e_c = \big(\|\mathbf{u}_d\|_2^2+\rho\big)\,\tau,
\label{eq:kino_edge_cost}
\end{equation}
其中 \(\rho>0\) 为时间惩罚系数。节点的累计实际代价 \(g_c\) 为从起点到该节点的所有边代价之和。为加速收敛，引入基于最优控制解析解的启发式代价 \(h_c\)。忽略障碍与输入饱和，只考虑\autoref{eq:kino_ct} 的自由系统，从当前状态 \(\mathbf{x}_c\) 到目标状态 \(\mathbf{x}_f\) 的“最小能量 + 时间”问题可写为
\begin{equation}
\min_{\mathbf{u}(\cdot),T>0}\ \int_0^T \|\mathbf{u}(t)\|_2^2\,\mathrm{d}t+\rho\,T,
\label{eq:kino_ocp}
\end{equation}
其最优控制在每个轴上均为时间的仿射函数；利用庞特里亚金极小值原理可得到解析形式的最优代价 \(J^*(\mathbf{x}_c,\mathbf{x}_f)\)。将该最优代价作为启发式
\begin{equation}
h_c = J^*(\mathbf{x}_c,\mathbf{x}_f),
\label{eq:kino_heur}
\end{equation}
即可保证启发函数可采纳且具有较强引导性。混合 A* 的评价函数定义为
\begin{equation}
f_c = g_c + h_c.
\label{eq:kino_f}
\end{equation}

算法同样维护开放表 \(\mathrm{OPEN}\) 与关闭表 \(\mathrm{CLOSED}\)。初始化时根据当前 UAV 状态设置 \(g_c(\mathbf{x}_{\mathrm{start}})=0\)，将起始节点放入 \(\mathrm{OPEN}\)。每次迭代从 \(\mathrm{OPEN}\) 中取出 \(f_c\) 最小的节点并展开其所有可行动作基元，更新或插入新的候选节点。当某一节点的末位置落入目标栅格邻域、且其到目标的解析最优连接满足碰撞与约束检验时，认为到达目标并终止搜索。

本文并不直接使用整条离散路径，而是根据当前感知范围和规划步长，从搜索得到的路径中选取一个局部引导目标点。具体地，若全局目标点 \(v_{\mathrm{goal}}\) 落在当前感知栅格内，则沿最优路径从起点向前行走一固定弧长或时间窗，取对应节点的末位置作为局部目标；若全局目标在感知范围之外，则从搜索树上选取距离 UAV 最近、且在自由空间中的末状态作为局部目标。该局部目标点随后作为 NMPC 中的位置参考约束，保证在传感半径有限、全局目标遥远的情况下，规划始终沿着全局可行的方向推进，而不会因为缺乏合理参考而出现发散或停滞。

在实际系统中，混合 A* 生成的离散状态序列会在时间轴上重新采样并作为无人机的参考轨迹，其效果如 \autoref{fig:kino_astar_demo} 所示。图中彩色方块表示由体素占据图提取的障碍区域，青色圆点连成的曲线为混合 A* 搜索得到的动力学可行路径，起点位于左端的绿色坐标轴处，终点为右端的红色球标记位置。可以看到，路径在保证与障碍保持安全间距的同时，自然地满足最大速度与最大加速度约束，不出现尖锐折线或不可跟踪的急转。

\begin{figure}[htbp]
    \centering
    \includegraphics[width=0.7\linewidth]{chapter_3/kino_astar.png} 
    \caption{混合 A* 生成的动力学可行引导路径示意图}
    \label{fig:kino_astar_demo}
\end{figure}

\section{基于离散控制障碍函数的避障约束构建} \label{sec:dcbf}
上一节的 混合 A* 在体素占据图上搜索得到满足动力学约束的无碰撞路径。本文并不直接将整条离散路径约束到优化变量上，而是仅从搜索结果中选取当前感知范围内距离全局目标最近的一个路径点，作为滚动时域内的局部引导目标。机器人之间以及机器人与障碍之间的实际碰撞避免均由本节给出的离散控制障碍函数（Discrete Control Barrier Function, DCBF）\cite{dcbf}实现。具体做法是：在统一的离散动力学模型下，将上一节构造的最小包络椭圆柱环境表示转化为关于状态的可微距离函数，并通过 DCBF 一步条件构造不等式约束，从而在 NMPC 中显式保证预测域内的轨迹安全性。

\subsection{经典离散控制障碍函数}

考虑与前文一致的离散动力学系统
\begin{equation}
\mathbf{x}_{k+1}=f(\mathbf{x}_k,\mathbf{u}_k),\qquad k=0,1,\dots
\label{eq:dcbf_dyn_general}
\end{equation}
其中 \(\mathbf{x}_k\in\mathbb{R}^n\) 为系统状态，\(\mathbf{u}_k\) 为控制输入。设障碍集合为 \(\mathcal{O}\subset\mathbb{R}^n\)，记状态到障碍的最小距离为 \(d_{\min}(\mathbf{x},\mathcal{O})\)，在需要时可以扣除一个安全裕度 \(\rho\ge 0\)，得到控制障碍函数
\begin{equation}
h(\mathbf{x}) = d_{\min}(\mathbf{x},\mathcal{O})-\rho.
\label{eq:dcbf_h_def}
\end{equation}
由此定义安全集
\begin{equation}
\mathcal{C}=\big\{\mathbf{x}\in\mathbb{R}^n \,\big|\, h(\mathbf{x})\ge 0\big\}.
\label{eq:dcbf_safe_set}
\end{equation}

若初值满足 \(\mathbf{x}_0\in\mathcal{C}\)，希望在后续所有采样时刻都保持 \(\mathbf{x}_k\in\mathcal{C}\)。离散控制障碍函数的核心思想是通过对 \(h(\mathbf{x}_k)\) 的变化率施加约束，使安全集在系统演化下保持前向不变，即存在类-\(\mathcal{K}\) 函数 \(\alpha(\cdot)\) 使得
\begin{equation}
h(\mathbf{x}_{k+1})-h(\mathbf{x}_k)\ \ge\ -\,\alpha\big(h(\mathbf{x}_k)\big),\qquad k=0,1,\dots
\label{eq:dcbf_core}
\end{equation}
在本文中采用最常用的线性形式 \(\alpha(h)=\gamma h\)，其中 \(\gamma\in(0,1)\)，则有
\begin{equation}
h(\mathbf{x}_{k+1})\ \ge\ (1-\gamma)\,h(\mathbf{x}_k).
\label{eq:dcbf_linear_alpha}
\end{equation}
该式刻画了“向安全边界的最大允许衰减率”：当系统远离障碍时，\(h(\mathbf{x}_k)\) 较大，允许的下降幅度也较大；当系统逐渐接近障碍边界时，\(h(\mathbf{x}_k)\) 变小，约束自动收紧，从而避免越界进入障碍集合 \(\mathcal{O}\)。

为了在优化框架中获得关于控制输入的可计算形式，对 \(h(\mathbf{x}_{k+1})\) 在 \(\mathbf{x}_k\) 处进行一阶近似：
\begin{equation}
h(\mathbf{x}_{k+1})\approx h(\mathbf{x}_k)+\nabla h(\mathbf{x}_k)^\top\!\big(\mathbf{x}_{k+1}-\mathbf{x}_k\big).
\label{eq:dcbf_first_order}
\end{equation}
将 \autoref{eq:dcbf_dyn_general} 代入，可以得到关于 \(\mathbf{u}_k\) 的仿射不等式约束
\begin{equation}
h(\mathbf{x}_k)+\nabla h(\mathbf{x}_k)^\top\!\big(f(\mathbf{x}_k,\mathbf{u}_k)-\mathbf{x}_k\big)\ \ge\ (1-\gamma)\,h(\mathbf{x}_k).
\label{eq:dcbf_affine}
\end{equation}
当动力学为线性或仿射形式（例如前文统一二阶积分模型的离散形式）时，\autoref{eq:dcbf_affine} 对控制量是线性的；在一般非线性 \(f\) 的情况下，也可以在 NMPC 中采用逐次线性化进行处理。对于多障碍场景，本文为每一个障碍 \(i\) 构造对应的 \(h_i(\mathbf{x})\)，在优化中同时施加所有约束。综上所述，在预测域 \(H\) 内对 \(k=0,\dots,H-1\) 施加 \autoref{eq:dcbf_linear_alpha}，即可在考虑系统动力学与输入限幅的同时，将“始终处于安全集内”的要求系统性地嵌入滚动时域优化过程。

\subsection{基于最小包络椭圆柱的几何避障约束}

此前已经通过最小包络椭圆构建了障碍物的椭圆柱表示。基于该表示，本节给出从机器人位置到椭圆边界的最小距离计算方法，并将其嵌入 DCBF 框架，形成可微的几何避障约束。对于第 \(i\) 个障碍，记其最小包络椭圆柱参数为
\begin{equation}
\mathrm{obs}^i=\big[x^{i}_{\mathrm{ob}},\,y^{i}_{\mathrm{ob}},\,z^{i}_{\mathrm{ob}},\,a_i,\,b_i,\,\theta_i,\,z_i^{-},\,z_i^{+}\big],
\label{eq:obs_param}
\end{equation}
其中 \((x^{i}_{\mathrm{ob}},y^{i}_{\mathrm{ob}},z^{i}_{\mathrm{ob}})\) 为柱体中心，\(a_i\ge b_i>0\) 为长短半轴，\(\theta_i\) 为长轴相对世界坐标系 \(x\) 轴的偏转角，\([z_i^{-},z_i^{+}]\) 为高度范围。令时刻 \(k\) 机器人位置为
\(\mathbf{p}_k=[x_k,y_k,z_k]^\top\)，则平面位置相对障碍中心的差为
\begin{equation}
\mathbf{d}_i(k)=
\begin{bmatrix}
x_k-x^{i}_{\mathrm{ob}}\\
y_k-y^{i}_{\mathrm{ob}}
\end{bmatrix},\qquad
r_i(k)=\big\|\mathbf{d}_i(k)\big\|_2.
\label{eq:d_i_k}
\end{equation}
平面连线与世界系 \(x\) 轴的夹角为
\begin{equation}
\phi_i(k)=\mathrm{atan2}\!\big(y_k-y^{i}_{\mathrm{ob}},\,x_k-x^{i}_{\mathrm{ob}}\big),
\label{eq:phi_def}
\end{equation}
与椭圆长轴的相对角度为
\begin{equation}
\vartheta_i(k)=\phi_i(k)-\theta_i ,
\label{eq:theta_rel_def}
\end{equation}
其中 \(\vartheta_i(k)\) 规范到 \((-\pi,\pi]\)。

沿方向 \(\vartheta_i(k)\) 从椭圆中心到边界的径向距离可以写成封闭式
\begin{equation}
\ell_i\big(\vartheta_i(k)\big)
=\frac{a_i b_i}{\sqrt{\,b_i^{2}\cos^{2}\!\vartheta_i(k)+a_i^{2}\sin^{2}\!\vartheta_i(k)\,}},
\label{eq:ellipse_radius_cos_form}
\end{equation}
等价的形式为
\begin{equation}
\ell_i\big(\vartheta_i(k)\big)
=\sqrt{\frac{a_i^{2}b_i^{2}\big(1+\tan^{2}\!\vartheta_i(k)\big)}{\,b_i^{2}+a_i^{2}\tan^{2}\!\vartheta_i(k)\,}}.
\label{eq:ellipse_radius_tan_form}
\end{equation}
若将机器人在平面上近似为半径为 \(d_{\mathrm{safe}}>0\) 的圆，则机器人到第 \(i\) 个椭圆柱侧壁的平面距离为
\begin{equation}
d_{\mathrm{clr},i}(k)=r_i(k)-\ell_i\big(\vartheta_i(k)\big)-d_{\mathrm{safe}}.
\label{eq:clearance_def}
\end{equation}
仅当机器人高度满足
\begin{equation}
z_k\in\big[z_i^{-},\,z_i^{+}\big]
\label{eq:height_active}
\end{equation}
时将该障碍视为“激活”，否则认为机器人在该障碍的上下方，不施加相关约束。由此可以定义针对第 \(i\) 个椭圆柱的离散控制障碍函数
\begin{equation}
h_i(\mathbf{x}_k)=d_{\mathrm{clr},i}(k),
\label{eq:cbf_hi}
\end{equation}
其安全集为 \(\{\mathbf{x}_k\mid h_i(\mathbf{x}_k)\ge 0\}\)。在预测域内，对所有被激活的障碍施加一步 DCBF 约束
\begin{equation}
h_i(\mathbf{x}_{k+1})\ \ge\ (1-\gamma)\,h_i(\mathbf{x}_k),
\qquad \gamma\in(0,1),\quad k=0,\ldots,H-1,
\label{eq:dcbf_constraint_ec}
\end{equation}
即可在保证统一动力学模型 \(\mathbf{x}_{k+1}=f(\mathbf{x}_k,\mathbf{u}_k)\) 的同时，实现对椭圆柱障碍的碰撞避免。

通过上述构造，最小包络椭圆柱参数 \(\{a_i,b_i,\theta_i,z_i^{-},z_i^{+}\}\) 与机器人位置之间的关系以封闭形式给出了到障碍边界的最小距离。将这些几何距离函数代入 DCBF 约束 \autoref{eq:dcbf_constraint_ec}，并配合上一节混合 A* 选取的局部引导目标点，NMPC 在整个预测域内可以同时考虑动力学可行性和空间安全性，从而在未知且复杂的环境中实现实时、可靠的协同避障规划。

\section{多目标约束的实时协同轨迹生成} \label{sec:traj_opt}
基于前述统一离散动力学模型与椭圆柱环境表达，本节构建以 UAV 为 leader、UGV 为 follower 的协同优化框架。UAV 根据 混合 A* 给出的全局无碰撞路径，在每个规划周期选取局部引导路点并作为轨迹优化的主要参考；各 UGV 通过与 UAV 及彼此之间的相对位置约束实现从动协同，形成具有一定几何结构的编队。优化目标综合考虑引导路点跟踪、编队几何保持、空地间视线连通质量以及控制能耗等多项子目标，约束条件则由统一动力学方程、速度与加速度限幅、基于 DCBF 的椭圆柱避障约束以及机器人间安全距离约束共同构成。整个问题在滚动时域框架下以 NMPC 形式在线求解，使系统能够在环境持续变化和任务目标更新时实现实时重规划与稳定协同。

\subsection{无人机引导路点约束构建}

上一节的 混合 A* 在体素占据图中给出了从当前 UAV 位置到全局目标的一条满足动力学与碰撞约束的离散路径。直接在优化中强制整个路径通常会导致约束规模过大。为此，本文在每个规划周期只从该离散路径中提取一个位于当前感知范围边界附近的“局部引导路点”，并在预测域内将 UAV 逐步吸引到该路点附近；随着滚动时域向前推进，引导路点也不断沿 混合 A* 路径向全局目标移动，从而实现全局导航。

设在当前规划时刻由 混合 A* 输出的三维路径折线为
\(\{\mathbf{w}_j\}_{j=0}^{L}\subset\mathbb{R}^3\)，其中 \(\mathbf{w}_0\) 与 UAV 当前姿态在位置上对齐。根据统一动力学模型，规划层对 UAV 采用固定采样周期 \(\Delta t\) 和预测步数 \(H\)。记 UAV 在平面上的最大速度为 \(\|\mathbf{v}_{\max}^{\mathrm{uav}}\|_2\)，选取一个折减系数 \(\eta\in(0,1]\)，得到预测域内可达到的保守平均速度
\begin{equation}
\bar v=\eta\,\big\|\mathbf{v}_{\max}^{\mathrm{uav}}\big\|_2.
\label{eq:wp_vbar}
\end{equation}
据此估计在单次预测域内理论上能够前进的最大弧长
\begin{equation}
L_{\mathrm{hor}}=\bar v\,H\Delta t .
\label{eq:wp_Lhor}
\end{equation}

对离散路径按弧长进行累加：令 \(s_0=0\)，并递推
\(s_j=s_{j-1}+\|\mathbf{w}_j-\mathbf{w}_{j-1}\|_2\)。给定预测域可推进弧长上界 \(L_{\mathrm{hor}}\)，取
\begin{equation}
s^\star=\min\{L_{\mathrm{hor}},\,s_L\},
\label{eq:wp_sstar}
\end{equation}
在路径上寻找第一个满足 \(s_J\ge s^\star\) 的索引 \(J\)。当 \(s_{J-1}<s^\star<s_J\) 时，在段
\([\mathbf{w}_{J-1},\mathbf{w}_J]\) 上按弧长进行线性插值，
\begin{equation}
\lambda=\frac{s^\star-s_{J-1}}{s_J-s_{J-1}},\qquad
\mathbf{w}^\star=(1-\lambda)\,\mathbf{w}_{J-1}+\lambda\,\mathbf{w}_J,
\label{eq:wp_interp}
\end{equation}
从而得到位于当前感知与控制能力边界附近的局部引导路点 \(\mathbf{w}^\star\in\mathbb{R}^3\)。当全局路径较短时，\(s^\star=s_L\)，此时 \(\mathbf{w}^\star\) 退化为终点。

在统一动力学模型下，UAV 的未来位置变量记为
\(\mathbf{p}^{a}_k\in\mathbb{R}^3\)（\(k=0,\dots,H\)），其中
\(\mathbf{p}^{a}_0\) 由当前状态给定。为了避免在预测的前期阶段对局部目标路点约束过紧，从而保留一定的机动空间，本文采用随时间单调增大的权重序列。具体地，引入参数
\(\lambda^{\min}_{\mathrm{wp}}>0\)、\(\lambda^{\max}_{\mathrm{wp}}>\lambda^{\min}_{\mathrm{wp}}\) 与指数 \(\nu\ge 1\)，定义
\begin{equation}
\lambda_k=\lambda^{\min}_{\mathrm{wp}}
+\big(\lambda^{\max}_{\mathrm{wp}}-\lambda^{\min}_{\mathrm{wp}}\big)\!
\left(\frac{k}{H}\right)^{\nu},
\qquad k=0,\dots,H-1,
\label{eq:wp_lambda}
\end{equation}
并在预测末端给出更大的终端权重
\(\lambda_{\mathrm{term}}=\kappa\,\lambda^{\max}_{\mathrm{wp}}\)（\(\kappa\gg 1\)）。这样，在靠近当前时刻的位置点上，UAV 可以在不违背其他约束的前提下自由调整轨迹；而在预测域后半段以及终端时刻，优化会越来越强调与引导路点 \(\mathbf{w}^\star\) 的一致性。

基于上述权重设计，引导路点跟踪代价写为
\begin{equation}
J_{\mathrm{wp}}=
\sum_{k=0}^{H-1}\lambda_k\,\big\|\mathbf{p}^{a}_k-\mathbf{w}^\star\big\|_2^2
+\lambda_{\mathrm{term}}\,\big\|\mathbf{p}^{a}_H-\mathbf{w}^\star\big\|_2^2 .
\label{eq:wp_cost}
\end{equation}
该代价项使得在每个规划周期内，UAV 既沿 混合 A* 给出的全局无碰撞通道推进，又能够在局部范围内根据最新地图信息进行微调。

\subsection{空地系统间可见性约束构建}

\begin{figure}[htbp]
    \centering
    \includegraphics[width=0.65\linewidth]{chapter_3/可见性约束.png} 
    \caption{空地系统间可见性约束示意图}
    \label{fig:visable_constraint}
\end{figure}

为在优化中显式度量并提升空地平台之间的视距质量，如\autoref{fig:visable_constraint}所示，本节设计空地系统间可见性约束，记时刻 \(k\) 的 UAV 位置为
\[
\mathbf{p}^{a}_k=[x^a_k,y^a_k,z^a_k]^\top,
\]
第 \(m\) 辆 UGV 的位置为
\[
\mathbf{p}^{g_m}_k=[x^{g_m}_k,y^{g_m}_k,z^{g_m}_k]^\top.
\]
二者在当前时刻的连线段写作
\[
\mathcal{L}^{(m)}_k(t)=\mathbf{p}^{a}_k+t\,\mathbf{d}^{(m)}_k,\qquad 
\mathbf{d}^{(m)}_k=\mathbf{p}^{g_m}_k-\mathbf{p}^{a}_k,\ \ t\in[0,1]. 
\]

对任意一个最小包络椭圆柱
\[
\mathrm{obs}^i=[x^{i}_{\mathrm{ob}},y^{i}_{\mathrm{ob}},z^{i}_{\mathrm{ob}},a_i,b_i,\theta_i,z_i^{-},z_i^{+}]
\]
首先计算其中心点
\(
\mathbf{c}_i=[x^{i}_{\mathrm{ob}},y^{i}_{\mathrm{ob}},z^{i}_{\mathrm{ob}}]^\top
\)
到连线段 \(\mathcal{L}^{(m)}_k\) 的最近点。线段参数取为
\[
t_i^\star=\Pi_{[0,1]}\!\left(\frac{(\mathbf{c}_i-\mathbf{p}^{a}_k)^\top \mathbf{d}^{(m)}_k}{\|\mathbf{d}^{(m)}_k\|_2^2}\right),
\]
最近点为
\[
\mathbf{q}_{i,m}(k)=\mathcal{L}^{(m)}_k(t_i^\star),
\]
其中 \(\Pi_{[0,1]}(\cdot)\) 表示对区间 \([0,1]\) 的截断算子。只有当最近点在高度方向与柱体有重叠时，该障碍才会对视线产生潜在遮挡作用，因此定义高度激活指示量
\[
\chi_{i,m}(k)=
\begin{cases}
1, & q_{i,m}^z(k)\in[z_i^{-},\,z_i^{+}],\\[2pt]
0, & \text{otherwise},
\end{cases}
\]
其中 \(q_{i,m}^z(k)\) 为 \(\mathbf{q}_{i,m}(k)\) 的 \(z\) 分量。

在被激活的情形下，将 \(\mathbf{q}_{i,m}(k)\) 投影到地平面，得到从柱体中心指向最近点的平面向量
\[
\mathbf{v}_{i,m}(k)=
\begin{bmatrix}
q_{i,m}^x(k)-x^{i}_{\mathrm{ob}}\\[2pt]
q_{i,m}^y(k)-y^{i}_{\mathrm{ob}}
\end{bmatrix},\qquad
r_{i,m}(k)=\big\|\mathbf{v}_{i,m}(k)\big\|_2.
\]
记该向量相对于世界系 \(x\) 轴的方位角为 \(\phi_{i,m}(k)=\mathrm{atan2}(v_{i,m,y}(k),v_{i,m,x}(k))\)，与椭圆长轴的相对角为
\[
\vartheta_{i,m}(k)=\phi_{i,m}(k)-\theta_i,
\]
并采用与\autoref{eq:ellipse_radius_cos_form} 相同的表达式计算沿方向 \(\vartheta_{i,m}(k)\) 到椭圆边界的径向距离
\[
\ell_i\big(\vartheta_{i,m}(k)\big)
=\frac{a_i b_i}{\sqrt{\,b_i^{2}\cos^{2}\!\vartheta_{i,m}(k)+a_i^{2}\sin^{2}\!\vartheta_{i,m}(k)\,}}.
\]
由此得到连线对第 \(i\) 个椭圆柱侧壁的平面最小距离
\[
d^{\mathrm{LoS}}_{i,m}(k)=r_{i,m}(k)-\ell_i\big(\vartheta_{i,m}(k)\big).
\]

设 \(d_{\mathrm{safe}}>0\) 为视线安全裕度。当 \(d^{\mathrm{LoS}}_{i,m}(k)\) 小于该阈值时，认为 UAV 与第 \(m\) 辆 UGV 的连线被第 \(i\) 个障碍“挤压”，存在潜在遮挡风险，需要通过优化调整双方位置。对单个时间步 \(k\) 定义可见性代价项
\[
J_{\mathrm{vis}}(k)=\lambda_{\mathrm{vis}}\sum_{m=1}^{M}\sum_{i=1}^{N_c}
\chi_{i,m}(k)\,\big[\,d_{\mathrm{safe}}-d^{\mathrm{LoS}}_{i,m}(k)\,\big]_+^{\,2},\qquad [x]_+=\max(x,0),
\]
其中 \(\lambda_{\mathrm{vis}}\) 为可见性损失的权重。\([\,d_{\mathrm{safe}}-d^{\mathrm{LoS}}_{i,m}(k)\,]_+^{2}\) 仅在最小距离小于安全裕度时产生平方惩罚，使得当 LoS 足够“宽松”时不额外约束机器人的机动空间，而在连线逐渐逼近障碍时自动放大梯度，驱动 UAV 与 UGV 远离连线遮挡区域。

将可见性代价在整个预测域和所有机器人上进行求和，得到最终的可见性目标函数
\begin{equation}
J_{\mathrm{vis}}
=\sum_{k=0}^{H-1}\lambda_{\mathrm{vis}}\sum_{m=1}^{M}\sum_{i=1}^{N_c}
\chi_{i,m}(k)\,\big[\,d_{\mathrm{safe}}-d^{\mathrm{LoS}}_{i,m}(k)\,\big]_+^{\,2}.
\label{eq:J_vis_total}
\end{equation}
该项与引导路点代价 \(J_{\mathrm{wp}}\) 以及后续的编队与能耗代价一同构成整体优化目标，使 \(\mathbf{p}^{a}_k\) 与 \(\{\mathbf{p}^{g_m}_k\}_{m=1}^{M}\) 在复杂障碍环境中既能完成任务推进，又能尽量维持稳定的空地视距联系。

\subsection{空地队形约束构建}
为在跟随 UAV 推进的同时保持地面编队的几何稳定性，本文从尺度、形状和空地距离三个方面对空地系统的队形进行建模，如\autoref{fig:formation_constraint}所示。尺度约束限制任意两车之间的距离落在给定的可行带内，从而避免过密导致的碰撞风险以及过疏造成的联络退化。形状约束使由四辆车构成的四边形接近矩形结构，防止队形在几何上发生共线或极度扁长的情况。空地距离约束则直接约束 UAV 与各 UGV 的相对距离，避免 UAV 长时间脱离地面编队包络，保证 UWB 与视距观测几何的可观测性。

\begin{figure}[htbp]
    \centering
    \includegraphics[width=0.6\linewidth]{chapter_3/队形约束.png} 
    \caption{空地系统队形约束示意图}
    \label{fig:formation_constraint}
\end{figure}

在预测时刻 \(k\)，记四辆 UGV 的三维位置为
\[
\mathbf{p}^{g_i}_k=
\begin{bmatrix}
x^{g_i}_k & y^{g_i}_k & z^{g_i}_k
\end{bmatrix}^\top\in\mathbb{R}^3,\quad i\in\{0,1,2,3\},
\]
UAV 的位置为 \(\mathbf{p}^a_k=[x^a_k,y^a_k,z^a_k]^\top\)。本节在每个时间步上构造编队代价
\begin{equation}
\ell_{\mathrm{form}}(k)=\ell_{\mathrm{dis}}(k)+\ell_{\mathrm{shape}}(k)+\ell_{\mathrm{bound}}(k),
\label{eq:l_form_k}
\end{equation}
其中三项分别对应尺度、形状和空地距离约束。

首先给出尺度项 \(\ell_{\mathrm{dis}}(k)\)。设允许的车间距离带为 \([d_{\mathrm{low}},\,d_{\mathrm{high}}]\)。当两车距离处于该区间内部时不产生惩罚；当距离小于下界时按与下界的差值平方惩罚；当距离大于上界时按与上界的差值平方惩罚。定义任意两车之间的距离
\[
r_{ij}(k)=\big\|\mathbf{p}^{g_i}_k-\mathbf{p}^{g_j}_k\big\|_2,\quad i<j,
\]
并引入分段函数
\[
\psi(r)=
\begin{cases}
(d_{\mathrm{low}}-r)^2, & r<d_{\mathrm{low}},\\[2pt]
0, & d_{\mathrm{low}}\le r\le d_{\mathrm{high}},\\[2pt]
(r-d_{\mathrm{high}})^2, & r>d_{\mathrm{high}}.
\end{cases}
\]
则尺度代价写为
\begin{equation}
\ell_{\mathrm{dis}}(k)=\sum_{i<j}\psi\!\big(r_{ij}(k)\big).
\label{eq:l_dis_k}
\end{equation}
这一项在允许距离带内部保持队形自由度，在越界时提供平滑的二次回带驱动力，既保证安全间距，也避免编队被过度拉伸。

然后给出形状项 \(\ell_{\mathrm{shape}}(k)\)。目标是使四边形的四个内角接近直角，从而抑制共线与扁长退化。将四辆 UGV 在地平面上的投影按顺时针排序为顶点索引 \((v_0,v_1,v_2,v_3)\)。对顶点 \(v_i\) 及其相邻顶点 \(v_{i^-},v_{i^+}\)（下标按模 4 取值）定义三条边长
\[
a_i(k)=\big\|\mathbf{p}^{g_{v_i}}_k-\mathbf{p}^{g_{v_{i^-}}}_k\big\|_2,\quad
b_i(k)=\big\|\mathbf{p}^{g_{v_{i^+}}}_k-\mathbf{p}^{g_{v_i}}_k\big\|_2,\quad
c_i(k)=\big\|\mathbf{p}^{g_{v_{i^+}}}_k-\mathbf{p}^{g_{v_{i^-}}}_k\big\|_2.
\]
根据余弦定理，顶点 \(v_i\) 对应的内角 \(\delta_i(k)\) 的余弦为
\begin{equation}
\cos\big(\delta_i(k)\big)
=\frac{a_i(k)^2+b_i(k)^2-c_i(k)^2}{2\,a_i(k)\,b_i(k)},\qquad i=0,1,2,3.
\label{eq:cos_delta_k}
\end{equation}
当 \(\delta_i(k)=\tfrac{\pi}{2}\) 时有 \(\cos(\delta_i(k))=0\)。据此构造形状代价
\begin{equation}
\ell_{\mathrm{shape}}(k)=\sum_{i=0}^{3}\big(\cos(\delta_i(k))\big)^2.
\label{eq:l_shape_k}
\end{equation}
当四车趋于共线或形成极度扁长的形状时，某些顶点的 \(|\cos(\delta_i(k))|\) 会显著增大，从而推高 \(\ell_{\mathrm{shape}}(k)\)。该项在优化中推动四边形向近似直角结构回归，与尺度项共同保持编队既不过密也不过疏。

在此基础上，引入空地距离项 \(\ell_{\mathrm{bound}}(k)\) 来约束 UAV 与各 UGV 之间的相对距离。设希望 UAV 与每辆 UGV 的距离不超过上界 \(d_{\mathrm{bound}}>0\)。当 UAV 距离某一 UGV 小于等于该上界时不产生惩罚；当距离超出上界时按照超出量的平方给予代价。定义
\[
r^{ag}_i(k)=\big\|\mathbf{p}^a_k-\mathbf{p}^{g_i}_k\big\|_2,\quad i\in\{0,1,2,3\},
\]
以及分段函数
\[
\varphi(r)=
\begin{cases}
\big(r-d_{\mathrm{bound}}\big)^2, & r>d_{\mathrm{bound}},\\[2pt]
0, & r\le d_{\mathrm{bound}}.
\end{cases}
\]
则空地距离代价为
\begin{equation}
\ell_{\mathrm{bound}}(k)=\sum_{i=0}^{3}\varphi\!\big(r^{ag}_i(k)\big).
\label{eq:l_bound_k}
\end{equation}
该项鼓励 UAV 始终位于由 UGV 编队构成的有限空间邻域内。当 UAV 试图脱离队形包络时，对应的 \(r^{ag}_i(k)\) 会超过 \(d_{\mathrm{bound}}\)，从而触发二次惩罚，使优化结果将 UAV 拉回到与地面编队相对接近的区域，有利于维持良好的 UWB 和视距观测几何。

最后，将三项按权重并入队形目标。设 \(\lambda_{\mathrm{dis}}>0\)、\(\lambda_{\mathrm{shape}}>0\)、\(\lambda_{\mathrm{bound}}>0\) 分别为尺度、形状和空地距离项的权重，则在整个预测域上的编队总代价写为
\begin{equation}
J_{\mathrm{form}}
=\sum_{k=0}^{H-1}\Big(
\lambda_{\mathrm{dis}}\,\ell_{\mathrm{dis}}(k)
+\lambda_{\mathrm{shape}}\,\ell_{\mathrm{shape}}(k)
+\lambda_{\mathrm{bound}}\,\ell_{\mathrm{bound}}(k)
\Big).
\label{eq:J_form_final}
\end{equation}
在合适选取权重的情形下，地面编队可以在 UAV 领航推进过程中保持稳定的矩形拓扑和合理尺度，同时 UAV 与各 UGV 之间的距离被限制在给定范围内，为协同任务和相对定位提供稳健的几何结构支持。

\subsection{非线性优化问题建模}
本小节将前述Follower–Leader体系中的多目标需求统一为一个带硬约束与软代价的非线性优化问题。在每个规划周期内，设预测域长度为 \(H\)，采样间隔为 \(\Delta t\)。记机器人集合为
\[
\mathcal{R}=\{a,g_1,\dots,g_M\},
\]
其中 \(a\) 表示 UAV，\(g_m\) 表示第 \(m\) 辆 UGV。对任意 \(r\in\mathcal{R}\)，离散状态与控制分别记为
\[
\mathbf{x}^{(r)}_k=
\begin{bmatrix}
\mathbf{p}^{(r)}_k\\[2pt]
\mathbf{v}^{(r)}_k
\end{bmatrix}\in\mathbb{R}^{2d_r},\qquad
\mathbf{u}^{(r)}_k\in\mathbb{R}^{d_r},\qquad k=0,\dots,H-1,
\]
其中 \(\mathbf{p}^{(r)}_k\in\mathbb{R}^{d_r}\) 为位置，\(\mathbf{v}^{(r)}_k\in\mathbb{R}^{d_r}\) 为速度，UAV 取 \(d_a=3\)，UGV 取 \(d_{g_m}=2\) 表示平面运动。动力学采用统一的二阶积分离散模型
\begin{equation}
\mathbf{x}^{(r)}_{k+1}=\mathbf{A}_r\,\mathbf{x}^{(r)}_{k}+\mathbf{B}_r\,\mathbf{u}^{(r)}_{k},
\label{eq:nlp_dyn}
\end{equation}
其中 \(\mathbf{A}_r,\mathbf{B}_r\) 由\autoref{eq:disc_AB} 给出，仅与 \(d_r\) 与 \(\Delta t\) 有关。

将全体机器人变量按时间与个体堆叠为
\[
\mathcal{X}=\big[\mathbf{x}^{(a)}_1,\dots,\mathbf{x}^{(a)}_H,\;
\mathbf{x}^{(g_1)}_1,\dots,\mathbf{x}^{(g_M)}_H\big],
\]
\[
\mathcal{U}=\big[\mathbf{u}^{(a)}_0,\dots,\mathbf{u}^{(a)}_{H-1},\;
\mathbf{u}^{(g_1)}_0,\dots,\mathbf{u}^{(g_M)}_{H-1}\big],
\]
并以 \(\mathcal{Z}=[\mathcal{X}^\top\ \mathcal{U}^\top]^\top\) 作为优化变量。

目标函数由四部分组成。第一项为 UAV 的引导路点跟踪代价 \(J_{\mathrm{wp}}\)，其在\autoref{sec:traj_opt} 中已给出，对应 UAV 预测轨迹 \(\{\mathbf{p}^a_k\}_{k=0}^{H}\) 向引导点 \(w^\star\) 的渐进收敛。第二项为空地可见性软约束 \(J_{\mathrm{vis}}\)，在每个时间步上根据 UAV–UGV 连线与各椭圆柱之间的最近净空构造铰链平方损失，用于惩罚潜在遮挡并促使无人机与地面编队保持良好视距。第三项为 UGV 的编队代价 \(J_{\mathrm{form}}\)，包括车间距离带项、四边形直角化项以及 UAV–UGV 距离约束，用于保持编队尺度与形状的几何稳定性。第四项为控制正则项，用于抑制能耗并提升数值平滑性：
\begin{equation}
J_{\mathrm{ctrl}}
=\lambda_{\mathrm{ctrl}}\sum_{r\in\mathcal{R}}\sum_{k=0}^{H-1}\big\|\mathbf{u}^{(r)}_k\big\|_2^{2}
+\lambda_{\Delta u}\sum_{r\in\mathcal{R}}\sum_{k=0}^{H-2}
\big\|\mathbf{u}^{(r)}_{k+1}-\mathbf{u}^{(r)}_{k}\big\|_2^{2},
\label{eq:J_ctrl}
\end{equation}
其中\(\lambda_{\mathrm{ctrl}}>0\) 为控制信号的平滑权重，\(\lambda_{\Delta u}>0\) 为相邻控制增量的平滑权重。

综合上述各项，单周期内的非线性优化问题写为
\begin{equation}
\min_{\mathcal{X},\mathcal{U}}\ \ 
J_{\mathrm{wp}}
+J_{\mathrm{vis}}
+J_{\mathrm{form}}
+J_{\mathrm{ctrl}}
\label{eq:nlp_obj}
\end{equation}

约束由离散动力学等式、状态与控制约束、初值约束以及基于 DCBF 的安全不等式组成。对任意 \(r\in\mathcal{R}\) 与 \(k=0,\dots,H-1\) 有
\begin{equation}
\begin{aligned}
&\mathbf{x}^{(r)}_{k+1}=\mathbf{A}_r\,\mathbf{x}^{(r)}_{k}+\mathbf{B}_r\,\mathbf{u}^{(r)}_{k},\\
&\mathbf{x}^{(r)}_{\mathrm{lb}}\le \mathbf{x}^{(r)}_k\le \mathbf{x}^{(r)}_{\mathrm{ub}},\qquad
\mathbf{u}^{(r)}_{\mathrm{lb}}\le \mathbf{u}^{(r)}_k\le \mathbf{u}^{(r)}_{\mathrm{ub}},\\
&\mathbf{x}^{(r)}_0=\mathbf{x}^{(r)}_{\mathrm{init}},\\
&h_{i,r}\!\left(\mathbf{x}^{(r)}_{k+1}\right)\ge(1-\gamma)\,
h_{i,r}\!\left(\mathbf{x}^{(r)}_{k}\right),\quad i\in\mathcal{I}^{+}_k.
\end{aligned}
\label{eq:nlp_cons}
\end{equation}
其中 \(\mathbf{x}^{(r)}_{\mathrm{lb}},\mathbf{x}^{(r)}_{\mathrm{ub}}\) 为状态上下界（位置边界由地图范围给定，速度上界由不同平台的最大速度给定），\(\mathbf{u}^{(r)}_{\mathrm{lb}},\mathbf{u}^{(r)}_{\mathrm{ub}}\) 为控制输入上下界（加速度上界由平台最大加速度给定），\(\mathbf{x}^{(r)}_{\mathrm{init}}\) 由上一时刻的状态估计与 ESKF 输出提供，\(\gamma\in(0,1)\) 为 DCBF 衰减系数，\(\mathcal{I}^{+}_k\) 为在时刻 \(k\) 由高度窗口激活的椭圆柱索引集合，\(h_{i,r}(\cdot)\) 为第 \(i\) 个障碍相对于机器人 \(r\) 的距离型控制障碍函数，其构造见\autoref{sec:dcbf}。

为提高收敛速度并稳定实时性，本文在滚动时域执行中采用热启动策略构造优化变量的初始猜测。若为首次求解，即当前周期 \(t=0\)，则对每个机器人 \(r\) 与每个预测步 \(k\) 取
\[
\mathbf{p}^{(r),0}_k=\mathbf{p}^{(r)}_{\mathrm{init}},\qquad
\mathbf{v}^{(r),0}_k=\mathbf{0},\qquad
\mathbf{u}^{(r),0}_k=\mathbf{0},
\]
据此组装 \(\mathbf{x}^{(r),0}_k=[\,\mathbf{p}^{(r),0}_k{}^\top\ \mathbf{v}^{(r),0}_k{}^\top\,]^\top\)。当从上一规划周期 \((t-1)\) 已得到可行解
\(\{\widehat{\mathbf{x}}^{(r)}_{j|t-1},\widehat{\mathbf{u}}^{(r)}_{j|t-1}\}_{j=0}^{H}\) 时，在当前周期 \(t\) 通过时间平移进行热启动：对 \(k=0,\dots,H-1\) 令
\[
\mathbf{x}^{(r),0}_k=\widehat{\mathbf{x}}^{(r)}_{k+1|t-1},\qquad
\mathbf{u}^{(r),0}_k=\widehat{\mathbf{u}}^{(r)}_{k+1|t-1},
\]
并将末端步长缺失的量以上一周期的终端解填充。最后将 \(\mathbf{x}^{(r),0}_k,\mathbf{u}^{(r),0}_k\) 投影到各自的约束区间，以保证初始猜测满足边界条件。该热启动方式在状态与控制上保留了上一周期的时空结构，并与新的观测和引导点对齐，从而显著减少求解迭代次数并提升滚动优化的实时性。

\subsection{基于滚动时域的实时优化策略}
在上述非线性优化模型的基础上，本文采用滚动时域策略在每个规划周期在线求解\eqref{eq:nlp_obj}–\eqref{eq:nlp_cons} 所定义的问题，并将优化解转化为各机器人的连续可执行轨迹与控制命令。优化变量、等式约束和不等式约束在 CasADi 优化器中以符号函数形式构建，目标函数由 \(J_{\mathrm{wp}}\)、\(J_{\mathrm{vis}}\)、\(J_{\mathrm{form}}\) 和 \(J_{\mathrm{ctrl}}\) 组成，所有项的梯度以及需要时的拉格朗日 Hessian 由自动微分给出。数值求解采用 IPOPT 的原始–对偶内点法，启用精确雅可比与 Hessian，并配合上一周期解的热启动以提高迭代收敛速度，系统运行流程如\autoref{fig:system_framework_1}所示。

\begin{figure}[htbp]
    \centering
    \includegraphics[width=1\linewidth]{chapter_3/系统框架图.png} 
    \caption{滚动时域优化流程图}
    \label{fig:system_framework_1}
\end{figure}

在每个规划周期内，执行流程如下。首先利用最新的感知结果更新三维占据栅格，并基于点云聚类结果更新最小包络椭圆柱集合，由此重构 DCBF 所需的安全集与几何距离函数。随后在当前体素地图上运行混合 A*搜索，在考虑动力学约束和障碍物的前提下得到一条无碰撞的引导路径，并从中选取位于当前感知范围边界附近的最远可达点 \(w^\star\) 作为本周期的 UAV 引导目标点，同时按照新的位姿与障碍物分布更新可见性和编队相关的代价项。接着将 ESKF 输出的当前状态作为 \(\{\mathbf{x}^{(r)}_0\}\) 的初值，并通过时间平移方式将上一规划周期得到的可行解 \(\{\widehat{\mathbf{x}}^{(r)}_{k|t-1},\widehat{\mathbf{u}}^{(r)}_{k|t-1}\}\) 映射到当前预测域，形成 \(\{\mathbf{x}^{(r),0}_k,\mathbf{u}^{(r),0}_k\}\) 作为 IPOPT 的初始猜测。求解器在满足动力学约束、输入上下界约束和 DCBF 约束的前提下，输出当前周期的最优状态与控制序列 \(\{\widehat{\mathbf{x}}^{(r)}_{k},\widehat{\mathbf{u}}^{(r)}_{k}\}_{k=0}^{H}\)。

为在执行层获得平滑可跟踪的连续轨迹，本文对每个机器人在每个坐标轴上采用五次多项式进行轨迹时间参数化。设当前控制周期对应的时间区间为 \([0,\Delta t]\)，从优化解中读取该区间起点和终点的状态
\(\{\mathbf{p}^{(r)}_0,\mathbf{v}^{(r)}_0,\mathbf{u}^{(r)}_0\}\) 与
\(\{\mathbf{p}^{(r)}_1,\mathbf{v}^{(r)}_1,\mathbf{u}^{(r)}_1\}\)。
对任一标量分量构造
\[
p^{(r)}(t)=\alpha_0+\alpha_1 t+\alpha_2 t^2+\alpha_3 t^3+\alpha_4 t^4+\alpha_5 t^5,\quad t\in[0,\Delta t],
\]
以位置、速度和加速度在 \(t=0\) 与 \(t=\Delta t\) 处的六个边界条件唯一确定系数向量 \(\boldsymbol{\alpha}=[\alpha_0,\ldots,\alpha_5]^\top\)。该构造保证轨迹在位置、速度和加速度上均为 \(C^2\) 连续，使离散优化解在时间上自然平滑。随后以控制频率在多项式上采样得到 \(\{p^{(r)}(t),\dot p^{(r)}(t),\ddot p^{(r)}(t)\}\)，作为参考信号分别下发给 UAV 与 UGV 的底层控制器进行跟踪控制。

该滚动时域策略的核心在于“感知—预测—执行”。在每个规划周期内，实际执行当前时间区间 \([0,\Delta t]\) 内对应的多项式轨迹，将尚未执行的部分作为下一周期的热启动，并在新的感知信息、障碍物集合和引导点 \(w^\star\) 到来后重新构建并求解优化问题。借助 CasADi 的自动微分能力和 IPOPT 的高效内点法，该策略能够在多目标权重的调节下，持续产生满足动力学可行性与安全约束的空地协同轨迹，从而在未知、拥挤和缓慢变化的环境中实现在线协同规划与控制。

% 尽量增加一个实物实验
\section{仿真实验} \label{sec:sim}
本节将通过构建仿真平台对所提出的基于点云聚类的协同规划方法进行性能评估，实验重点在于验证该方法在未知障碍物环境下生成协同轨迹时系统内部可视情况和避障安全性。在不同的障碍物密度情况下对多优化目标构建的有效性进行分析，并设置消融实验来验证。

\subsection{实验环境设置}

为评估所提出协同规划方法在不同环境复杂度下的性能，本文在尺寸为 \(30\times30\times3\,\mathrm{m}^3\) 的三维仿真场地内构建两类场景：稀疏障碍场景与稠密障碍场景。体素分辨率均为 \(0.1\,\mathrm{m}\)，障碍物在整个场地范围内随机布置，满足稀疏环境中障碍物间的平均间距约为 \(2.0\,\mathrm{m}\)，稠密环境中障碍物间的平均间距约为 \(1.2\,\mathrm{m}\)。空地系统由 \(1\) 架 UAV 与 \(4\) 辆 UGV 组成。初始时刻，UAV 位于 \((0,0)\,\mathrm{m}\) 上方，四辆 UGV 的位置为\(\mathrm{UGV}_0(-1,1)\)、\(\mathrm{UGV}_1(-1,-1)\)、\(\mathrm{UGV}_2(1,-1)\)、\(\mathrm{UGV}_3(1,1)\)。UAV 全局目标点设在场景中 \((32,0)\,\mathrm{m}\) 附近，引导空地编队穿越整个障碍区域。仿真环境基于Gazebo搭建，UAV的仿真模型基于PX4仿真器实现，UGV的模型基于Nexus-4WD-Mecanum 仿真器进行修改，并在其上搭载 MID-360 LiDAR 以模拟真实传感器数据采集，仿真环境如 \autoref{fig:sim_env} 所示。

\begin{figure}[htbp]
    \centering
    \includegraphics[width=0.75\linewidth]{chapter_3/仿真环境.png} 
    \caption{Gazebo仿真环境示意图}
    \label{fig:sim_env}
\end{figure}

非线性优化问题基于 CasADi 的符号函数构建，目标函数与约束通过自动微分进行计算，并由 IPOPT 的对偶内点法求解。规划循环以固定频率触发，底层控制器以更高频率采样五次多项式轨迹并生成底层控制命令。实际仿真中，单次 NMPC 求解在平均与最坏情况下均能满足规划频率要求，保证了滚动时域协同优化的实时性。关键优化与求解相关参数汇总如\autoref{tab:opt_params} 所示。

\begin{table}[htbp]
    \centering
    \caption{协同优化与求解相关参数设置}
    \label{tab:opt_params}
    \begin{tabularx}{\linewidth}{c X<{\centering}}
        \toprule
        参数 & 数值 \\
        \midrule
        采样时间 $\Delta t$ & $0.2\,\mathrm{s}$ \\
        预测步长 $H$ & $10$ \\
        UAV 速度/加速度限幅 & $2.5\,\mathrm{m/s}$，$3.0\,\mathrm{m/s^2}$ \\
        UGV 速度/加速度限幅 & $1.4\,\mathrm{m/s}$，$1.0\,\mathrm{m/s^2}$ \\
        碰撞判定安全半径 & $d_{\mathrm{safe}}=0.3\,\mathrm{m}$ \\
        DCBF 衰减系数 & $\gamma=0.15$ \\
        UGV之间的距离 & $[d_{\mathrm{low}},d_{\mathrm{high}}]=[4.0,8.0]\,\mathrm{m}$ \\
        UAV–UGV 距离 & 最大允许距离 $6.0\,\mathrm{m}$\\
        UAV 路点跟踪权重 & $\lambda_{\mathrm{wp}}^{\min}=1.0$，$\lambda_{\mathrm{wp}}^{\max}=4.0$\\
        可见性代价权重 & $\lambda_{\mathrm{vis}}=1.2$ \\
        UGV 距离代价权重 & $\lambda_{\mathrm{dis}}=1.0$ \\
        UAV–UGV 距离代价权重 & $\lambda_{\mathrm{bound}}=1.4$ \\
        队形形状代价权重 & $\lambda_{\mathrm{shape}}=0.9$ \\
        控制能量与控制增量权重 & $\lambda_{\mathrm{ctrl}}=0.5$，$\lambda_{\Delta u}=0.2$ \\
        IPOPT 最大迭代次数 & $500$ \\
        IPOPT 收敛容忍度 &  $10^{-5}$ \\
        优化问题重规划频率 & $20\,\mathrm{Hz}$ \\
        底层控制频率 & $100\,\mathrm{Hz}$ \\
        单次求解平均耗时 & $0.04\,\mathrm{s}$ \\
        \bottomrule
    \end{tabularx}
\end{table}

\subsection{实验结果分析}

本节在稀疏与稠密两类障碍物环境下，对比“启用可见性约束（Full）”与“去掉可见性约束（NoVis）”两种配置的协同规划效果。\autoref{fig:sparse_traj_compare} 和 \autoref{fig:dense_traj_compare} 展示了两类环境中 UAV 与四辆 UGV 的平面轨迹，同时 \autoref{tab:sparse_stat} 与 \autoref{tab:dense_stat} 分别给出了平均可视率、队形偏差、任务成功率以及路径长度等统计指标。平均可视率定义为在任务执行期间，UAV 与每一辆 UGV 之间的连线在无任一障碍物遮挡条件下保持可见的时间占整段运行时间的比例，对四辆 UGV 分别计算后取算术平均作为该次试验的综合可视率指标。任务成功率定义为在多次重复试验中，所有机器人在整个运行过程中均未发生与障碍物或其他机器人碰撞的次数占总试验次数的比例。队形偏差中，“径向偏差”度量各 UGV 相对理想矩形顶点的距离偏差，“角度偏差”度量四边形内角相对 $90^\circ$ 的偏差，“边长偏差”度量相邻边长度相对理想边长的偏差；轨迹长度一列给出了 UAV 路径长度以及四辆 UGV 路径长度的平均值。

\autoref{fig:sim_runtime_record} 直观展示了本章协同规划方法在复杂障碍场景下的在线规划过程与空地协同关系。图中，基于混合 A* 生成的动力学可行全局参考路径以蓝色曲线给出，用于从当前位置指向远端目标点；当前规划周期内优化得到的 UAV 局部参考轨迹以红色点列表示，对应 UAV 在滚动时域中的待跟踪轨迹；四辆 UGV 的目标轨迹以青色点列给出，体现出其在 UAV 领航下的协同推进效果。场景中的彩色点云及其外层的几何包络体对应前文通过聚类与最小包络椭圆柱得到的障碍物表示，反映了环境感知与凸集建模的过程。同时连接UAV与UGV之间的绿色直线表示两者之间是可见的，如果不可见则以红色直线表示。通过该图可以清晰看到，在点云增量更新与障碍物动态重建的条件下，空地多机器人系统如何依托混合 A* 引导与非线性优化联合作用，实现对轨迹的实时重规划与协同运动。

\begin{figure}[htbp]
    \centering
    \includegraphics[width=0.75\linewidth]{chapter_3/runtime_record2.png} 
    \caption{空地系统协同规划过程示意图}
    \label{fig:sim_runtime_record}
\end{figure}

从稀疏障碍物环境的结果（\autoref{fig:sparse_traj_compare} 与 \autoref{tab:sparse_stat}）可以看出，两种配置下 UAV 与 UGV 均能绕开障碍顺利到达全局目标点，任务成功率均为 $100\%$，轨迹长度相近：UAV 约为 $34\sim 35\,\mathrm{m}$，UGV 平均约为 $28\sim 29\,\mathrm{m}$。启用可见性约束时，平均可视率略高（$0.87$ 对比 $0.85$），说明在稀疏场景中可见性代价能够在不显著增加路径长度的前提下，轻微提升空地之间的视距质量。然而，由于障碍分布较为宽松，系统有足够的空间同时满足避障和队形约束，去掉可见性项后，四车之间的径向和角度偏差反而更小（径向偏差由 $0.24\,\mathrm{m}$ 降至 $0.13\,\mathrm{m}$，角度偏差由 $16.5^\circ$ 降至 $8.6^\circ$），多边形形状略为更规整。这表明在障碍稀疏时，可见性项对 LoS 的增益有限，而队形项本身已经足以保持较好的几何结构。

\begin{figure}[htbp]
  \centering
  \begin{subfigure}{0.95\linewidth}
    \centering
    \includegraphics[width=\linewidth]{chapter_3/sparse_env_traj_full_constrain.png}
    \caption{启用可见性约束}
    \label{fig:sparse_full}
  \end{subfigure}
  \vskip 0.5em
  \begin{subfigure}{0.95\linewidth}
    \centering
    \includegraphics[width=\linewidth]{chapter_3/sparse_env_traj_noVisConstrain.png}
    \caption{去掉可见性约束}
    \label{fig:sparse_novis}
  \end{subfigure}
  \caption{稀疏障碍环境下的空地系统轨迹对比}
  \label{fig:sparse_traj_compare}
\end{figure}

\begin{table}[htbp]
  \caption{稀疏障碍物环境下不同约束配置的统计结果}
  \label{tab:sparse_stat}
  \centering
  \begin{tabularx}{\linewidth}{X<{\centering} c c c c}
    \toprule
    优化问题 & 可视率 & 队形偏差 & 成功率 & 轨迹长度 (m)\\
    \midrule
    启用可见性约束 &
    0.87 &
    $0.24 / 16.5^\circ / 0.20$ &
    100\% &
    $34.1 / 28.7$ \\
    去掉可见性约束 &
    0.85 &
    $0.13 / 8.6^\circ / 0.20$ &
    100\% &
    $34.7 / 28.3$ \\
    \bottomrule
  \end{tabularx}
\end{table}

在稠密障碍物环境下（\autoref{fig:dense_traj_compare} 与 \autoref{tab:dense_stat}），情况有所不同。首先，两种配置下 UAV 的路径长度依然维持在约 $35\,\mathrm{m}$，UGV 平均路径在 $28\sim 29\,\mathrm{m}$ 左右，说明即便在更拥挤的场景中，优化器仍然能够找到长度相近的可行轨迹。但从可视率和成功率可以看出可见性约束的作用：启用可见性项后，平均可视率提高到 $0.77$（相较去掉可见性时的 $0.66$），空地之间的视距遮挡明显减少；同时队形偏差也略有改善，径向偏差从 $0.21\,\mathrm{m}$ 降至 $0.15\,\mathrm{m}$，角度偏差从 $12.1^\circ$ 降至 $10.7^\circ$，边长偏差略有下降。这说明在障碍密集、视线容易被遮挡的条件下，可见性代价能够引导 UAV 与 UGV 在障碍缝隙间协同机动，使队形既保持基本矩形结构，又尽量避免 LoS 穿过障碍物。

\begin{figure}[htbp]
  \centering
  \begin{subfigure}{0.95\linewidth}
    \centering
    \includegraphics[width=\linewidth]{chapter_3/dense_env_traj_full_constrain.png}
    \caption{启用可见性约束}
    \label{fig:dense_full}
  \end{subfigure}
  \vskip 0.5em
  \begin{subfigure}{0.95\linewidth}
    \centering
    \includegraphics[width=\linewidth]{chapter_3/dense_env_traj_noVisConstrain.png}
    \caption{去掉可见性约束}
    \label{fig:dense_novis}
  \end{subfigure}
  \caption{稠密障碍环境下的空地系统轨迹对比}
  \label{fig:dense_traj_compare}
\end{figure}

\begin{table}[htbp]
  \caption{稠密障碍物环境下不同约束配置的统计结果}
  \label{tab:dense_stat}
  \centering
  \begin{tabularx}{\linewidth}{X<{\centering} c c c c}
    \toprule
    优化问题 & 可视率 & 队形偏差 & 成功率 & 轨迹长度 (m)\\
    \midrule
    启用可见性约束 &
    0.77 &
    $0.15 / 10.7^\circ / 0.29$ &
    20\% &
    $34.0 / 29.0$ \\
    去掉可见性约束 &
    0.66 &
    $0.21 / 12.1^\circ / 0.31$ &
    60\% &
    $34.9 / 28.6$ \\
    \bottomrule
  \end{tabularx}
\end{table}

另一方面，稠密环境中启用可见性约束的任务成功率从 $60\%$ 降至 $20\%$。这反映了在高密度障碍场景下，多目标之间存在更强的冲突：为了同时兼顾避障、安全距离、队形保持和视距质量，局部可行域被进一步压缩，部分仿真中会出现优化不可行或无法在限定时间内收敛的情况；而不考虑可见性时，可行域更大、成功率略高，但以更低的平均可视率和更差的队形质量为代价。结合轨迹图可以观察到，去掉可见性约束时 UAV 更倾向于沿局部最短路径穿越障碍“缝隙”，导致部分时间段内空地之间被障碍完全遮挡；启用可见性后，UAV 会在横向上适度绕行，以保证绝大多数时间内至少有一条视距通道保持畅通。

综合两类环境的结果，可以认为在障碍较稀疏时，可见性项对整体路径和队形的影响较为温和，但仍能在不降低成功率的前提下略微提升可视率；在障碍密集、视距容易被遮挡的场景中，可见性约束显著提升了空地系统间的平均可视率和队形几何质量，同时带来一定的成功率折损。总体上，所提出的多目标协同优化框架能够在路径长度、避障安全、队形保持和空地可见性之间实现可调节的权衡，为后续在真实复杂环境中的在线部署提供了参数选择与性能预期的依据。

\section{本章小结} \label{sec:summary_opt1}

本章围绕“轻量环境表达与多目标统一优化”的思路，提出了一种面向空地多机器人系统的协同轨迹规划方法，并在仿真环境中验证了其有效性。首先，在动力学层面针对空地异构平台进行了统一建模。对于 UAV，基于其微分平坦特性，以质心位置和偏航角作为平坦输出，给出了由平坦输出到姿态、推力和力矩的解析映射；对于麦克纳姆轮 UGV，则从全向运动学关系出发，说明了通过机体系平移速度即可在平面内实现任意方向运动。在此基础上，将 UAV 的三维运动与 UGV 的二维平面运动统一抽象为以“位置–速度”为状态、以“加速度”为控制的二阶积分离散模型，并施加平台相关的速度与加速度限幅约束，为后续协同优化提供了统一的状态与控制变量。  

在环境建模方面，本章提出了基于 DBSCAN 聚类与最小包络椭圆柱的障碍物表达方法。通过对共享点云进行密度聚类，将障碍物从噪声中分离，并在地平面上对每个聚类簇求解最小包络椭圆，再结合高度范围扩展为椭圆柱，以少量几何参数紧凑描述复杂环境。基于该凸集表示，构造了到椭圆边界的解析距离及其梯度，并与离散控制障碍函数结合，形成可直接嵌入 NMPC 的几何避障约束，在保证安全集前向不变性的同时兼顾数值可解性。为获得具有动力学一致性的引导信息，本文利用混合 A* 在当前栅格地图上搜索满足速度与加速度约束的无碰撞路径，从中选取位于感知边界附近的局部目标点作为 UAV 在预测域内的引导路点，从而避免全局目标点远离可见区域时引发的“盲目追踪”。  

在多目标协同优化方面，本章在统一动力学和环境表示的基础上，建立了“UAV 领航、UGV 跟随”的非线性模型预测控制框架。代价函数同时包含 UAV 对引导路点的跟踪项、空地之间基于椭圆柱几何的可见性代价、四辆 UGV 的队形尺度与形状代价以及控制能量与控制增量的正则项，并额外加入 UAV 与各 UGV 之间的距离约束以抑制队形过度拉伸。通过调整各子目标的权重，可以在路径长度、可视质量、编队稳定性和能耗之间实现可调的折中。整个优化问题在 CasADi 中以符号形式构建，采用 IPOPT 的原始–对偶内点法在滚动时域内实时求解，并利用五次多项式对相邻采样时刻的解进行时间参数化，将离散解转化为平滑、可执行的连续参考轨迹。  

仿真结果表明，该方法在 \(30\times 30\times 3\,\mathrm{m}^3\) 的稀疏和稠密障碍环境中均能完成一架 UAV 与四辆 UGV 的协同穿越任务，单次求解时间保持在实时范围内。对比启用和去除可见性约束的结果可以看出，在稀疏场景下，可见性约束在几乎不增加路径代价的情况下略微提升了平均可视率；在稠密场景下，可见性约束显著改善了空地系统间的视距联系，同时对成功率和轨迹长度带来一定影响，体现出在可视质量与通过效率之间存在明确的权衡空间。总体来看，本章提出的基于点云聚类、椭圆柱凸表示和离散控制障碍函数的协同规划框架，能够在复杂三维障碍环境中生成安全、保持队形并兼顾可见性的空地协同轨迹。


% 复杂环境规划
\chapter{基于序列可视空间的两阶段协同规划方法研究}

\section{引言}

在复杂未知环境中实现空地多机器人系统的协同规划，需要在有限算力条件下同时兼顾动力学可执行性、空地可视性和地面编队稳定性，并保证整体规划的实时性。若在全局尺度上直接以占据栅格为约束，将避障约束、可视性约束和多机耦合统一纳入单一优化问题，则状态与控制维度会随时间和机器人数量急剧增长，约束数量也随障碍物和可视性关系极速扩张，容易导致优化问题难收敛、数值求解不稳定等问题。

上一章在点云聚类与最小包络椭圆柱近似的基础上，构建了基于椭圆柱几何表示和离散控制障碍函数的协同规划方法。该方法在中等稀疏环境下可以在统一 NMPC 框架中同时处理引导跟踪、避障、安全距离和编队几何等多目标，并在仿真中展现出良好的实时性。然而在障碍更加稠密或几何更加复杂的场景下，直接在全局优化问题中对每一个椭圆柱施加 DCBF 约束，会带来以下问题：其一，可行域在时间和空间上高度碎片化，导致优化容易陷入局部最优甚至不可行；其二，可见性与编队等长期几何目标与避障在同一层面进行协调，权重调节空间变窄，稳定性与可解释性下降；其三，当任务尺度或环境复杂度进一步增大时，约束数量随障碍和时间步线性增加，单次求解时间和收敛鲁棒性难以保证，从而限制了方法在更大规模、更加稠密环境中的泛化能力。

为缓解上述问题并提升在稠密场景下的稳定性和对复杂环境的泛化能力，本文在上一章方法的基础上提出两阶段协同规划框架，将“全局一致性目标”和“局部安全执行目标”在时间尺度和决策层次上进行解耦。第一阶段为全局规划阶段，在共享占据栅格和统一二阶积分动力学模型的基础上，采用混合 A* 搜索生成满足速度与加速度限幅的全局引导路径，使参考轨迹在动力学上可跟踪且与任务目标方向一致。沿该引导路径进行等时采样，在每一个采样点构建包含可见性与编队约束的可视空间，并按时间顺序将这些可视空间组织为序列可视空间（Sequential Star-Convex Polytope, SSCP）：
\begin{equation}
    \mathrm{SSCP} = \big\{ \mathcal{V}_k \subset \mathbb{R}^3 \,\big|\, k = 0,\dots,K \big\},
    \label{eq:sscp_def}
\end{equation}
其中每个 $\mathcal{V}_k$ 为一处在自由空间内部的可视集，描述了 UAV 与 UGV 在该时刻应满足的视距关系。SSCP 的形式定义如 \autoref{eq:sscp_def} 所示。由于 SSCP 已通过环境几何和占据信息排除了碰撞体素，全局阶段不再显式施加细粒度避障约束，而是将 UAV 渐近逼近引导路径、空地可见性保持以及空地队形几何稳定性作为可微代价和软约束，在滚动时域内进行中优化，从而保持问题规模紧凑并提升数值收敛稳定性。

第二阶段为局部规划与跟踪阶段，面向本地执行与安全保证。全局阶段输出的其对应的离散参考位姿序列在每个控制周期内被时间参数化为 MINCO \cite{wang2022geometrically}多项式轨迹，局部跟踪器在短预测域上最小化对参考序列的同步误差，并在同一层面处理避障与控制能量问题。可视保持和队形几何主要由全局规划负责，局部优化则专注于避障和能量开销控制，两者通过滚动时域在时间上紧密衔接。

基于上述设计，本章结构安排如下：\autoref{sec:vs+kino} 介绍序列可视空间的构建方法，包括可视空间的构建方法以及基于混合 A* 搜索的引导路径和序列可视空间生成步骤；\autoref{sec:sscp-global} 在 SSCP 上建立全局实时协同轨迹生成模型，给出优化变量初始化策略、无人机引导路点约束、序列可视空间约束以及完整的优化问题建模与求解方法；\autoref{sec:local-tracking} 围绕局部跟踪与避障，构建路径点同步跟踪约束、局部轨迹避障方法和最小能量轨迹优化模型；\autoref{sec:two-stage-arch} 从系统角度总结两阶段协同规划方法的整体架构，说明全局规划与局部执行之间的信息流；\autoref{sec:experiments} 给出仿真与实物实验结果，对比分析与上一章方法在稠密场景下的性能差异，验证两阶段框架在稳定性与泛化能力方面的优势；最后在 \autoref{sec:summary_opt2} 对本章工作进行总结。

\section{序列可视空间的构建方法}\label{sec:vs+kino}
本节给出序列可视空间（SSCP）的构建方法。首先在共享占据栅格和统一动力学模型的基础上，利用上一章第~\ref{sec:cluster_path} 节中基于混合 A* 的方法，在全局地图上得到满足速度与加速度限幅的动力学可行引导路径。该路径以无人机平台的机动能力为约束，保证每一段都可由后续执行层稳定跟踪。然后对该引导路径按照等时间间隔进行采样，形成与滚动预测域一致的离散参考序列。围绕每一个采样位姿，基于局部点云构建对应的可视空间，用于刻画该时刻 UAV 与 UGV 所应处于的“视距良好且无障碍遮挡”的空间区域。最后将这些按时间排序的可视空间串联，即得到序列可视空间 \(\mathrm{SSCP}\)，作为全局协同优化中刻画可视性的核心约束表示。

\subsection{可视空间的构建}
本小节围绕引导路径等时采样得到的参考点 \(c_i\in\mathbb{R}^3\) 构建其附近的可视空间，用于在全局层面以可微方式约束可视性与遮挡。\autoref{fig:scp_construct} 以二维视角展示了单个可视空间的构建流程。图中 \(p(t)\) 与本节符号 \(c_i\) 对应，为当前引导点；半径 \(R\) 的红色虚线圆给出局部点云检索半径；更大的半径 \(r\) 为球翻转半径。绿色实心点为局部采样点，红色实心点为在球面上增广的人工点，二者组成局部点集；这些点经球翻转后得到蓝色点云，其凸包（蓝色多边形）近似可视边界。绿色星形多边形表示反变换回原空间后的星凸可视区域，位于该区域内部的点对中心 \(p(t)\) 可见，而图中箭头所示灰色多边形障碍物遮挡了部分方向，使得某些点（黄色五角星）落在可视区域之外。图中 \(x\) 与 \(\hat x\) 分别表示原空间和翻转空间中的对应点，\(d_i\) 表示点到可视边界的距离。

\begin{figure}[htbp]
    \centering
    \includegraphics[width=0.7\linewidth]{chapter_4/scp.png}
    \caption{单个可视空间的构建示意图}
    \label{fig:scp_construct}
\end{figure}

在统一世界坐标系下，记半径 \(R>0\) 球邻域内的局部点集为
\[
M_v=\{m_j\}\subset\mathbb{R}^3.
\]
以参考点 \(c_i\) 为球心，在半径 \(R\) 邻域内从点云中检索观测点，并在球面上均匀增广若干人工采样点，用于填补稀疏区域。定义球翻转变换 \(F:\mathbb{R}^3\to\mathbb{R}^3\)。对任意点 \(x\)，记
\begin{equation}
y=x-c_i,\qquad r>R,
\end{equation}
其中 \(r\) 为翻转半径，则球翻转定义为
\begin{equation}
\hat{x}=F(x)
=y+2\big(r-\|y\|_2\big)\frac{y}{\|y\|_2},
\label{eq:ball_flip}
\end{equation}
当点靠近 \(c_i\) 时被推向半径 \(r\) 附近，当点接近半径 \(R\) 时变换趋近恒等。对 \(M_v\) 中的每个点做变换 \(\hat{m}_j=F(m_j)\)，在翻转空间中对所有 \(\hat{m}_j\) 计算凸包，得到多面体的 \(H\)-表示：
\begin{equation}
\mathcal{P}_i=\big\{\hat{x}\in\mathbb{R}^3\ \big|\ A\hat{x}\le b\big\},\qquad
A=\begin{bmatrix}n_1^\top\\ \vdots\\ n_K^\top\end{bmatrix},\quad
b=\begin{bmatrix}n_1^\top a_1\\ \vdots\\ n_K^\top a_K\end{bmatrix},
\label{eq:polytope_H}
\end{equation}
其中 \(K\) 为面数，\(n_k\in\mathbb{R}^3\) 为第 \(k\) 个面的外法向量，\(a_k\in\mathbb{R}^3\) 为该面上一点。对任意 \(\hat{x}\)，其相对于第 \(k\) 个面的有符号支持函数为
\begin{equation}
d_k(\hat{x})=n_k^\top(\hat{x}-a_k),\qquad k=1,\dots,K,
\label{eq:face_support}
\end{equation}
当 \(d_k>0\) 时点在面外侧，当 \(d_k<0\) 时在多面体内部。

为获得对 \(\max_k d_k(\hat{x})\) 的光滑近似方便求导，引入如下形式
\begin{equation}
\Xi(\hat{x})
=\frac{1}{\alpha}\log\!\Big(\sum_{k=1}^{K}e^{\alpha d_k(\hat{x})}\Big),
\qquad \alpha>0,
\label{eq:lse_def}
\end{equation}
当 \(\alpha\) 增大时 \(\Xi(\hat{x})\) 渐近于 \(\max_k d_k(\hat{x})\)。原空间中的点 \(x\) 经 \(\hat{x}=F(x)\) 映射到翻转空间后，若其 \(\Xi(\hat{x})\) 取值较大，则说明点在可视边界外部或靠近边界。

为在优化中构建可视裕度，设定最小可视距离 \(d_{\min}\geq 0\)，定义可视性违背量
\begin{equation}
\phi(x\mid S_i)
=d_{\min}-\Xi\big(F(x)\big),
\label{eq:violation}
\end{equation}
其中 \(S_i\) 表示以 \(c_i\) 为中心的可视空间。当 \(\phi(x\mid S_i)\le 0\) 时，说明点 \(x\) 相对可视边界具有不小于 \(d_{\min}\) 的裕度；当 \(\phi(x\mid S_i)>0\) 时，说明点 \(x\) 距离可视边界不足 \(d_{\min}\)。据此构造分段光滑的可视性代价
\begin{equation}
V_{\mathrm{scp}}(x\mid S_i)=
\begin{cases}
0,& \phi(x\mid S_i)\le 0,\\[3pt]
\lambda_{\mathrm{scp}}\ \phi(x\mid S_i)^{3},& \phi(x\mid S_i)>0,
\end{cases}
\label{eq:scp_penalty}
\end{equation}
其中 \(\lambda_{\mathrm{scp}}>0\) 为权重。三次多项式在边界处函数值与一阶导数连续，在违反程度增大时惩罚随 \(\phi^3\) 快速增长，有利于在保持数值平滑的同时提供足够的“推离”趋势。

为便于在优化中使用与求解，可进一步给出 \autoref{eq:scp_penalty} 的梯度形式。由 \autoref{eq:lse_def} 可得
\begin{equation}
w_k(\hat{x})=\frac{e^{\alpha d_k(\hat{x})}}{\sum_{j=1}^{K}e^{\alpha d_j(\hat{x})}},\qquad 
\bar{n}(\hat{x})=\sum_{k=1}^{K}w_k(\hat{x})\,n_k,
\label{eq:softmax_w}
\end{equation}
从而
\begin{equation}
\nabla_{\hat{x}}\Xi(\hat{x})=\bar{n}(\hat{x}).
\label{eq:grad_xi_hat}
\end{equation}
由 \autoref{eq:ball_flip}，球翻转的雅可比为
\begin{equation}
J_F(x)=\frac{\partial F(x)}{\partial x}
=\Big(\frac{2r}{\|y\|_2}-1\Big)\mathbf{I}_3-\frac{2r}{\|y\|_2^{3}}\,y\,y^\top,
\qquad y=x-c_i,
\label{eq:Jflip}
\end{equation}
其中 \(\mathbf{I}_3\) 为三维单位阵。由链式法则与 \autoref{eq:violation}、\autoref{eq:grad_xi_hat}、\autoref{eq:Jflip} 得
\begin{equation}
\nabla_x \phi(x\mid S_i)= -\,J_F(x)^\top\,\bar{n}\big(F(x)\big).
\label{eq:grad_phi}
\end{equation}
进而由 \autoref{eq:scp_penalty} 的分段定义可得
\begin{equation}
\nabla_x V_{\mathrm{scp}}(x\mid S_i)=
\begin{cases}
\mathbf{0},& \phi(x\mid S_i)\le 0,\\[3pt]
3\,\lambda_{\mathrm{scp}}\ \phi(x\mid S_i)^{2}\ \nabla_x \phi(x\mid S_i),& \phi(x\mid S_i)>0,
\end{cases}
\label{eq:grad_Vscp}
\end{equation}
代入 \autoref{eq:grad_phi} 可写成显式形式
\begin{equation}
\nabla_x V_{\mathrm{scp}}(x\mid S_i)=
\begin{cases}
\mathbf{0},& \phi(x\mid S_i)\le 0,\\[6pt]
-\,3\,\lambda_{\mathrm{scp}}\ \phi(x\mid S_i)^{2}\ J_F(x)^\top\,\bar{n}\big(F(x)\big),& \phi(x\mid S_i)>0,
\end{cases}
\label{eq:grad_Vscp_final}
\end{equation}
从而在给定 \(\{n_k,a_k\}_{k=1}^{K}\)、翻转半径 \(r\)、裕度 \(d_{\min}\) 和权重 \(\lambda_{\mathrm{scp}}\) 后，可以在优化中高效计算可视空间代价及其梯度。

为了直观展示基于上述流程得到的可视空间样式，\autoref{fig:vis_space_example} 给出了在不同环境和不同机器人参考位姿下构建的若干星凸可视区域示意。图中彩色多边形为各时刻的可视空间边界，灰色矩形为障碍物投影，可以看到可视区域自适应地在障碍物间“开洞”，并围绕参考点形成非规则但保持星凸性的自由区域。这些可视空间在时间上按引导路径顺序串联，即可构成后续全局优化所使用的序列可视空间。

\begin{figure}[htbp]
    \centering
    \includegraphics[width=0.6\linewidth]{chapter_4/vis_space_example.png}
    \caption{在不同参考位姿与环境下构建的可视空间示例}
    \label{fig:vis_space_example}
\end{figure}

\subsection{基于引导路径的序列可视空间生成}
上一章在第~\ref{sec:cluster_path} 节中基于混合 A* 在占据栅格上搜索得到从当前 UAV 位置至任务目标的三维动力学可行引导路径，记该路径的连续时间参数化为 \(\mathbf{p}_{\mathrm{ref}}(t)\)。在本章中将 \(\mathbf{p}_{\mathrm{ref}}(t)\) 视为引导路径序列，用于驱动可视空间的时序化构建，并用于可视代价的建立。

为与后续优化问题在时间上对齐，设预测步长为 \(\Delta t\)、预测长度为 \(H\)。在时间区间 \([t_1,t_H]\) 上按等间隔采样得到引导路点序列
\begin{equation}
    \mathrm{WS}
    =\big[\mathbf{p}_{\mathrm{ref}}(t_1),\mathbf{p}_{\mathrm{ref}}(t_2),\ldots,\mathbf{p}_{\mathrm{ref}}(t_H)\big],
    \qquad t_{k+1}-t_k=\Delta t,
    \label{eq:sscp_ws}
\end{equation}
其中 \(\mathrm{WS}\)（Waypoint Sequence）中的每一个 \(\mathbf{p}_{\mathrm{ref}}(t_k)\) 都对应预测域内的一个离散时间步。对每个采样点，将其记为可视空间核
\begin{equation}
    c_k=\mathbf{p}_{\mathrm{ref}}(t_k),\qquad k=1,\dots,H.
\end{equation}
在以 \(c_k\) 为中心、半径 \(R\) 的邻域内检索局部点集，并按照上一小节所述“点集检索与增广—球翻转—凸包—光滑代价”的流程构建以 \(c_k\) 为中心的可视空间 \(S_k\)。具体而言，在每个 \(k\) 上：首先从占据栅格与增量点云中提取局部点集并在球面上增广均匀采样点；随后依据 \autoref{eq:ball_flip} 对所有点进行球翻转，得到翻转点集；在翻转空间中计算凸包，获得如 \autoref{eq:polytope_H} 所示的 \(H\)-表示，并据此定义 \autoref{eq:lse_def}–\autoref{eq:scp_penalty} 的可视代价 \(V_{\mathrm{scp}}(x\mid S_k)\) 及其梯度；最后将 \((c_k,\mathcal{P}_k)\) 以及 \(R,r,d_{\min},\lambda_{\mathrm{scp}}\) 等参数打包，作为时刻 \(t_k\) 的可视空间描述。

\begin{figure}[htbp]
    \centering
    \includegraphics[width=0.75\linewidth]{chapter_4/sscp.png}
    \caption{序列可视空间示例}
    \label{fig:sscp_example}
\end{figure}

按时间顺序将所有可视空间组织为
\begin{equation}
    \mathrm{SSCP}
    =\big[S_1,S_2,\ldots,S_H\big],
    \label{eq:sscp_seq}
\end{equation}
它是上一节集合形式 \(\{\mathcal{V}_k\}\) 的具体实现，每个 \(S_k\) 对应一个以 \(c_k\) 为中心的星凸可视区域，最终构建的结果如\autoref{fig:sscp_example}所示。全局协同优化在预测域内对各 UGV 的位姿 \(\{p^{(r)}_k\}\) 在相应可视空间上评估代价
\(\sum_k V_{\mathrm{scp}}(p^{(r)}_k\mid S_k)\)，从而在不显式进行大量射线碰撞检测的前提下，用有限维、可微且梯度结构清晰的形式刻画沿引导路径随时间演化的视距保持要求。

\section{基于序列可视空间的全局实时协同轨迹生成} \label{sec:sscp-global}
本节在上一节给出的引导路径与序列可视空间基础上，构建面向全局层面的实时协同轨迹生成方法。与前一章的建模思路相比，本章在全局层面不再处理基于椭圆柱的障碍物避障约束，而是将优化重点集中在两类长期一致性目标上：其一是无人机对混合 A* 引导路点的跟踪，其二是在时间序列上维持空地系统的可视性与队形的几何稳定性。与局部避障与能量抑制等问题由后续局部轨迹生成模块负责，从而显著降低全局优化的维度和约束耦合度，更好地适应滚动时域下的实时求解需求。

\subsection{优化变量初始化与热启动方法}
本小节给出预测域内状态与控制向量的初始化与热启动策略。设预测步长为 \(H\)，采样间隔为 \(\Delta t\)。在每个时间步上，每辆 UGV 包含平面位置、速度、加速度三组二维变量，共 \(6\) 个标量；每架 UAV 包含三维位置、速度、加速度三组变量，共 \(9\) 个标量。记地面车辆数量为 \(N_g\)、空中平台数量为 \(N_a\)，则单步决策维数为 \(D=6N_g+9N_a\)，预测域总维数为 \(H D\)。为与引导路径和序列可视空间一致，初始化时序与混合 A* 引导路径的时间网格严格对齐。

首先为每个时间步构造“中心轨迹”。若第 \(t\) 步对应的可视空间已成功构建，则直接采用其星凸多面体的几何中心；否则退化为引导路径在该时刻的位置。记
\begin{equation}
\label{eq:center_choice}
\mathbf{c}_t=
\begin{cases}
\mathbf{c}_t^{\mathrm{poly}}, & \text{若可视空间 } S_t \text{ 有效},\\
\mathbf{p}_{\mathrm{ref}}(t), & \text{否则},
\end{cases}
\qquad t=0,\dots,H-1,
\end{equation}
其中 \(\mathbf{c}_t^{\mathrm{poly}}\) 表示 \(S_t\) 的多面体中心，\(\mathbf{p}_{\mathrm{ref}}(t)\) 为混合 A* 引导路径在时刻 \(t\) 的位置。对平面分量 \(\mathbf{c}_t^{xy}\) 采用一阶差分给出中心速度初值
\begin{equation}
\label{eq:center_vel}
\mathbf{v}^{\,c}_t=
\begin{cases}
\displaystyle \frac{\mathbf{c}_{t+1}^{xy}-\mathbf{c}_{t}^{xy}}{\Delta t}, & t=0,\dots,H-2,\\[8pt]
\displaystyle \frac{\mathbf{c}_{t}^{xy}-\mathbf{c}_{t-1}^{xy}}{\Delta t}, & t=H-1,
\end{cases}
\end{equation}
其中上标 \(xy\) 表示取位置向量的平面前两维。

为在中心轨迹周围构造合理的初始编队，对四辆 UGV 采用固定方向的交替模板并叠加随时间渐进的偏移。第 \(i\) 辆 UGV 的模板方向取自四个象限向量
\begin{equation}
\label{eq:fixed_dir}
\mathbf{d}_i\in\big\{(-1,1),\,(-1,-1),\,(1,-1),\,(1,1)\big\},
\end{equation}
并将其单位化。第 \(t\) 步的平面偏移量定义为
\begin{equation}
\label{eq:offset_series}
\boldsymbol{\delta}_{i,t}=
\begin{cases}
\mathbf{0}, & t=0,\\
r_0\,\mathrm{ramp}(t)\,\mathbf{d}_i, & t\ge 1,
\end{cases}
\qquad 
\mathrm{ramp}(t)=\min\!\left(1,\ \frac{t}{3}\right),
\end{equation}
其中 \(r_0>0\) 为期望编队半径。\(\mathrm{ramp}(t)\) 在前三个时间步内由 \(0\) 线性增长到 \(1\)，使编队在预测域前段平滑展开，避免从起点瞬间产生过大的横向分离。

在此基础上，各平台的状态初值按如下规则给出。

\textbf{UGV 初始化。} 对每辆 UGV 及每个时间步 \(t\)：
\begin{equation}
\label{eq:ugv_init}
\begin{aligned}
&\text{若 } t=0:\quad &&\mathbf{p}^{g}_{i,0}=\mathbf{p}^{g,\mathrm{cur}}_{i,xy},\quad \mathbf{v}^{g}_{i,0}=\mathbf{v}^{g,\mathrm{cur}}_{i,xy},\quad \mathbf{a}^{g}_{i,0}=\mathbf{0},\\
&\text{若 } t\ge 1:\quad &&\mathbf{p}^{g}_{i,t}=\mathbf{c}_{t}^{xy}+\boldsymbol{\delta}_{i,t},\quad \mathbf{v}^{g}_{i,t}=\mathbf{v}^{\,c}_t,\quad \mathbf{a}^{g}_{i,t}=\mathbf{0},
\end{aligned}
\end{equation}
其中 \(\mathbf{p}^{g,\mathrm{cur}}_{i,xy}\)、\(\mathbf{v}^{g,\mathrm{cur}}_{i,xy}\) 分别为第 \(i\) 辆 UGV 当前的平面位置与速度观测。

\textbf{UAV 初始化。} 对每架 UAV 及每个时间步 \(t\)：
\begin{equation}
\label{eq:uav_init}
\begin{aligned}
&\text{若 } t=0:\quad &&\mathbf{p}^{a}_{j,0}=\mathbf{p}^{a,\mathrm{cur}}_{j},\quad \mathbf{v}^{a}_{j,0}=\mathbf{v}^{a,\mathrm{cur}}_{j},\quad \mathbf{a}^{a}_{j,0}=\mathbf{0},\\
&\text{若 } t\ge 1:\quad &&\mathbf{p}^{a}_{j,t}=\mathbf{p}_{\mathrm{ref}}(t),\quad \mathbf{v}^{a}_{j,t}=\dot{\mathbf{p}}_{\mathrm{ref}}(t),\quad \mathbf{a}^{a}_{j,t}=\mathbf{0},
\end{aligned}
\end{equation}
即在首步严格对齐当前状态估计，从第二步起将无人机的位置和速度初值直接投影到混合 A* 引导路径及其切向速度上，并将加速度初始化为零，以获得平滑且数值稳定的初始轨迹。

其中\eqref{eq:center_choice}–\eqref{eq:center_vel} 使全体平台的初值与可视空间中心及全局引导趋势一致，减少优化过程中对大尺度偏移的修正；\eqref{eq:offset_series} 与 \eqref{eq:ugv_init} 通过逐步展开的小幅偏移打破多车同位姿退化，同时避免在预测初期产生过激的编队扩张，有利于可视性与编队相关代价项平稳起作用；\eqref{eq:uav_init} 则保证滚动时域在 \(t=0\) 与实时观测严格对齐，从而提高解的物理一致性。

当上一规划周期已经得到可行解时，可在上述初始化之上进一步采用热启动策略：将上一周期预测解沿时间轴整体前移一帧，对新出现的时间节点在相邻两帧之间做线性插值，得到初始序列 \(\{\mathbf{x}^{0}_{t},\mathbf{u}^{0}_{t}\}\)，并对超出状态或输入界限的分量投影回各自的约束区间。该热启动方式在保留上一周期轨迹时空结构的同时，与新的引导路径和可视空间自然对齐，可显著减少迭代次数并提高全局优化的实时性。

\subsection{无人机引导路点约束构建}
本小节以上一章中由混合 A* 搜索得到的离散参考序列为目标，在每个预测步约束无人机逐步逼近对应的参考状态。设预测域步数为 \(H\)，第 \(t\) 步第 \(j\) 架无人机的参考位置与速度分别为 \(\mathbf{p}^{a,\mathrm{ref}}_{j,t}\in\mathbb{R}^3\) 与 \(\mathbf{v}^{a,\mathrm{ref}}_{j,t}\in\mathbb{R}^3\)，对应的优化变量为 \(\mathbf{p}^a_{j,t}\) 与 \(\mathbf{v}^a_{j,t}\)。定义位置与速度的跟踪误差为
\begin{equation}
\label{eq:track_err}
\mathbf{e}^{p}_{j,t}=\mathbf{p}^a_{j,t}-\mathbf{p}^{a,\mathrm{ref}}_{j,t},\qquad
\mathbf{e}^{v}_{j,t}=\mathbf{v}^a_{j,t}-\mathbf{v}^{a,\mathrm{ref}}_{j,t},\qquad
t=0,\dots,H-1,\ j=1,\dots,N_a.
\end{equation}
由于 \(t=0\) 时刻的状态由当前观测固定，不作为决策变量参与优化，因此从 \(t=1\) 步开始施加引导路点跟踪代价。设统一权重为 \(\lambda_{\mathrm{wp}}>0\)，在末时刻 \(t=H-1\) 通过倍数 \(\kappa_{\mathrm{tail}}\ge 1\) 强化到达要求，则整个预测域上的引导路点跟踪代价写为
\begin{equation}
\label{eq:J_wp}
J_{\mathrm{wp}}
=\lambda_{\mathrm{wp}}\!\left[
\sum_{t=1}^{H-2}\sum_{j=1}^{N_a}\Big(\|\mathbf{e}^{p}_{j,t}\|_2^2+\|\mathbf{e}^{v}_{j,t}\|_2^2\Big)
+\ \kappa_{\mathrm{tail}}\sum_{j=1}^{N_a}\Big(\|\mathbf{e}^{p}_{j,H-1}\|_2^2+\|\mathbf{e}^{v}_{j,H-1}\|_2^2\Big)
\right].
\end{equation}
可以看出，\(\lambda_{\mathrm{wp}}\) 控制整体的引导强度，\(\kappa_{\mathrm{tail}}\) 用于在预测域末端进一步压缩末状态与参考状态之间的偏差，从而在不施加硬终端约束的前提下保证“按时到达”。

式~\eqref{eq:J_wp} 对决策变量的梯度具有简单的闭式形式，便于在求解器中使用。对中间时刻 \(t=1,\dots,H-2\) 有
\begin{equation}
\label{eq:grad_wp_mid}
\frac{\partial J_{\mathrm{wp}}}{\partial \mathbf{p}^a_{j,t}}
=2\,\lambda_{\mathrm{wp}}\big(\mathbf{p}^a_{j,t}-\mathbf{p}^{a,\mathrm{ref}}_{j,t}\big),\qquad
\frac{\partial J_{\mathrm{wp}}}{\partial \mathbf{v}^a_{j,t}}
=2\,\lambda_{\mathrm{wp}}\big(\mathbf{v}^a_{j,t}-\mathbf{v}^{a,\mathrm{ref}}_{j,t}\big),
\end{equation}
对末时刻 \(t=H-1\) 则为
\begin{equation}
\label{eq:grad_wp_tail}
\frac{\partial J_{\mathrm{wp}}}{\partial \mathbf{p}^a_{j,H-1}}
=2\,\lambda_{\mathrm{wp}}\,\kappa_{\mathrm{tail}}\big(\mathbf{p}^a_{j,H-1}-\mathbf{p}^{a,\mathrm{ref}}_{j,H-1}\big),\qquad
\frac{\partial J_{\mathrm{wp}}}{\partial \mathbf{v}^a_{j,H-1}}
=2\,\lambda_{\mathrm{wp}}\,\kappa_{\mathrm{tail}}\big(\mathbf{v}^a_{j,H-1}-\mathbf{v}^{a,\mathrm{ref}}_{j,H-1}\big).
\end{equation}

上述设计在整个预测域内对每一时刻的引导路点跟踪给予相同量级的权重，从而避免轨迹在时间窗口内部出现明显漂移；仅在末端通过 \(\kappa_{\mathrm{tail}}\) 进行额外放大，使无人机在预测域尾部更紧密地贴合混合 A* 给出的参考路径。结合滚动时域执行策略，当前窗口末端的强化约束同时也为下一窗口提供了良好的初始状态，从而改善整体轨迹的一致性与求解器的收敛稳定性。

\subsection{序列可视空间约束构建}

本小节利用之前构建的序列可视空间，将时序化的可视性信息显式写入全局协同优化。在每一预测时刻 \(t_k\)，可视空间 \(S(t_k)\) 刻画了以该时刻 UAV 参考位姿为中心、且对其具有良好视距条件的一片星凸区域。若某辆 UGV 的位置落在该区域外侧并留有一定裕度，则认为可视性良好；反之，当接近或穿越可视边界时，应在代价函数中给予惩罚以推动系统恢复可见状态。

对第 \(k\) 个时刻的可视空间 \(S(t_k)\)，其翻转空间中的凸多面体采用与\autoref{eq:polytope_H} 相同的 \(H\)-表示形式。记该多面体由 \(K\) 个面构成，第 \(i\) 个面的外法向为 \(n^{(k)}_i\in\mathbb{R}^3\)，面上一点为 \(a^{(k)}_i\in\mathbb{R}^3\)，则对任意翻转点 \(\hat{x}\) 有
\begin{equation}
\label{eq:sscp_face_dist}
d^{(k)}_i(\hat{x})=n^{(k)\top}_i\big(\hat{x}-a^{(k)}_i\big),\qquad i=1,\ldots,K.
\end{equation}
其对应的最大有符号距离为
\begin{equation}
\label{eq:sscp_dmax}
d_{\max}^{(k)}(\hat{x})=\max_{i=1,\ldots,K}d^{(k)}_i(\hat{x}).
\end{equation}
实际实现时，为获得平滑且可导的代价，本章采用\autoref{eq:lse_def} 所给出的形式对 \(\max\) 运算进行光滑近似。为书写简洁，以下仍以 \(d_{\max}^{(k)}(\hat{x})\) 表示该“最大外侧距离”的光滑近似值。

设第 \(i\) 辆 UGV 在时刻 \(t_k\) 的位置为 \(\mathbf{p}^{g}_{i,k}\in\mathbb{R}^3\)，其翻转点为
\begin{equation}
\hat{\mathbf{p}}^{\,g}_{i,k}=F\big(\mathbf{p}^{g}_{i,k}\big),
\end{equation}
其中 \(F(\cdot)\) 为\autoref{eq:ball_flip} 所定义的球翻转映射。给定可视安全裕度 \(d_{\min}>0\)，定义单点可视性代价为
\begin{equation}
\label{eq:sscp_point_cost}
V\!\big(\hat{\mathbf{p}}^{\,g}_{i,k}\mid S(t_k)\big)
=\Big[d_{\min}-d_{\max}^{(k)}\!\big(\hat{\mathbf{p}}^{\,g}_{i,k}\big)\Big]_+^{\,3},\qquad [z]_+=\max(z,0).
\end{equation}
当 \(\mathbf{p}^{g}_{i,k}\) 落在可视空间外部且距离边界不小于 \(d_{\min}\) 时，\(d_{\max}^{(k)}\) 足够大，有 \(d_{\min}-d_{\max}^{(k)}\le 0\)，此时代价为零，不对该时刻的可视性施加额外约束；当 UGV 靠近可视边界甚至穿越边界时，\(d_{\max}^{(k)}\) 减小，使得括号内取正值，代价按三次多项式快速增长，从而在优化中产生明显的“推回”趋势。

在整个预测域上，对所有 UGV 与全部时刻的单点代价进行加权累加，得到序列可视空间上的可视性代价
\begin{equation}
\label{eq:sscp_Jvis}
J_{\mathrm{vis}}
=\lambda_{\mathrm{vis}}\sum_{k=1}^{H-1}\sum_{i=1}^{N_g}
V\!\big(\hat{\mathbf{p}}^{\,g}_{i,k}\mid S(t_k)\big),
\end{equation}
其中 \(N_g\) 为 UGV 数量，\(\lambda_{\mathrm{vis}}>0\) 为可视性代价的权重系数。

\subsection{空地队形约束构建}

为提高基于 UWB 的相对定位稳定性与精度，本节在上一章空地队形约束的基础上进行优化。在每个离散时刻 \(k\)，以 UGV 的平面位置\(\{\mathbf{p}^{g}_{i,k}\in\mathbb{R}^2\}_{i=1}^{N_g}\) 为优化变量，从边长均衡、整体尺度、径向均衡以及空地距离四个方面构造复合队形损失，并通过一致的顶点排序避免随时间的标号混乱，如\autoref{fig:formation_example}所示。

\begin{figure}[htbp]
    \centering
    \includegraphics[width=0.75\linewidth]{chapter_4/formation.png}
    \caption{队形约束构建示意图}
    \label{fig:formation_example}
\end{figure}

首先在时刻 \(k\) 计算 UGV 多边形的质心
\begin{equation}
\label{eq:formation_centroid}
\mathbf{c}_k=\frac{1}{N_g}\sum_{i=1}^{N_g}\mathbf{p}^{g}_{i,k},
\end{equation}
并以 \(\mathbf{c}_k\) 为极点，计算各 UGV 的极角
\begin{equation}
\label{eq:formation_angle}
\varphi_i=\mathrm{atan2}\!\big(p^{g}_{i,k,y}-c_{k,y},\ p^{g}_{i,k,x}-c_{k,x}\big).
\end{equation}
按 \(\varphi_i\) 升序得到顶点序列 \((v_0,v_1,\dots,v_{N_g-1})\)，据此定义有向边 \((v_i,v_{i+1\bmod N_g})\)。该基于质心的极角排序在 UGV 发生交叉换位时仍能保持边的拓扑顺序一致，避免在后续边长与径向代价中出现“错边配点”的问题。

在此基础上引入四类子代价。第一类为边长均衡项，用于促使多边形边长接近一致，从而抑制过长或过短的边带来的几何偏置。记第 \(i\) 条边的长度
\begin{equation}
\label{eq:edge_length}
d_i=\big\|\mathbf{p}^{g}_{v_{i+1},k}-\mathbf{p}^{g}_{v_i,k}\big\|_2,\qquad 
\bar d=\frac{1}{N_g}\sum_{i=0}^{N_g-1}d_i,
\end{equation}
则边长均衡代价定义为
\begin{equation}
\label{eq:J_edge}
J_{\mathrm{edge}}(k)=\lambda_{\mathrm{edge}}\sum_{i=0}^{N_g-1}\big(d_i-\bar d\big)^2,
\end{equation}
其中 \(\lambda_{\mathrm{edge}}>0\) 控制对形状规则性的偏好。该项在所有边长一致时取零，当某些边明显偏长或偏短时会产生二次惩罚，将多边形整体推向更接近矩形或正多边形的形状。

第二类为整体尺度约束，用于限制平均边长处于期望区间，避免车距过近导致多径和安全风险，也避免过远导致测距退化。给定允许距离带 \([d_{\mathrm{low}},d_{\mathrm{high}}]\)，定义
\begin{equation}
\label{eq:J_mean}
J_{\mathrm{mean}}(k)=\lambda_{\mathrm{mean}}\Big(\,[d_{\mathrm{low}}-\bar d]_+^{\,2}+[\,\bar d-d_{\mathrm{high}}]_+^{\,2}\Big),
\qquad [z]_+=\max(z,0),
\end{equation}
其中 \(\lambda_{\mathrm{mean}}>0\) 为尺度权重。当 \(\bar d\) 处于可行带内部时该项为零；一旦整体编队被“压缩”或“拉伸”超出期望范围，代价按与带边界差值的平方增长，从而在保持一定机动自由度的同时限制整体尺度漂移。对各边长度 \(d_i\) 的偏导数可按
\begin{equation}
\label{eq:J_mean_dk}
\frac{\partial J_{\mathrm{mean}}}{\partial d_i}=\frac{\lambda_{\mathrm{mean}}}{N_g}
\begin{cases}
-2\,(d_{\mathrm{low}}-\bar d),& \bar d<d_{\mathrm{low}},\\
\ \ 2\,(\bar d-d_{\mathrm{high}}),& \bar d>d_{\mathrm{high}},\\
\ \ 0,& \text{otherwise},
\end{cases}
\end{equation}
再结合
\begin{equation}
\label{eq:edge_dpos}
\frac{\partial d_i}{\partial \mathbf{p}^{g}_{v_i,k}}=-\,\frac{\mathbf{p}^{g}_{v_{i+1},k}-\mathbf{p}^{g}_{v_i,k}}{d_i},\qquad
\frac{\partial d_i}{\partial \mathbf{p}^{g}_{v_{i+1},k}}=\ \ \frac{\mathbf{p}^{g}_{v_{i+1},k}-\mathbf{p}^{g}_{v_i,k}}{d_i},
\end{equation}
即可通过链式法则得到对各顶点位置的梯度。

第三类为径向均衡项，用于抑制“边长合适但个别点半径过大或过小”的退化配置，使各 UGV 相对质心的分布更为均匀。令
\begin{equation}
\label{eq:radial_def}
r_i=\big\|\mathbf{p}^{g}_{v_i,k}-\mathbf{c}_k\big\|_2,\qquad 
\bar r=\frac{1}{N_g}\sum_{i=0}^{N_g-1}r_i,
\end{equation}
则径向方差代价为
\begin{equation}
\label{eq:J_radial}
J_{\mathrm{rad}}(k)=\lambda_{\mathrm{rad}}\sum_{i=0}^{N_g-1}\big(r_i-\bar r\big)^2,
\end{equation}
其中 \(\lambda_{\mathrm{rad}}>0\) 控制对“环状展开”程度的约束。该项在所有 UGV 到质心距离接近时趋于零，当某些车辆距离质心明显偏离整体平均时会产生惩罚，有利于避免极度扁长或一侧塌缩的几何形状。其梯度可通过 \(r_i\) 对 \(\mathbf{p}^{g}_{v_i,k}\) 的导数按链式法则直接写出，这里不再展开。

第四类为空地距离项，用于约束 UAV 与各 UGV 之间的空间距离不过度增大，以保障 UWB 测距链路与空地协同观测质量。设 UAV 数量为 \(N_a\)，在时刻 \(k\) 的三维位置为 \(\{\mathbf{p}^{a}_{j,k}\in\mathbb{R}^3\}_{j=1}^{N_a}\)。将 UGV 的平面坐标嵌入三维并与地面高度对齐，记
\begin{equation}
\label{eq:bdry_dist}
d_{ij}=\big\|\mathbf{p}^{a}_{j,k}-[\mathbf{p}^{g}_{i,k};0]\big\|_2,
\end{equation}
其中 \([\cdot;0]\) 表示在 \(z\) 轴上补零。给定最大期望量程 \(d_{\mathrm{range}}>0\)，定义
\begin{equation}
\label{eq:J_bdry}
J_{\mathrm{bd}}(k)=\lambda_{\mathrm{bd}}\sum_{i=1}^{N_g}\sum_{j=1}^{N_a}\,[\,d_{ij}-d_{\mathrm{range}}\,]_+^{\,2},
\end{equation}
其中 \(\lambda_{\mathrm{bd}}>0\) 为空地距离项的权重。当 UAV 长时间远离某些 UGV 超出期望量程时，该项会对 \(\mathbf{p}^{a}_{j,k}\) 和 \(\mathbf{p}^{g}_{i,k}\) 产生将其相互拉近的梯度，从而鼓励编队保持一定的空地耦合强度。对应的偏导数可按
\begin{equation}
\label{eq:J_bdry_grad}
\frac{\partial J_{\mathrm{bd}}}{\partial \mathbf{p}^{a}_{j,k}}=
2\,\lambda_{\mathrm{bd}}\,(d_{ij}-d_{\mathrm{range}})\,\frac{\mathbf{p}^{a}_{j,k}-[\mathbf{p}^{g}_{i,k};0]}{d_{ij}},\qquad
\frac{\partial J_{\mathrm{bd}}}{\partial \mathbf{p}^{g}_{i,k}}= -\Big(\frac{\partial J_{\mathrm{bd}}}{\partial \mathbf{p}^{a}_{j,k}}\Big)_{xy},
\end{equation}
对所有 \((i,j)\) 满足 \(d_{ij}>d_{\mathrm{range}}\) 的对组合累加即可得到完整梯度。

将上述四类子代价在时域上求和，得到空地队形总损失
\begin{equation}
\label{eq:J_form_total}
J_{\mathrm{form}}=\sum_{k=0}^{H-1}\Big(J_{\mathrm{edge}}(k)+J_{\mathrm{mean}}(k)+J_{\mathrm{rad}}(k)+J_{\mathrm{bd}}(k)\Big),
\end{equation}
其中 \(\lambda_{\mathrm{edge}},\lambda_{\mathrm{mean}},\lambda_{\mathrm{rad}},\lambda_{\mathrm{bd}}>0\) 为各子项权重。由于每一项均由欧氏距离与二次函数等基础运算组成，其关于 \(\mathbf{p}^{g}_{i,k}\) 与 \(\mathbf{p}^{a}_{j,k}\) 的梯度都可以按\autoref{eq:J_edge}–\autoref{eq:J_bdry_grad} 中的解析表达式直接求得，并线性叠加构成 \(\partial J_{\mathrm{form}}/\partial (\cdot)\)。通过边长均衡、整体尺度、径向均衡与空地距离四个层面的约束，UGV 编队在全局优化中能够保持近似矩形且尺度合理的拓扑结构，同时与 UAV 维持合适的空间距离，从而为 UWB 相对定位与空地协同感知提供更加稳定、可观的几何布局。

\subsection{优化问题建模与求解}\label{sec:opt_que}

本节在预测域长度 \(H\) 和步长 \(\Delta t\) 下，对 UAV 与 UGV 统一建模并构建面向全局协同规划的非线性优化问题。每个机器人 \(i\) 在离散时刻 \(k\) 的决策变量为位置、速度与加速度
\begin{equation}
\label{eq:state_stack_final}
\mathbf{x}_{i,k}=\begin{bmatrix}\mathbf{p}_{i,k}\\[2pt] \mathbf{v}_{i,k}\\[2pt] \mathbf{a}_{i,k}\end{bmatrix},\qquad
\mathbf{p}_{i,k},\mathbf{v}_{i,k},\mathbf{a}_{i,k}\in\mathbb{R}^{d_i},\ d_i\in\{2,3\},
\end{equation}
其中 \(d_i=3\) 对应 UAV，\(d_i=2\) 对应在平面上运动的 UGV。加速度 \(\mathbf{a}_{i,k}\) 在全局层面同时扮演控制输入的角色，后续局部轨迹生成模块在此基础上进行连续时间参数化。为保持与前文建模一致，状态在离散时间上满足二阶积分形式
\begin{equation}
\label{eq:disc_dyn_final}
\mathbf{p}_{i,k+1}=\mathbf{p}_{i,k}+\Delta t\,\mathbf{v}_{i,k}+\tfrac{1}{2}\Delta t^2\,\mathbf{a}_{i,k},\qquad
\mathbf{v}_{i,k+1}=\mathbf{v}_{i,k}+\Delta t\,\mathbf{a}_{i,k},
\end{equation}
初始条件由当前状态估计给出
\begin{equation}
\label{eq:init_cond_final}
\mathbf{p}_{i,0}=\mathbf{p}^{\mathrm{init}}_{i},\qquad \mathbf{v}_{i,0}=\mathbf{v}^{\mathrm{init}}_{i},
\end{equation}
并在实现中通过直接赋值固定，不作为优化变量参与求解。上述抽象在保持 UAV 微分平坦可跟踪性和 UGV 平面全向可执行性的前提下，使不同平台共享统一的变量结构，便于在同一优化框架内叠加引导、可视、队形与能量等多种代价。

为将动力学约束和速度／加速度上下限约束纳入无约束求解框架，采用平方惩罚的方式将其转化为损失项并加入总目标函数。这样可以避免显式构造大规模等式与不等式约束集合，便于配合拟牛顿方向进行迭代求解，同时在有良好热启动的前提下将搜索过程限制在近似动态可行的邻域。

首先考虑动力学一致性。对每个机器人 \(i\) 和每个 \(k\ge 1\)，定义相邻时刻的动力学“缺陷”
\begin{equation}
\label{eq:dyn_defects_final}
\mathbf{c}^{(1)}_{i,k}=\mathbf{p}_{i,k}-\mathbf{p}_{i,k-1}-\Delta t\,\mathbf{v}_{i,k-1},\qquad
\mathbf{c}^{(2)}_{i,k}=\mathbf{v}_{i,k}-\mathbf{v}_{i,k-1}-\Delta t\,\mathbf{a}_{i,k-1},
\end{equation}
并采用二次惩罚
\begin{equation}
\label{eq:J_dyn_final}
J_{\mathrm{dyn}}=\lambda_{\mathrm{dyn}}\sum_{k=1}^{H-1}\sum_{i}\Big(\|\mathbf{c}^{(1)}_{i,k}\|_2^2+\|\mathbf{c}^{(2)}_{i,k}\|_2^2\Big),
\end{equation}
其中 \(\lambda_{\mathrm{dyn}}>0\) 控制动力学一致性的严格程度。当轨迹严格满足\autoref{eq:disc_dyn_final} 时，\(\mathbf{c}^{(1)}_{i,k}\) 与 \(\mathbf{c}^{(2)}_{i,k}\) 等于零，\(J_{\mathrm{dyn}}\) 消失；在数值求解中，通过逐步减小缺陷可将解推向动态可行区域。

然后处理速度与加速度的上下限。对每个机器人给定最大速度 \(v_{\max,i}>0\) 和最大加速度 \(a_{\max,i}>0\)，对任意 \((i,k)\) 定义分段平方惩罚
\begin{equation}
\label{eq:rho_funcs}
\rho_v(n;v_{\max,i})=\begin{cases}
\big(\tfrac{n}{v_{\max,i}}-1\big)^2,& n>v_{\max,i},\\[3pt]
0,& n\le v_{\max,i},
\end{cases}
\qquad
\rho_a(n;a_{\max,i})=\begin{cases}
\big(\tfrac{n}{a_{\max,i}}-1\big)^2,& n>a_{\max,i},\\[3pt]
0,& n\le a_{\max,i},
\end{cases}
\end{equation}
由此得到上下限损失
\begin{equation}
\label{eq:J_limit_final}
J_{\mathrm{lim}}=\lambda_{\mathrm{lim}}\sum_{k=0}^{H-1}\sum_{i}\Big[
\rho_v\big(\|\mathbf{v}_{i,k}\|_2;v_{\max,i}\big)+
\rho_a\big(\|\mathbf{a}_{i,k}\|_2;a_{\max,i}\big)\Big],
\end{equation}
其中 \(\lambda_{\mathrm{lim}}>0\) 为对应权重。当速度与加速度范数处于物理上限之内时该项为零，一旦超过上限则按相对超限比例的平方进行惩罚，从而在保持一定机动裕度的同时限制极端加减速行为。

控制能量与轨迹平滑性通过加速度能量和相邻步加速度差的平方进行正则化，这与前一章控制正则项的设计保持一致。定义
\begin{equation}
\label{eq:J_ctrl_final}
J_{\mathrm{ctrl}}=\lambda_{\mathrm{ctrl}}\sum_{k=0}^{H-1}\sum_{i}\|\mathbf{a}_{i,k}\|_2^{2}
+\lambda_{\Delta u}\sum_{k=0}^{H-2}\sum_{i}\|\mathbf{a}_{i,k+1}-\mathbf{a}_{i,k}\|_2^{2},
\end{equation}
其中第一项抑制整体控制能量，第二项抑制加速度在时间上的剧烈变化，\(\lambda_{\Delta u}>0\) 为平滑权重。对应的全局权重 \(\lambda_{\mathrm{ctrl}}\) 在实验部分与前一章保持一致，用于调节控制正则在总目标中的相对重要性。

在此基础上，前文已经分别给出了无人机引导路点代价 \(J_{\mathrm{wp}}\)、序列可视空间代价 \(J_{\mathrm{vis}}\) 和无人车队形代价 \(J_{\mathrm{form}}\)。综合上述各项，本章全局协同规划阶段的目标函数写为
\begin{equation}
\label{eq:J_total_final}
\min_{\{\mathbf{p}_{i,k},\mathbf{v}_{i,k},\mathbf{a}_{i,k}\}}~
J=J_{\mathrm{wp}}+J_{\mathrm{vis}}+J_{\mathrm{form}}
+J_{\mathrm{ctrl}}+J_{\mathrm{dyn}}+J_{\mathrm{lim}},
\end{equation}
其中 \(J_{\mathrm{wp}}\) 推动 UAV 渐近跟踪混合 A* 引导路点，\(J_{\mathrm{vis}}\) 通过 SSCP 在时间序列上保持空地可视性，\(J_{\mathrm{form}}\) 维持空地队形的形状与尺度，\(J_{\mathrm{ctrl}}\) 抑制能量与抖动，\(J_{\mathrm{dyn}}\) 和 \(J_{\mathrm{lim}}\) 将动态一致性与物理上下限以惩罚形式嵌入统一优化框架，使总代价在符号和含义上与前一章保持一致。

求解方面，考虑到预测域内状态与控制变量随机器人数量线性增长，而动力学一致性在时间上耦合、可视性和队形代价在同一时刻跨机器人耦合，直接构造大规模稀疏 KKT 系统进行二阶内点法求解会带来较高的实现与计算开销。为此，本章在实现中采用开源的 L\text{-}BFGS 方法\cite{LBFGS-Lite}。该方法通过保留有限条历史 \((s_\ell,y_\ell)\) 向量对构造近似逆 Hessian，在每一步利用两层循环公式计算搜索方向，并结合满足强 Wolfe 条件的线搜索确定步长，在平滑目标上具有较好的收敛效率和较低的内存占用。依托前文给出的各损失项解析梯度，能够在每个滚动周期内高效构建目标与梯度，并完成若干步迭代收敛到局部极小。

\section{局部跟踪避障轨迹生成} \label{sec:local-tracking}
本节在全局离散规划结果的基础上开展局部连续轨迹生成与优化。首先，将带时间戳的全局离散路径点按照其时间顺序初始化为本地多项式轨迹，采用 MINCO 形式\cite{wang2022geometrically}以同时满足位姿与导数的连续性，并将每一个全局路径点的时间标签保留用于后续时间对齐。随后，在该本地轨迹上引入避障与跟踪等目标进行连续优化：避障项采用基于碰撞段、锚点与排斥方向估计的力场方法，对处于障碍物内部或接近障碍边界的轨迹片段，沿由锚点指示的无碰撞方向对轨迹关键变量施加排斥代价与梯度，从而在不依赖全局 ESDF 的条件下高效剔除碰撞\cite{zhou2021egoplanner} 。在此基础上，增加路径点同步约束：对每一个全局时间戳 \(t_k\)，在待优化的本地轨迹上取同一时刻的点 \(\mathbf{r}(t_k)\)，并将 \(\|\mathbf{r}(t_k)-\mathbf{w}_k\|_2^2\) 的平方偏差作为跟踪损失加入整体目标，使局部解在时空上严格贴合全局离散解。最后，叠加速度与加速度的上下限惩罚以及能量与平滑项，形成面向实时执行的局部轨迹优化问题。

\subsection{局部轨迹避障方法}

本小节在本地 MINCO 轨迹上直接构建可微的避障代价。\autoref{fig:local_avoid} 给出了基于 \(\{\mathbf{p},\mathbf{v}\}\) 对的几何直观：首先在障碍物内部或邻近处选取碰撞采样点及其切向速度，由此在障碍物表面上寻找局部无碰撞路径，再根据得到的推出方向对轨迹施加排斥力场，使最终轨迹被“推”出障碍物外侧。

\begin{figure}[htbp]
    \centering
    \includegraphics[width=0.7\linewidth]{chapter_4/local_avoidance_method.jpg}
    \caption{基于 \(\{\mathbf{p},\mathbf{v}\}\) 对的局部避障方法示意图\cite{egoswarm2021}}
    \label{fig:local_avoid}
\end{figure}

局部轨迹采用 MINCO 形式表示，由分段时长
\(\mathbf{T}=[T_1,\dots,T_M]^\top\) 与内部路径点参数
\(\mathbf{q}=[\mathbf{q}_1,\dots,\mathbf{q}_{M-1}]\) 唯一定义，记
\begin{equation}
\label{eq:minco_local_rep_all}
\mathbf{p}(t)=\mathcal{M}\!\big(\mathbf{q},\mathbf{T};t\big),\qquad 
\mathbf{v}(t)=\dot{\mathbf{p}}(t),\qquad 
t\in[t_0,t_M],\ \ t_M-t_0=\textstyle\sum_{m=1}^{M}T_m .
\end{equation}
为将连续时间代价离散化，在区间 \([t_0,t_M]\) 上等间隔采样
\(t_i=t_0+i\,\Delta t_s\)（\(i=0,\dots,K\)），并令时间权重 \(w_i=\Delta t_s\)。

在每个采样时刻，利用轨迹的切向方向与障碍边界构造一组局部锚点与排斥方向。记
\(\mathbf{Q}_i=\mathbf{p}(t_i)\) 为第 \(i\) 个采样点，
\(\mathbf{R}_i=\mathbf{v}(t_i)/\|\mathbf{v}(t_i)\|_2\) 为单位切向，在
\(\mathbf{Q}_i\) 处建立切片平面
\begin{equation}
\label{eq:slice_plane_all}
\Psi_i=\Big\{\mathbf{x}\in\mathbb{R}^3\ \big|\ (\mathbf{x}-\mathbf{Q}_i)^\top \mathbf{R}_i=0\Big\}.
\end{equation}
将 \(\Psi_i\) 与障碍边界 \(\Gamma\) 相交得到若干线段，将 \(\mathbf{Q}_i\) 正交投影到每条线段并取最近点作为候选锚点 \(\mathbf{p}_{ij}\)。对每个候选构造归一化方向
\begin{equation}
\label{eq:pair_dir_all}
\tilde{\mathbf{v}}_{ij}=\frac{\mathbf{p}_{ij}-\mathbf{Q}_i}{\|\mathbf{p}_{ij}-\mathbf{Q}_i\|_2},
\end{equation}
再通过一次占据检查将方向修正为指向自由空间的一侧：若
\(\mathbf{p}_{ij}+\varepsilon\,\tilde{\mathbf{v}}_{ij}\)（\(\varepsilon>0\) 很小）仍在障碍内部，则置 \(\mathbf{v}_{ij}=-\tilde{\mathbf{v}}_{ij}\)，否则置 \(\mathbf{v}_{ij}=\tilde{\mathbf{v}}_{ij}\)。这样，每一对 \(\{\mathbf{p}_{ij},\mathbf{v}_{ij}\}\) 都定义了一张通过 \(\mathbf{p}_{ij}\)、法向为 \(\mathbf{v}_{ij}\) 的局部“零距离平面”。

相对于第 \(j\) 个候选对，采样点 \(\mathbf{Q}_i\) 的有符号距离写为
\begin{equation}
\label{eq:pair_dist_all}
d_{ij}=\big(\mathbf{Q}_i-\mathbf{p}_{ij}\big)^\top \mathbf{v}_{ij}.
\end{equation}
从所有候选中选取距离最小的一对
\begin{equation}
\label{eq:active_pair_all}
d_i=\min_{j} d_{ij},\qquad 
\mathbf{n}_i=\mathbf{v}_{i j^\star},\quad j^\star=\arg\min_{j} d_{ij},
\end{equation}
即可得到与当前环境最一致的推出方向 \(\mathbf{n}_i\) 及最近距离 \(d_i\)。\autoref{fig:local_avoid} 中的 \(\Gamma\)、\(\mathbf{p}\) 与 \(\mathbf{v}\) 分别对应障碍物边界、锚点与排斥方向的几何示意。

给定几何安全裕度 \(C_o>0\)，在 \(d_i<C_o\) 时认定轨迹与障碍过近，并施加三次距离惩罚，沿时间累加得到避障代价
\begin{equation}
\label{eq:Jo_all}
J_{o}=\lambda_{o}\sum_{i=0}^{K} w_i\,\big(C_o-d_i\big)_{+}^{3},\qquad [x]_+=\max(x,0),\ \ \lambda_o>0,
\end{equation}
其中 \(d_i\) 关于 \(\mathbf{Q}_i=\mathbf{p}(t_i)\) 线性，因而 \(J_o\) 对轨迹是可微的。当 \(d_i<C_o\) 时，其对采样位置的梯度为
\begin{equation}
\label{eq:Jo_grad_p_all}
\frac{\partial J_o}{\partial \mathbf{p}(t_i)}=-3\,\lambda_{o}\,w_i\,\big(C_o-d_i\big)^{2}\,\mathbf{n}_i,
\end{equation}
否则梯度为零。利用 MINCO 的解析雅可比，可将梯度回传到参数空间
\begin{equation}
\label{eq:chain_rule_all}
\frac{\partial J_o}{\partial \mathbf{q}}
=\sum_{i=0}^{K}\!\Big(\frac{\partial \mathbf{p}(t_i)}{\partial \mathbf{q}}\Big)^{\!\top}\frac{\partial J_o}{\partial \mathbf{p}(t_i)},\qquad
\frac{\partial J_o}{\partial \mathbf{T}}
=\sum_{i=0}^{K}\!\Big(\frac{\partial \mathbf{p}(t_i)}{\partial \mathbf{T}}\Big)^{\!\top}\frac{\partial J_o}{\partial \mathbf{p}(t_i)}.
\end{equation}

上述构造在每个采样时刻仅激活一对 \(\{\mathbf{p}_{ij},\mathbf{v}_{ij}\}\)，计算与存储开销与候选数线性相关；通过方向校正保证排斥力始终指向障碍外部；三次惩罚在边界处导数连续，有利于与 MINCO 轨迹的可导表示一同集成进局部轨迹优化问题，实现高效、光滑的避障调整。

\subsection{路径点同步跟踪约束构建}

全局前端输出带时间标记的离散参考序列
\(\mathcal{W}=\{(\mathbf{r}_s,\tau_s)\}_{s=0}^{S}\)，其中
\(\mathbf{r}_s\in\mathbb{R}^3\) 为第 \(s\) 个全局路径点，\(\tau_s\) 为其对应时间戳。
本地轨迹采用 MINCO 形式 \(\mathbf{p}(t)=\mathcal{M}(\mathbf{q},\mathbf{T};t)\)，
定义局部时间区间 \([t_0,t_M]\)，并使用投影算子
\(\Pi_{[t_0,t_M]}(\cdot)\) 将所有参考时间映射到该区间内，记
\(\tilde{\tau}_s=\Pi_{[t_0,t_M]}(\tau_s)\)。在每个映射时刻直接在
\(\mathbf{p}(t)\) 上取样，与对应的全局路径点做差，构造同步跟踪损失
\begin{equation}
\label{eq:J_sync_def}
J_{\mathrm{sync}}
=\lambda_{\mathrm{sync}}\sum_{s=0}^{S}\big\|\mathbf{p}(\tilde{\tau}_s)-\mathbf{r}_s\big\|_2^{2},
\qquad \lambda_{\mathrm{sync}}>0.
\end{equation}
该设计将时间与位置约束统一在同一代价项中，使局部 MINCO 轨迹在给定时间戳上紧贴全局离散路径，从而保持节奏一致并抑制沿时间轴的漂移。

对采样位置的梯度为
\begin{equation}
\label{eq:J_sync_grad_p}
\frac{\partial J_{\mathrm{sync}}}{\partial \mathbf{p}(\tilde{\tau}_s)}
=2\,\lambda_{\mathrm{sync}}\big(\mathbf{p}(\tilde{\tau}_s)-\mathbf{r}_s\big).
\end{equation}
利用 MINCO 的解析雅可比将其回传到参数空间。令
\begin{equation}
\label{eq:minco_jac_sync}
\mathbf{J}^{p}_{s,q}=\frac{\partial \mathbf{p}(\tilde{\tau}_s)}{\partial \mathbf{q}},\qquad
\mathbf{J}^{p}_{s,T}=\frac{\partial \mathbf{p}(\tilde{\tau}_s)}{\partial \mathbf{T}},
\end{equation}
则有
\begin{equation}
\label{eq:J_sync_grad_q}
\frac{\partial J_{\mathrm{sync}}}{\partial \mathbf{q}}
=\sum_{s=0}^{S}\big(\mathbf{J}^{p}_{s,q}\big)^{\!\top}
\frac{\partial J_{\mathrm{sync}}}{\partial \mathbf{p}(\tilde{\tau}_s)},
\qquad
\frac{\partial J_{\mathrm{sync}}}{\partial \mathbf{T}}
=\sum_{s=0}^{S}\big(\mathbf{J}^{p}_{s,T}\big)^{\!\top}
\frac{\partial J_{\mathrm{sync}}}{\partial \mathbf{p}(\tilde{\tau}_s)}.
\end{equation}
在实现中，当 \(\tilde{\tau}_s\) 落在第 \(m\) 段多项式的局部时间内时，
\(\mathbf{p}(\tilde{\tau}_s)\) 由该段基函数直接计算，
\(\mathbf{J}^{p}_{s,q}\) 与 \(\mathbf{J}^{p}_{s,T}\) 也由 MINCO 的线性映射显式给出，无需数值差分。
将 \eqref{eq:J_sync_def}–\eqref{eq:J_sync_grad_q} 与局部避障项及能量平滑项一并加入目标函数后，
优化器即可在严格时间对齐的前提下调整轨迹形状，保证局部执行对全局规划的时空一致跟踪。

\subsection{最小能量轨迹优化方法}

本地轨迹优化的目标是在满足避障代价 \(J_{o}\) 与路径点同步代价 \(J_{\mathrm{sync}}\) 的前提下，使轨迹同时具有良好的能量效率和动态可执行性。为此，在 MINCO 参数空间 \((\mathbf{q},\mathbf{T})\) 上引入平滑度、总时间和动态可行三类代价，与前述两项共同最小化，得到可直接用于控制的平滑轨迹。

令 \(\mathbf{p}(t)=\mathcal{M}(\mathbf{q},\mathbf{T};t)\) 为分段多项式轨迹，取导数阶 \(s\in\{3,4\}\)（\(s=3\) 为最小 snap），在全时域 \([t_0,t_M]\) 上惩罚高阶导数的二范数积分
\begin{equation}
\label{eq:Js}
J_s=\int_{t_0}^{t_M}\big\|\mathbf{p}^{(s)}(t)\big\|_2^{2}\,{\rm d}t.
\end{equation}
依托 MINCO 的分段多项式结构，\eqref{eq:Js} 可解析化为 \((\mathbf{q},\mathbf{T})\) 的二次型并得到闭式梯度，无需数值积分。

其次为鼓励较短的飞行时间并预留控制裕度，引入总时长惩罚
\begin{equation}
\label{eq:Jt}
J_t=\sum_{m=1}^{M}T_m=t_M-t_0 ,
\end{equation}
其梯度简单为 \(\partial J_t/\partial T_m=1\)。

再次为满足平台的速度、加速度与跃度上限，在离散采样集 \(t_i=t_0+i\Delta t_s\)（\(i=0,\dots,K\)）上对越界进行惩罚。记 \(w_i=\Delta t_s\)，
\(\mathbf{v}_i=\dot{\mathbf{p}}(t_i)\)、\(\mathbf{a}_i=\ddot{\mathbf{p}}(t_i)\)、\(\mathbf{j}_i=\dddot{\mathbf{p}}(t_i)\)，允许上限为 \(v_m,a_m,j_m>0\)，则
\begin{align}
\label{eq:Jdv}
J_{d,v}&=\sum_{i=0}^{K} w_i\,\Big[\,\|\mathbf{v}_i\|_2^{2}-v_m^{2}\,\Big]_{+}^{3},\\
\label{eq:Jda}
J_{d,a}&=\sum_{i=0}^{K} w_i\,\Big[\,\|\mathbf{a}_i\|_2^{2}-a_m^{2}\,\Big]_{+}^{3},\\
\label{eq:Jdj}
J_{d,j}&=\sum_{i=0}^{K} w_i\,\Big[\,\|\mathbf{j}_i\|_2^{2}-j_m^{2}\,\Big]_{+}^{3},\\
\label{eq:Jd_all}
J_d&=J_{d,v}+J_{d,a}+J_{d,j},
\end{align}
其中 \([x]_+=\max(x,0)\)。当任一导数量平方超过阈值时，代价以三次形式快速增长，阈值处导数连续，有利于数值稳定。记
\(s_i=\|\mathbf{y}_i\|_2^{2}-y_m^{2}\)（\(\mathbf{y}_i\) 依次代指 \(\mathbf{v}_i,\mathbf{a}_i,\mathbf{j}_i\)，\(y_m\) 代指 \(v_m,a_m,j_m\)），若 \(s_i>0\) 有
\begin{equation}
\label{eq:Jd_grad_y}
\frac{\partial}{\partial \mathbf{y}_i}\sum w_i\,[s_i]_{+}^{3}
=6\,w_i\,s_i^{2}\,\mathbf{y}_i ,
\end{equation}
再通过 MINCO 的解析雅可比
\(\partial \dot{\mathbf{p}}/\partial(\mathbf{q},\mathbf{T})\)、
\(\partial \ddot{\mathbf{p}}/\partial(\mathbf{q},\mathbf{T})\)、
\(\partial \dddot{\mathbf{p}}/\partial(\mathbf{q},\mathbf{T})\)
回传到 \((\mathbf{q},\mathbf{T})\)。

综合避障、同步、平滑、时长与动态可行性，本地轨迹的总目标函数写为
\begin{equation}
\label{eq:J_total_minco_final}
\min_{\mathbf{q},\,\mathbf{T}}\;
J_{\mathrm{total}}
=J_{o}
+J_{\mathrm{sync}}
+\lambda_{s} J_{s}
+\lambda_{t} J_{t}
+\lambda_{d} J_{d},
\end{equation}
其中 \(\lambda_{o},\lambda_{\mathrm{sync}},\lambda_{s},\lambda_{t},\lambda_{d}>0\) 为各子项权重。\(
J_o\) 的位置梯度由局部平面法向给出并按式~\eqref{eq:chain_rule_all} 回传，
\(J_{\mathrm{sync}}\) 的梯度由式~\eqref{eq:J_sync_grad_p}–\eqref{eq:J_sync_grad_q} 得到，
\(J_s\) 由多项式积分的闭式表达计算，
\(J_t\) 的梯度为常数向量，
\(J_d\) 的梯度由式~\eqref{eq:Jd_grad_y} 在三阶导数上类推得到。

求解采用 L\text{-}BFGS，在仅依赖一阶信息的前提下利用有限内存的拟牛顿近似与 Wolfe 线搜索对光滑目标实现快速收敛，计算与存储开销随采样数线性增长，且便于从上一周期解热启动。优化得到的 MINCO 轨迹按时间持续输出 \(\mathbf{p}(t),\mathbf{v}(t),\mathbf{a}(t)\) 给各机器人控制器：对无人机，通过微分平坦映射将 \(\{\mathbf{p},\dot{\mathbf{p}},\ddot{\mathbf{p}}\}\) 转为推力与角速度参考；对无人车，直接作为平面速度与加速度参考，从而获得平滑、可行且能量友好的局部跟踪轨迹。

\section{两阶段协同规划方法架构设计} \label{sec:two-stage-arch}

本节从系统角度梳理全局规划与局部规划之间的分工与数据流。整体架构如图~\ref{fig:two_stage_arch} 所示，系统在外环执行基于序列可视空间的全局协同规划，在内环执行基于 MINCO 的局部轨迹优化，两者通过带时间戳的参考序列进行耦合，实现“全局可视一致性 + 本地连续可执行”的协同。

\begin{figure}[htbp]
    \centering
    \includegraphics[width=1\linewidth]{chapter_4/two_stage_arch.png}
    \caption{基于序列可视空间的两阶段协同规划总体架构示意图}
    \label{fig:two_stage_arch}
\end{figure}

在全局阶段，首先由定位与感知模块融合 UAV 与多辆 UGV 的位姿信息，构建共享占据栅格地图。以当前 UAV 位置为起点、任务目标为终点，基于混合 A* 生成满足动力学约束的全局引导路径及其时间参数化。沿该引导路径按固定时间步长采样，并在每个采样点附近构建可视空间，进而形成时序化的序列可视空间（SSCP）。在\autoref{sec:sscp-global} 中给出的建模基础上，外环在预测域长度 \(H_g\) 内，对所有机器人在每个离散时刻的状态 \(\{\mathbf{p}_{i,k},\mathbf{v}_{i,k},\mathbf{a}_{i,k}\}\) 求解全局优化问题 \eqref{eq:J_total_final}，综合考虑无人机引导路点跟踪 \(J_{\mathrm{wp}}\)、序列可视性代价 \(J_{\mathrm{vis}}\)、地面队形代价 \(J_{\mathrm{form}}\) 以及能量、动力学与上下限惩罚等项，得到未来一段时间的离散参考轨迹：
\[
\mathcal{W}^{(r)}_k=\Big\{\big(\mathbf{p}^{\mathrm{ref}(r)}_{k,s},\ \mathbf{v}^{\mathrm{ref}(r)}_{k,s},\ \tau_{k,s}\big)\ \big|\ s=0,\dots,H_g-1\Big\},
\]
其中 \(r\) 为机器人索引，\(\tau_{k,s}=t_k+s\Delta t_g\) 为显式时间戳。该集合是外环对内环的输出接口。

在局部阶段，每个机器人独立运行局部轨迹优化器。对于第 \(r\) 个机器人，收到其对应的参考序列 \(\mathcal{W}^{(r)}_k\) 后，在时间区间 \([t_k,\ t_k+T_\ell]\) 上构建 MINCO 轨迹
\(\mathbf{p}^{(r)}(t)=\mathcal{M}(\mathbf{q}^{(r)},\mathbf{T}^{(r)};t)\)，并按照\autoref{sec:local-tracking} 中的建模建立局部目标函数：路径点同步项 \(J_{\mathrm{sync}}^{(r)}\) 在各时间戳 \(\tau_{k,s}\) 处强制对齐全局参考点，可微避障项 \(J_{o}^{(r)}\) 由局部 \(\{\mathbf{p},\mathbf{v}\}\) 对与零距离平面构造，高阶导平滑度 \(J_{s}^{(r)}\)、总时长 \(J_{t}^{(r)}\) 与动态可行性项 \(J_{d}^{(r)}\) 则约束能量与物理边界，并使用 L\text{-}BFGS 结合显式梯度进行求解，得到连续时间的最优轨迹 \(\mathbf{p}^{(r)}(t),\dot{\mathbf{p}}^{(r)}(t),\ddot{\mathbf{p}}^{(r)}(t)\)。

求解得到的局部轨迹会在执行层被实时采样并转化为控制指令：对 UAV，将轨迹通过微分平坦映射生成推力与角速度参考；对 UGV 直接使用平面速度与加速度参考驱动底层控制器。外环以较低频率运行，主要负责全局可视性与队形结构的一致性；内环以较高频率运行，重点保证连续轨迹的动态可行性与局部避障。通过时间戳同步与滚动时域更新，外环每次优化覆盖至少一个内环预测区间，实现全局与局部的稳定衔接，进而在复杂未知环境下兼顾全局协同、可视保持与实时安全执行。

\section{实验验证} \label{sec:experiments}

\subsection{实验设置}

本章通过仿真与实物实验对所提出的两阶段协同规划方法进行验证。仿真实验环境与上一章相同，构造稀疏与稠密两类障碍场景，用于考察在同一飞行任务下，不同约束配置对空地系统平均可视率、地面队形几何偏差、任务成功率以及整体轨迹长度等指标的影响。实物实验在约 $12\times16\times3\,\mathrm{m}^3$ 的室内场地中进行，场景包含一组固定障碍布置，障碍物之间的平均间距约为 $1.5\,\mathrm{m}$，任务目标点设置在 $(15,0)\,\mathrm{m}$ ，到达任务点后返回原点，视作任务完成，用于验证算法在真实传感噪声和通信延迟下的可行性与鲁棒性。

\begin{figure}[htbp]
    \centering
    \includegraphics[width=0.65\linewidth]{chapter_4/realworld_ugvuav.png}
    \caption{实物实验机器人平台}
    \label{fig:real_expr_uav_ugv}
\end{figure}

实物实验平台由一架自组装小型四旋翼无人机（UAV）和四辆麦克纳姆轮无人车（UGV）组成，如\autoref{fig:real_expr_uav_ugv} 所示。UAV 搭载 UWB 标签与机载计算机（Nvidia Jetson Orin），用于接收 UGV 上 UWB 锚点的测距信息并完成相对定位，同时进行本地规划；UGV 搭载 Livox MID-360 激光雷达、UWB 基站以及机载计算机（Nvidia Jetson Orin），负责环境建图与本地规划。空地系统整体定位与地图构建采用第二章中的 UWB–LiDAR 融合方法，实现空地系统间的统一坐标系与相对位姿估计。全局协同规划在一台本地笔记本电脑上运行，规划周期为 $2\,\mathrm{s}$，负责更新全局参考路径与序列可视空间；局部轨迹优化在各机器人板载计算机上独立运行，周期为 $0.5\,\mathrm{s}$，负责将全局离散参考转化为可执行的连续轨迹。底层运动控制频率为$100\,\mathrm{Hz}$，高于局部规划频率，由各平台的底层控制器根据轨迹插值得到期望状态。

协同优化采用 L\text{-}BFGS 作为求解器，其主要参数设置如表~\ref{tab:lbfgs_param} 所示；全局协同规划与局部轨迹优化中各子目标的权重分别见表~\ref{tab:global_weights} 和表~\ref{tab:local_weights}，仿真与真实实验的参数设置保持一致。

\begin{table}[htbp]
    \centering
    \caption{L\text{-}BFGS 优化器主要参数设置}
    \label{tab:lbfgs_param}
    \begin{tabular}{ccc}
        \toprule
        参数名 & 数值 & 说明 \\
        \midrule
        \texttt{mem\_size}       & $8$              & 反向更新中保存的历史 $(s,y)$ 对数量 \\
        \texttt{max\_iterations} & $600$            & 每次调用的最大迭代次数 \\
        \texttt{min\_step}       & $1\times10^{-45}$& 线搜索允许的最小步长 \\
        \texttt{max\_linesearch} & $600$            & 每次迭代的最大线搜索次数 \\
        \texttt{past}            & $18$             & 用于收敛判据的历史窗口长度 \\
        \texttt{delta}           & $1\times10^{-1}$ & 相邻窗口目标值变化阈值 \\
        \bottomrule
    \end{tabular}
\end{table}

\begin{table}[htbp]
    \centering
    \caption{全局协同规划中各子目标权重设置}
    \label{tab:global_weights}
    \begin{tabular}{ccc}
        \toprule
        代价项                         & 权重符号                          & 数值 \\
        \midrule
        动力学一致性 $J_{\mathrm{dyn}}$         & $\lambda_{\mathrm{dyn}}$          & $1000.0$ \\
        能量与平滑项 $J_{\mathrm{eng}}$         & $\lambda_{e},\,\lambda_{s}$       & $2.0$ \\
        速度/加速度上下限 $J_{\mathrm{lim}}$    & $\lambda_{\mathrm{bound}}$        & $1000.0$ \\
        可视性（SSCP）$J_{\mathrm{vis}}$        & $\lambda_{\mathrm{vis}}$          & $10.0$ \\
        UAV 引导路点跟踪 $J_{\mathrm{wp}}$      & $\lambda_{\mathrm{wp}}$           & $6.0$ \\
        UGV 边长均衡 $J_{\mathrm{edge}}$        & $\lambda_{\mathrm{edge}}$         & $4.0$ \\
        UGV 平均边长尺度 $J_{\mathrm{mean}}$    & $\lambda_{\mathrm{mean}}$         & $4.0$ \\
        UGV 径向均衡 $J_{\mathrm{rad}}$         & $\lambda_{\mathrm{rad}}$          & $4.0$ \\
        空地测距边界 $J_{\mathrm{bd}}$          & $\lambda_{\mathrm{bd}}$           & $5.0$ \\
        \bottomrule
    \end{tabular}
\end{table}

\begin{table}[htbp]
    \centering
    \caption{局部轨迹优化中各子目标权重设置}
    \label{tab:local_weights}
    \begin{tabular}{ccc}
        \toprule
        代价项                           & 权重符号                 & 数值 \\
        \midrule
        局部避障项 $J_{o}$                & $\lambda_{o}$            & $10000$ \\
        路径点同步项 $J_{\mathrm{sync}}$  & $\lambda_{\mathrm{sync}}$& $100$ \\
        高阶平滑度项 $J_{s}$             & $\lambda_{s}$            & $0.5$ \\
        总时间项 $J_{t}$                 & $\lambda_{t}$            & $1000$ \\
        动态可行性惩罚项 $J_{d}$         & $\lambda_{d}$            & $10000$ \\
        \bottomrule
    \end{tabular}
\end{table}

\subsection{仿真实验}

本节在与上一章相同的随机矩形障碍栅格环境中，构造稀疏与稠密两种障碍密度，并在每种环境下分别开启与关闭序列可视空间代价 \(J_{\mathrm{vis}}\)（即 SSCP 约束），共得到四组对比场景。四组实验均采用一架 UAV 与四辆 UGV，从相同起点出发，执行相同的导航任务，使用本章的两阶段方法在仿真环境中运行的过程如\autoref{fig:sim_expr_gazebo} 所示。其中图a展示了Gazebo的仿真环境，图b为对应的 RViz 显示界面，包含 3D 占据栅格地图及各机器人的历史轨迹。图c为同一场景的 2D 示意图，将障碍物和机器人位置投影到了地平面上。图d中的红色轨迹展示了各机器人在当前规划时域内需跟随的优化路点序列。直线连接了无人机与各无人车，绿色连线表示相应的空地（UAV-UGV）组合处于视距内，否则显示为红色。

\begin{figure}[htbp]
    \centering
    \includegraphics[width=0.8\linewidth]{chapter_4/sim_running.png}
    \caption{仿真环境运行样例}
    \label{fig:sim_expr_gazebo}
\end{figure}

机器人在全局层面由第~\ref{sec:sscp-global} 节提出的协同优化生成离散参考，再由第~\ref{sec:local-tracking} 节的局部轨迹进行跟踪与避障。稀疏环境下的轨迹对比如 \autoref{fig:traj_sparse} 所示，稠密环境下的轨迹对比如 \autoref{fig:traj_dense} 所示。图中颜色表示速度大小，红色小方块为膨胀后的障碍物栅格。为定量评估两阶段方法的效果，使用与之前相同的四类评价指标。四组实验的统计结果汇总于 \autoref{tab:sim_metrics}。从 \autoref{fig:traj_sparse}、\autoref{fig:traj_dense} 与 \autoref{tab:sim_metrics} 可以得到以下结果。

% --- 第一组图：稀疏环境 ---
\begin{figure}[htbp]
    \centering
    % 第一张子图
    \begin{subfigure}{\textwidth}
        \centering
        \includegraphics[width=0.85\textwidth]{chapter_4/sparse_env_traj_full_constrain.png}
        \caption{稀疏环境，启用可视性}
        \label{fig:sparse_full}
    \end{subfigure}
    \vspace{1em} % 增加垂直间距
    
    % 第二张子图
    \begin{subfigure}{\textwidth}
        \centering
        \includegraphics[width=0.85\textwidth]{chapter_4/sparse_env_traj_noVisConstrain.png}
        \caption{稀疏环境，关闭可视性}
        \label{fig:sparse_novis}
    \end{subfigure}
    
    \caption{稀疏环境下 UAV 与 UGV 的轨迹对比}
    \label{fig:traj_sparse}
\end{figure}

% --- 第二组图：稠密环境 ---
\begin{figure}[htbp]
    \centering
    % 第三张子图
    \begin{subfigure}{\textwidth}
        \centering
        \includegraphics[width=0.85\textwidth]{chapter_4/dense_env_traj_full_constrain.png}
        \caption{稠密环境，启用可视性}
        \label{fig:dense_full}
    \end{subfigure}
    \vspace{1em}
    
    % 第四张子图
    \begin{subfigure}{\textwidth}
        \centering
        \includegraphics[width=0.85\textwidth]{chapter_4/dense_env_traj_noVisConstrain.png}
        \caption{稠密环境，关闭可视性}
        \label{fig:dense_novis}
    \end{subfigure}
    
    \caption{稠密环境下 UAV 与 UGV 的轨迹对比}
    \label{fig:traj_dense}
\end{figure}

首先，在两种环境下启用 SSCP 可视性约束后，平均可视率均显著提升：稀疏环境中从 0.67 提升到 0.80，稠密环境中从 0.65 提升到 0.75；同时任务成功率分别由 40\% 与 20\% 提高到 100\% 与 80\%。对应轨迹中可以看到，启用可视性时 UAV 倾向于在 UGV 队形上方保持“伴飞”，而不再为了缩短自身路径而远离地面平台，从而减小因遮挡导致的定位退化与任务失败。

其次，队形保持指标表明，可视性约束并未打乱 UGV 的整体几何结构，反而在稠密场景下显著抑制了退化队形。以稠密环境（\autoref{fig:traj_dense}）为例，启用 SSCP 后半径偏差由 0.546\,m 降至 0.217\,m，边长偏差由 0.950\,m 降至 0.593\,m，说明 UGV 锚点在平面上的分布更均匀，也有利于 UWB 相对定位的条件数。稀疏场景（\autoref{fig:traj_sparse}）下两种配置的边长偏差接近，但启用可视性时角度偏差显著减小，队形旋转更加平衡。

\begin{table}[htbp]
    \centering
    \caption{不同环境与是否启用可视性约束下的仿真实验统计结果}
    \label{tab:sim_metrics}
    \resizebox{\textwidth}{!}{%
        \begin{tabular}{cccccccc}
            \toprule
            \multirow{2}{*}{场景} & \multirow{2}{*}{可视性} & 平均 & 半径偏差 & 角度偏差 & 边长偏差 & \multirow{2}{*}{成功率} & 轨迹平均长度 \\
             & & 可视率 & ($\mathrm{m}$) & ($^\circ$) & ($\mathrm{m}$) & & (UAV/UGV) / $\mathrm{m}$ \\
            \midrule
            稀疏 & 启用 & 0.80 & 0.282 & 11.756 & 0.493 & 100\% & 32.1 / 26.2 \\
            稀疏 & 关闭 & 0.67 & 0.341 & 17.875 & 0.509 & 40\%  & 38.3 / 23.9 \\
            稠密 & 启用 & 0.75 & 0.217 & 26.134 & 0.593 & 80\%  & 40.4 / 32.3 \\
            稠密 & 关闭 & 0.65 & 0.546 & 38.240 & 0.950 & 20\%  & 53.1 / 32.4 \\
            \bottomrule
        \end{tabular}%
    }
\end{table}

再次，从轨迹长度来看，引入可视性约束后 UAV 在稀疏场景中的路径略有增加（32.1\,m vs. 38.3\,m），在稠密场景中则明显缩短（40.4\,m vs. 53.1\,m）。这是因为在不使用 SSCP 时，全局协同优化缺乏对障碍几何与可视性的显式刻画，容易将同步路点拖拽到障碍物间隙甚至占据体素内部；局部规划器在试图同时满足这些“不可达”同步点与避障约束时，只能产生大幅绕行或陷入求解困难，导致整体路径拉长甚至任务失败。启用 SSCP 后，全局层面对 UGV 队形施加了时序化的可视空间软约束，将参考状态限制在可见且远离障碍的星凸区域内，局部规划只需在该可行通道中进行微调，从而兼顾了路径长度、可视率与收敛稳定性。

综上，在稀疏与稠密两种典型场景中，所提出的两阶段协同规划框架配合 SSCP 可视性建模，可以在保持队形几何合理性的前提下显著提高空地系统的视距保持率与任务成功率，尤其在障碍物稠密、几何关系复杂的环境中优势更加明显。

\subsection{实物实验}\label{sec:real-exp}
在实验设置中所述的 $12\times16\times3\,\mathrm{m}^3$ 室内场地内，我们进行了多轮空地协同实物实验，以验证两阶段协同规划在真实传感噪声、通信时延和动力学不确定性下的有效性。场地中布置多组柱状障碍物，平均间距约为 $1.5\,\mathrm{m}$。

\begin{figure}[htbp]
    \centering
    \includegraphics[width=\linewidth]{chapter_4/real_expr.png}
    \caption{实物实验场景及运行过程}
    \label{fig:real_expr}
\end{figure}

\autoref{fig:real_expr} 给出了其中一组实验的空间示意和关键时刻的局部占据栅格视图。图中上半部分展示了 UAV 与四辆 UGV 在障碍物间协同机动的侧视图，下半部分的 (1)–(4) 依次给出 UAV 轨迹在局部栅格地图中的位姿与运行轨迹：从初始的 LoS 阶段 (1)，到部分 UGV 被柱体遮挡形成 NLoS (2)，再到 UAV 主动调整轨迹恢复 LoS (3)，最终在保持队形几何的前提下完成穿越任务 (4)。这一过程直观展示了序列可视空间约束在实物系统中的作用。

为量化可视性表现，\autoref{fig:vis_curve} 展示了整段运行过程的可视率随归一化时间的变化曲线，其中实线为平均可视率，阴影带为 $\pm 1$ 标准误差。整体平均可视率为 $0.83$，在障碍物密集区域虽出现多次下降，但大部分时间保持在 $0.75$ 以上。这表明在真实建图误差和测量噪声存在的情况下，基于 SSCP 的全局规划仍能有效驱动空地系统实现 LoS 视距。

队形保持采用与仿真实验相同的指标进行评估。统计所有时间步和多轮实验数据，径向偏差、角度偏差和边长偏差的整体平均值分别为
\(\bar e_r=0.200\,\mathrm{m}\)、\(\bar e_\varphi=8.262^\circ\) 和 \(\bar e_{\mathrm{edge}}=0.504\,\mathrm{m}\)，说明四辆 UGV 在复杂障碍环境中仍能维持较为规整的多边形队形。整组实验的任务成功率为 $80\%$，典型成功任务中 UAV 轨迹长度约为 $53.8\,\mathrm{m}$，UGV 轨迹平均长度约为 $27.7\,\mathrm{m}$，与仿真结果在量级和趋势上保持一致。

此外，为验证空地协同定位性能，在每轮实验结束后，我们测量 UAV 实际着陆点与目标着陆点在地面坐标系下的误差。统计结果表明，UAV 终端定位误差在 $x$ 轴和 $y$ 轴上的平均值分别为
\(\bar e_x= \textbf{7.8~cm} \) 和 \(\bar e_y=\textbf{8.84~cm}\)。这表明基于动态 UWB 锚点与队形约束的定位方案在实物环境中能够稳定提供优于 $10\,\mathrm{cm}$ 的定位精度，为上层协同规划与控制提供了可靠的状态基础。

\begin{figure}[htbp]
    \centering
    \includegraphics[width=\linewidth]{chapter_4/vis_compare.png}
    \caption{实物实验中可视率随归一化时间的变化曲线}
    \label{fig:vis_curve}
\end{figure}

\section{本章小结} \label{sec:summary_opt2}
本章围绕“基于序列可视空间的两阶段协同规划”展开研究，面向空地异构多机器人系统在复杂障碍环境中的长期协同导航问题，给出了从可视空间建模到全局–局部一体化规划的完整方法。

首先，针对传统基于射线检测的 LoS 约束维度高、不可导的问题，本章使用基于球翻转和凸包逼近的星型凸可视空间表示方法。以混合 A* 生成的动力学可行引导路径为骨架，在每个等时采样点周围构建局部星型凸多面体，并通过对数–指数光滑近似得到可微的可视性代价及其梯度，进而在时间轴上串联形成序列可视空间（SSCP），实现了对时序可视约束的紧凑参数化表达。

其次，在全局层面，基于统一的二阶积分离散动力学，对 UAV 与多辆 UGV 的位置、速度和加速度进行联合建模，引入引导路点跟踪、序列可视空间约束和 空地队形约束三类协同目标，并叠加动力学一致性、速度/加速度上下限及能量与平滑度惩罚，构建了多目标非线性优化问题。为适应长预测域与多机器人耦合，本章采用 L\text{-}BFGS 进行无约束优化，通过热启动与滚动时域执行在保证数值稳定性的同时实现了实时可解的全局协同轨迹生成。

再次，在局部层面，本章采用 MINCO 作为连续轨迹表示，引入基于局部切片与锚点对的可微避障代价，以及对带时间戳全局路点的同步跟踪约束，并结合最小高阶导平滑度、总时长和动态可行性惩罚，构造了一阶可导的局部轨迹优化目标。局部优化同样使用 L\text{-}BFGS 在各机器人板载执行，实现了在局部地图上的实时避障与轨迹细化。

在此基础上，本章最后给出了两阶段协同规划的体系架构：外环以较低频率运行全局协同优化，输出带时间戳的多机器人参考序列；内环以较高频率在各机器人上独立执行局部 MINCO 优化，通过时间戳对齐实现“全局可视一致性与队形协同”与“局部连续可执行与安全避障”的解耦与闭环。本章方法在仿真与实物实验中表现出较高的任务成功率、可视率和队形保持性能，充分验证了算法的可行性与有效性。

\chapter{总结与展望}

\section{本文工作总结}

随着无人机与无人车技术的发展，空地协同系统在复杂环境下的应用潜力日益凸显，但在缺乏外部基础设施的未知与遮挡环境中，如何实现高精度的相对定位、主动维护通信与视线连接并生成安全高效的协同轨迹，仍是制约其大规模落地的关键瓶颈。针对上述问题，本文围绕“动态 UWB 支撑的相对定位”与“可见性驱动的协同规划”两条主线，构建了一套完整的空地协同系统框架。首先，针对传统 UWB 定位依赖固定基站且部署繁琐的问题，本文提出了基于移动 UWB 基站的空地相对定位与环境共享方法。该方法在测距层面设计了高效的多基站 TWR 时序协议，通过“单次轮询—顺序应答”机制显著降低了通信开销，在多基站配置下实现了百赫兹级的高频测距；在状态估计层面，提出了基于多边定位初值与误差状态扩展卡尔曼滤波（ESKF）的融合算法，利用 UGV 的 LiDAR-SLAM 全局坐标系作为统一参考，将机载 IMU 的高频特性与 UWB 的无偏位置观测紧耦合，实现了无人机在无 GPS 环境下的高精度相对位姿估计；同时，建立了基于三维占据栅格的地图共享机制，通过局部子图裁剪与世界坐标系对齐融合，实现了空地平台间低带宽下的环境信息同步，为后续协同规划提供了统一的状态与环境输入。

在解决定位与感知问题的基础上，本文针对中等复杂度环境下的空地协同导航问题，提出了基于点云聚类与椭圆柱近似的可见性协同规划方法。该方法利用 DBSCAN 算法对环境点云进行聚类，并提取最小包络椭圆柱作为障碍物的凸几何近似，实现了非结构化环境的轻量化参数表达。在此几何表示之上，结合离散控制障碍函数（DCBF）构建了可微的几何避障约束，并设计了空地视线（LoS）连通性代价与地面编队几何约束，将避障安全、可见性保持与队形稳定性统一纳入非线性模型预测控制（NMPC）框架中。仿真结果表明，该单层优化方法在稀疏至中等密度的障碍物环境中，能够实时生成平滑、安全且兼顾空地视距质量的协同轨迹。

为了进一步解决单层优化在稠密障碍物环境中求解困难、收敛性差的问题，提升系统在复杂环境下的泛化能力，本文提出了基于序列可视空间的两阶段协同规划架构。该架构将规划问题解耦为“全局一致性规划”与“局部安全执行”两个层面。在全局规划层，引入序列可视空间（SSCP）概念，通过沿混合 A* 引导路径构建时序化的星凸可视区域，将复杂的避障与视距约束转化为光滑的可视性代价函数，利用 L-BFGS 求解器快速生成满足空地可见性与队形几何要求的全局参考序列。在局部规划层，各机器人独立运行基于 MINCO 轨迹参数化的局部优化器，在严格跟踪全局时间戳的同时，利用基于局部点对的避障力场进行微调，保证了轨迹的动力学可行性与安全性。仿真与实物实验证明，该两阶段架构显著提升了系统在复杂稠密环境中的任务成功率与可视率，验证了所提方案在真实场景中的有效性。

\section{未来工作展望}

尽管本文在空地协同定位与规划方面取得了一定进展，但受限于实验条件与理论假设，相关研究仍存在进一步深化与拓展的空间。首先，目前的规划框架主要采用集中式或半分布式架构，假设空地之间通信链路畅通且带宽满足需求，然而在实际大范围作业中，通信时延、丢包以及带宽限制不可忽略。未来的工作可以探索完全分布式的协同规划算法，结合一致性理论或分布式 ADMM 方法，仅依靠邻居间的局部信息交互实现全局协同，并将通信质量（如信号强度、误码率）显式纳入规划代价，实现“通信—感知—运动”的深度耦合。

其次，本文主要处理静态障碍物环境，对于移动障碍物（如行人、其他车辆）的处理主要依赖高频重规划进行被动规避，这在动态环境下的安全性略显不足。未来可引入对动态障碍物的轨迹预测模块，构建基于时空联合的避障约束，或者利用博弈论思想，研究多机器人与动态环境实体之间的交互式规划策略，提升系统在人机共存等高动态场景下的安全性和适应性。

此外，本文的实验验证主要基于“一机四车”的小规模编队，随着任务复杂度的提升，系统可能需要扩展至更大规模的异构集群。在大规模集群下，如何设计可扩展的队形控制策略、降低计算复杂度以及解决高维状态空间下的优化求解问题是亟待解决的难点，未来可考虑引入分层控制架构或基于学习的群体策略，以提高系统对规模扩展的适应能力。

最后，在定位方面，本文主要依赖 UWB 与 IMU 的融合并以激光 SLAM 地图为基准，但在强电磁干扰或长走廊等极端环境下，定位系统的鲁棒性仍面临挑战。未来可以研究融合视觉特征、语义信息以及多机间的相互观测约束，构建视觉-激光-惯性-UWB 全源紧耦合的协同定位系统，利用异构传感器的互补特性进一步提升系统在全天候、全场景下的状态估计精度与可靠性。
